% Leçon 29 — Éléments socioculturels : différence entre la monnaie chinoise et la monnaie camerounaise — Chinois Tle A — Cours signé OnBuch+
% Programme officiel MINESEC. Compiler avec : tectonic ZH29-socio-difference-entre-la-monnaie-chinoise-et-la-m.tex
\newcommand{\DOCMATIERE}{Chinois}
\newcommand{\DOCNIVEAU}{Tle A}
\newcommand{\DOCMODULE}{Module 4 : Activités économiques et monde du travail}
\newcommand{\DOCLECON}{ZH29}
\newcommand{\DOCTITRE}{Leçon 29 — Éléments socioculturels : différence entre la monnaie chinoise et la monnaie camerounaise}
% OnBuch+ — préambule « cours signés » ENRICHI (compilé avec tectonic / XeLaTeX)
% Système visuel « Pop » d'OnBuch+ : fond crème (jamais blanc), bordures encre
% épaisses, ombres « dures » décalées, titres Archivo Black en capitales.
% Version MATHÉMATIQUES : graphes (pgfplots/tikz), boîtes pédagogiques
% (définition, propriété, méthode, attention, le savais-tu, exercice résolu,
% exercice, TP, bilan), page de garde et signature OnBuch+ ancrée en bas.
%
% Macros attendues dans main.tex AVANT \input{preamble} :
%   \DOCMATIERE  \DOCNIVEAU  \DOCMODULE  \DOCLECON (numéro)  \DOCTITRE
\documentclass[11pt]{article}

\usepackage{fontspec}
\usepackage[french]{babel} % français = langue principale ; le chinois via \zh{..} / textebox (xeCJK)
\usepackage{xeCJK}
\usepackage{pifont} % \ding{51} \ding{55} utilisés par les modèles
\usepackage[a4paper,margin=2cm,top=3.4cm,bottom=2.8cm,headheight=1.4cm,footskip=1.3cm]{geometry}
\usepackage{amsmath,amssymb,amsfonts}
\usepackage{mathtools} % \xmapsto, \coloneqq... souvent utilisés par les modèles
\usepackage{cancel} % simplifications barrées (bilans de forces)
% \abs{z} (module, valeur absolue) et \norm{u} (norme), utilisés par les modèles
\providecommand{\abs}{}\let\abs\relax\DeclarePairedDelimiter{\abs}{\lvert}{\rvert}
\providecommand{\norm}{}\let\norm\relax\DeclarePairedDelimiter{\norm}{\lVert}{\rVert}
\usepackage[table]{xcolor} % [table] : fournit aussi \rowcolors (alternance de lignes)
\usepackage{pagecolor}
\usepackage{tikz}
\usetikzlibrary{arrows.meta,positioning,calc,shapes.geometric,shapes.misc,shapes.arrows,decorations.pathmorphing,decorations.pathreplacing,decorations.markings,patterns,shadows,mindmap,trees,fit,backgrounds,matrix,intersections,angles,quotes}
\usepackage{pgfplots}
\usepackage[version=4]{mhchem} % équations nucléaires \ce{^{235}_{92}U}
\usepackage[european,siunitx]{circuitikz} % schémas électriques
\pgfplotsset{compat=1.18}
\usepgfplotslibrary{fillbetween,groupplots} % aires entre courbes (calcul intégral)
\usepackage{enumitem}
\usepackage{fancyhdr}
% [most] charge toutes les bibliothèques tcolorbox (listings, minted, raster,
% documentation...) et gonfle inutilement la mémoire TeX sur un document long
% (~300 boîtes) : on ne charge que ce qui est réellement utilisé.
\usepackage[skins,breakable]{tcolorbox}
\usepackage{graphicx}
\usepackage{colortbl}
\usepackage{booktabs}
\usepackage{tabularx}
\usepackage{diagbox} % case d'en-tête diagonale des tableaux à double entrée
% Colonnes extensibles centrée / gauche / droite (C, L, R), souvent utilisées par les modèles
\newcolumntype{C}{>{\centering\arraybackslash}X}
\newcolumntype{L}{>{\raggedright\arraybackslash}X}
\newcolumntype{R}{>{\raggedleft\arraybackslash}X}
\usepackage{array}
\usepackage{multicol}
\usepackage{multirow}
\usepackage{siunitx}
% parse-numbers=false : accepte aussi \SI{2,5 \times 10^{-3}}{...}
\sisetup{locale=FR,output-decimal-marker={,},group-minimum-digits=4,parse-numbers=false}
\DeclareSIUnit\ohm{\ensuremath{\symup{\Omega}}} % U+2126 absent de la police
\DeclareSIUnit\division{div} % réglages d'oscilloscope (ms/div, V/div)
\DeclareSIUnit\tr{tr} % tours (vitesse de rotation, ex. \SI{1500}{\tr\per\minute})
\DeclareSIUnit\molar{M} % concentration molaire (ex. \SI{3}{\molar}, \SI{1}{\micro\molar})
\DeclareSIUnit\an{an} % année (le modèle confond parfois avec \year, un compteur TeX) — ex. \SI{5}{\kilo\meter\per\an}
\DeclareSIUnit\FCFA{FCFA} % franc CFA (ex. \SI{500}{\FCFA\per\kilo\gram})
\DeclareSIUnit\degreeBrix{\textdegree Brix} % teneur en sucre d'un sirop/jus (réfractométrie)
\DeclareSIUnit\month{mois} % le modèle utilise parfois \month, un compteur TeX, comme unité de durée
\usepackage{qrcode}
% \degree : commande fréquemment utilisée par les modèles pour un angle ou une
% température en dehors de \SI{}{} (non fournie par les paquets chargés).
\providecommand{\degree}{^{\circ}}
% \qed (fin de démonstration), utilisé par les modèles sans amsthm
\providecommand{\qed}{\hfill$\square$}
% \barre : double barre des valeurs interdites dans un tableau de variations (tkz-tab)
\providecommand{\barre}{\ensuremath{\|}}
% environnement proof (amsthm non chargé) utilisé par les modèles
\ifdefined\proof\else\newenvironment{proof}[1][Démonstration]{\par\noindent\textit{#1.}\ }{\qed\par}\fi
\usepackage{lastpage}
\usepackage{hyperref}
\hypersetup{hidelinks}

% ============================================================
% Palette « Pop » OnBuch+ — mêmes hex que la classe P (plus_ui.dart)
% ============================================================
\definecolor{poporange}{HTML}{F59321}
\definecolor{popdark}{HTML}{E07A0C}
\definecolor{popcream}{HTML}{FFF6E8}
\definecolor{popink}{HTML}{1C1714}
\definecolor{poppurple}{HTML}{7A5AE0}
\definecolor{popgreen}{HTML}{2E9E6A}
\definecolor{poppink}{HTML}{E0457B}
\definecolor{popblue}{HTML}{2F7DE1}
\definecolor{popgold}{HTML}{C98A12}
\definecolor{popmuted}{HTML}{8A7F76}
\definecolor{popline}{HTML}{EADFD2}
\definecolor{popgray}{HTML}{8A7F76}
\colorlet{popgrey}{popgray}
% Alias tolérants (le modèle nomme parfois les couleurs différemment)
\colorlet{popwhite}{white}
\colorlet{popblack}{popink}
\colorlet{popred}{poppink}
\colorlet{popyellow}{popgold}
% Teintes claires (fonds de tableaux/encadrés) — le modèle nomme parfois
% les couleurs avec un suffixe L (ex. popblueL) sans les avoir définies.
\colorlet{poporangeL}{poporange!20}
\colorlet{popdarkL}{popdark!20}
\colorlet{poppurpleL}{poppurple!20}
\colorlet{popgreenL}{popgreen!20}
\colorlet{poppinkL}{poppink!20}
\colorlet{popblueL}{popblue!20}
\colorlet{popgoldL}{popgold!20}
\colorlet{popredL}{popred!20}
\colorlet{popgrayL}{popgray!20}
% Teintes claires pour fonds de tableaux / zones
\colorlet{popblueL}{popblue!10!white}
\colorlet{poporangeL}{poporange!14!white}
\colorlet{popgreenL}{popgreen!12!white}
\colorlet{poppurpleL}{poppurple!12!white}
\colorlet{poppinkL}{poppink!12!white}
\colorlet{popgoldL}{popgold!14!white}

\pagecolor{popcream}

% ============================================================
% Typographie
% ============================================================
\defaultfontfeatures{Ligatures=TeX}
\setmainfont{PlusJakartaSans-Medium.ttf}[Path=../../../fonts/,BoldFont=PlusJakartaSans-Bold.ttf]
% Symboles, API et emojis absents de la police principale : repli sur DejaVu Sans (caractères actifs)
\newfontface\symfont{DejaVuSans.ttf}[Path=../../../fonts/]
\makeatletter
\newcommand\symactive[1]{\catcode#1=13 \begingroup\lccode`\~=#1 \lowercase{\endgroup\def~}{{\symfont\char#1}}}
\newcommand\symblank[1]{\catcode#1=13 \begingroup\lccode`\~=#1 \lowercase{\endgroup\def~}{}}
\newcommand\symhyph[1]{\catcode#1=13 \begingroup\lccode`\~=#1 \lowercase{\endgroup\def~}{\mbox{-}}}
\newcount\symc
\newcommand\symrange[2]{\symc=#1 \@whilenum\symc<#2 \do{\expandafter\symactive\expandafter{\the\symc}\advance\symc by 1 }}
\newcommand\symblankrange[2]{\symc=#1 \@whilenum\symc<#2 \do{\expandafter\symblank\expandafter{\the\symc}\advance\symc by 1 }}
\makeatother
\symrange{"0250}{"02FF}\symactive{"2713}\symactive{"2714}\symactive{"2717}\symactive{"2718}\symactive{"2610}\symactive{"2611}\symactive{"2612}
\symrange{"2600}{"27BF}
\catcode"2705=13 \begingroup\lccode`\~="2705 \lowercase{\endgroup\def~}{{\symfont\char"2713}}\symblank{"26BD}\symblank{"26A1}
\symrange{"2460}{"257F}\symactive{"223C}\symactive{"2248}\symactive{"2260}
\symactive{"01F8}\symactive{"01F9}\symactive{"1E3E}\symactive{"1E3F}
\symhyph{"2011}\symhyph{"2E1A}\symblankrange{"1F300}{"1FAFF}\symblank{"FE0F}
% Caractères chinois : WenQuanYi Zen Hei (sous-ensemble GB2312 + pinyin), gras/italique simulés
\setCJKmainfont{WenQuanYiZenHei-subset.ttf}[Path=../../../fonts/,AutoFakeBold=3,AutoFakeSlant=0.2]
\xeCJKsetup{CJKspace=false}
% Pinyin : voyelles à caron (ǎ ǐ ǒ ǔ ǖ ǘ ǚ ǜ) absentes de la police principale -> caractères actifs
\newfontface\pyfont{WenQuanYiZenHei-subset.ttf}[Path=../../../fonts/]
\newcommand\pyactive[1]{\catcode#1=13 \begingroup\lccode`\~=#1 \lowercase{\endgroup\def~}{{\pyfont\char#1}}}
\pyactive{"01CD}\pyactive{"01CE}\pyactive{"01CF}\pyactive{"01D0}\pyactive{"01D1}\pyactive{"01D2}\pyactive{"01D3}\pyactive{"01D4}\pyactive{"01D5}\pyactive{"01D6}\pyactive{"01D7}\pyactive{"01D8}\pyactive{"01D9}\pyactive{"01DA}\pyactive{"01DB}\pyactive{"01DC}
% Maths en sans-serif (Fira Math) avec chiffres et lettres droites de Plus
% Jakarta, pour que les formules se fondent dans le texte.
\usepackage[math-style=ISO,bold-style=ISO]{unicode-math}
\setmathfont{FiraMath-Regular.otf}[Path=../../../fonts/]
\setmathfont{PlusJakartaSans-Medium.ttf}[Path=../../../fonts/,range={up/num,up/{Latin,latin},\mathup}]
\usepackage{icomma} % virgule décimale sans espace en mode math
% Caractères mathématiques Unicode absents des polices de texte (Plus Jakarta,
% Archivo Black) : rendus par leur équivalent mathématique, en texte comme en
% formule — sinon ils s'affichent en carré vide (ex. « Arithmétique dans ℤ »).
\usepackage{newunicodechar}
\newunicodechar{ℝ}{\ensuremath{\mathbb{R}}}
\newunicodechar{ℤ}{\ensuremath{\mathbb{Z}}}
\newunicodechar{ℂ}{\ensuremath{\mathbb{C}}}
\newunicodechar{ℕ}{\ensuremath{\mathbb{N}}}
\newunicodechar{ℚ}{\ensuremath{\mathbb{Q}}}
\newunicodechar{𝔻}{\ensuremath{\mathbb{D}}}
\newunicodechar{α}{\ensuremath{\alpha}}
\newunicodechar{β}{\ensuremath{\beta}}
\newunicodechar{γ}{\ensuremath{\gamma}}
\newunicodechar{δ}{\ensuremath{\delta}}
\newunicodechar{ε}{\ensuremath{\varepsilon}}
\newunicodechar{θ}{\ensuremath{\theta}}
\newunicodechar{λ}{\ensuremath{\lambda}}
\newunicodechar{μ}{\ensuremath{\mu}}
\newunicodechar{σ}{\ensuremath{\sigma}}
\newunicodechar{τ}{\ensuremath{\tau}}
\newunicodechar{φ}{\ensuremath{\varphi}}
\newunicodechar{ψ}{\ensuremath{\psi}}
\newunicodechar{ω}{\ensuremath{\omega}}
\newunicodechar{Σ}{\ensuremath{\Sigma}}
% majuscules produites par \MakeUppercase dans les titres (α-aminés -> Α)
\newunicodechar{Α}{\ensuremath{\alpha}}
\newunicodechar{Β}{\ensuremath{\beta}}
\newunicodechar{Γ}{\ensuremath{\Gamma}}
\newunicodechar{Δ}{\ensuremath{\Delta}}
\newunicodechar{Ε}{\ensuremath{\varepsilon}}
\newunicodechar{Θ}{\ensuremath{\Theta}}
\newunicodechar{Λ}{\ensuremath{\Lambda}}
\newunicodechar{Μ}{\ensuremath{\mu}}
\newunicodechar{Τ}{\ensuremath{\tau}}
\newunicodechar{Φ}{\ensuremath{\Phi}}
\newunicodechar{Ψ}{\ensuremath{\Psi}}
\newunicodechar{Ω}{\ensuremath{\Omega}}
\newunicodechar{↦}{\ensuremath{\mapsto}}
\newunicodechar{∈}{\ensuremath{\in}}
\newunicodechar{∉}{\ensuremath{\notin}}
\newunicodechar{∧}{\ensuremath{\wedge}}
\newunicodechar{⇔}{\ensuremath{\Leftrightarrow}}
\newunicodechar{⟺}{\ensuremath{\Longleftrightarrow}}
\newunicodechar{⇒}{\ensuremath{\Rightarrow}}
\newunicodechar{⟹}{\ensuremath{\Longrightarrow}}
\newunicodechar{∘}{\ensuremath{\circ}}
\newunicodechar{∪}{\ensuremath{\cup}}
\newunicodechar{∩}{\ensuremath{\cap}}
\newunicodechar{≡}{\ensuremath{\equiv}}
\newunicodechar{⊂}{\ensuremath{\subset}}
\newunicodechar{⊥}{\ensuremath{\perp}}
\newunicodechar{∃}{\ensuremath{\exists}}
\newunicodechar{∀}{\ensuremath{\forall}}
\newunicodechar{∛}{\ensuremath{\sqrt[3]{\,}}}
\newunicodechar{∜}{\ensuremath{\sqrt[4]{\,}}}
\newunicodechar{⃗}{\ensuremath{{}^{\to}}}
\newunicodechar{ⁿ}{\ifmmode^{n}\else\textsuperscript{\ensuremath{n}}\fi}
\newunicodechar{ˣ}{\ifmmode^{x}\else\textsuperscript{\ensuremath{x}}\fi}
\newunicodechar{ᵇ}{\ifmmode^{b}\else\textsuperscript{\ensuremath{b}}\fi}
\newunicodechar{ʳ}{\ifmmode^{r}\else\textsuperscript{\ensuremath{r}}\fi}
\newunicodechar{ᵘ}{\ifmmode^{u}\else\textsuperscript{\ensuremath{u}}\fi}
\newunicodechar{ʸ}{\ifmmode^{y}\else\textsuperscript{\ensuremath{y}}\fi}
\newunicodechar{ᵃ}{\ifmmode^{a}\else\textsuperscript{\ensuremath{a}}\fi}
\newunicodechar{ᵉ}{\ifmmode^{e}\else\textsuperscript{\ensuremath{e}}\fi}
\newunicodechar{ᵏ}{\ifmmode^{k}\else\textsuperscript{\ensuremath{k}}\fi}
\newunicodechar{ᵖ}{\ifmmode^{p}\else\textsuperscript{\ensuremath{p}}\fi}
\newunicodechar{ᴺ}{\ifmmode^{N}\else\textsuperscript{\ensuremath{N}}\fi}
\newunicodechar{ᶜ}{\ifmmode^{c}\else\textsuperscript{\ensuremath{c}}\fi}
\newunicodechar{ʰ}{\ifmmode^{h}\else\textsuperscript{\ensuremath{h}}\fi}
\newunicodechar{ᵠ}{\ifmmode^{q}\else\textsuperscript{\ensuremath{q}}\fi}
\newunicodechar{ₙ}{\ifmmode_{n}\else\textsubscript{\ensuremath{n}}\fi}
\newunicodechar{ₐ}{\ifmmode_{a}\else\textsubscript{\ensuremath{a}}\fi}
\newunicodechar{ᵢ}{\ifmmode_{i}\else\textsubscript{\ensuremath{i}}\fi}
\newunicodechar{ᵣ}{\ifmmode_{r}\else\textsubscript{\ensuremath{r}}\fi}
\newunicodechar{ₖ}{\ifmmode_{k}\else\textsubscript{\ensuremath{k}}\fi}
\newunicodechar{ᵦ}{\ifmmode_{\beta}\else\textsubscript{\ensuremath{\beta}}\fi}
\newunicodechar{ₓ}{\ifmmode_{x}\else\textsubscript{\ensuremath{x}}\fi}
\newfontfamily\popdisplay{ArchivoBlack-Regular.ttf}[Path=../../../fonts/]
\newfontfamily\poplogo{SpaceGrotesk-Bold.ttf}[Path=../../../fonts/]
\newfontfamily\popsemi{PlusJakartaSans-SemiBold.ttf}[Path=../../../fonts/]
\newfontfamily\popxbold{PlusJakartaSans-ExtraBold.ttf}[Path=../../../fonts/]
\newfontfamily\popxboldls{PlusJakartaSans-ExtraBold.ttf}[Path=../../../fonts/,LetterSpace=3.0]

\setlist[enumerate]{leftmargin=1.7em,itemsep=5pt,topsep=4pt}
\setlist[itemize]{leftmargin=1.7em,itemsep=4pt,topsep=4pt,label={\color{poporange}\textbullet}}
\setlength{\parindent}{0pt}
\setlength{\parskip}{5pt}
\renewcommand{\arraystretch}{1.25}

% pgfplots : style Pop par défaut
\pgfplotsset{
  every axis/.append style={axis line style={popink,line width=1pt},tick style={popink},
    grid style={popline,line width=0.5pt},label style={font=\small},tick label style={font=\footnotesize}},
  popcurve/.style={poporange,line width=1.8pt,no markers,smooth},
}
% Même style disponible dans un \draw TikZ simple (hors pgfplots)
\tikzset{popcurve/.style={poporange,line width=1.8pt}}
\tikzset{popgrid/.style={popline,very thin}}
\colorlet{popcurve}{poporange} % le nom du style est parfois utilisé comme couleur

% ============================================================
% Pill / squiggle
% ============================================================
\newcommand{\poppill}[3]{%
  \tikz[baseline=-0.55ex]\node[draw=popink,line width=1.2pt,rounded corners=2.6pt,%
    fill=#2,inner xsep=6.5pt,inner ysep=3pt]{%
    \popxboldls\fontsize{7}{7}\selectfont\color{#3}\MakeUppercase{#1}};%
}
\newcommand{\popsquiggle}[1]{%
  \tikz[baseline=-0.4ex]{
    \draw[#1,line width=2.6pt,line cap=round]
      plot [smooth] coordinates {(0,0.09) (0.34,0.01) (0.68,0.21) (1.02,0.01) (1.36,0.10)};
  }%
}
% Mot-clé surligné (marqueur orange sous le texte)
\newcommand{\cle}[1]{{\popxbold\color{popdark}#1}}

% ============================================================
% En-tête / pied
% ============================================================
\pagestyle{fancy}
\fancyhf{}
\renewcommand{\headrulewidth}{0pt}
\renewcommand{\footrulewidth}{0pt}
\lhead{\raisebox{-0.15ex}{\poplogo\fontsize{13}{13}\selectfont\color{popink}OnBuch\color{poporange}+}}
\rhead{\poppill{\DOCMATIERE\ \textbullet\ \DOCNIVEAU\ \textbullet\ Leçon \DOCLECON}{white}{popink}}
\lfoot{\popxbold\fontsize{7.6}{7.6}\selectfont\color{popmuted}\MakeUppercase{Cours signé OnBuch+}\ \textbullet\ {\popsemi\DOCTITRE}}
\rfoot{\tikz[baseline=-0.6ex]\node[rounded corners=8.5pt,fill=popink,minimum height=17pt,minimum width=17pt,inner xsep=5pt,inner sep=0]{\popxbold\fontsize{8}{8}\selectfont\color{popcream}\thepage\,/\,\pageref*{LastPage}};}

% ============================================================
% Titres
% ============================================================
\newcommand{\chaptitre}[2]{%
  \begin{tcolorbox}[enhanced,colback=poporange,colframe=popink,boxrule=2.6pt,arc=11pt,
      drop shadow=popink,left=15pt,right=15pt,top=12pt,bottom=11pt,before skip=3pt,after skip=18pt]
    \poppill{#2}{popink}{popcream}\par\vspace{9pt}
    {\raggedright\hyphenpenalty=10000\exhyphenpenalty=10000\popdisplay\fontsize{22}{24}\selectfont\color{popcream}\MakeUppercase{#1}\par}
  \end{tcolorbox}
}
% Section numérotée : pastille orange avec le numéro + titre Archivo
\newcounter{popsec}
\newcommand{\coursec}[1]{%
  \stepcounter{popsec}\setcounter{popsub}{0}%
  \par\vspace{16pt}\noindent
  \begin{minipage}{\linewidth}
  \tikz[baseline=-0.7ex]\node[draw=popink,line width=1.6pt,fill=poporange,rounded corners=4pt,minimum size=22pt,inner sep=0]{\popdisplay\fontsize{12}{12}\selectfont\color{popcream}\thepopsec};%
  \hspace{8pt}{\popdisplay\fontsize{15}{17}\selectfont\color{popink}\MakeUppercase{#1}}\par
  \vspace{4pt}\noindent{\color{poporange}\rule{36pt}{3.6pt}}
  \end{minipage}\par\nopagebreak\vspace{8pt}
}
\newcounter{popsub}
\newcommand{\courssub}[1]{%
  \stepcounter{popsub}%
  \par\vspace{9pt}\noindent{\popxbold\fontsize{11.5}{13}\selectfont\color{popink}{\color{poporange}\thepopsec.\thepopsub}\hspace{6pt}#1}\par\nopagebreak\vspace{4pt}
}

% ============================================================
% Boîtes pédagogiques — même recette : fond blanc, bordure épaisse colorée,
% ombre dure encre, badge plein. Titre optionnel [#1].
% ============================================================
\tcbset{popbase/.style={breakable,enhanced,colback=white,boxrule=2.3pt,arc=9pt,
  shadow={3pt}{-3pt}{0pt}{popink},left=14pt,right=14pt,top=11pt,bottom=11pt,before skip=11pt,after skip=15pt,
  fontupper=\popsemi\fontsize{10.6}{14.5}\selectfont\color{popink},
  fontlower=\popsemi\fontsize{10.6}{14.5}\selectfont\color{popink}}}
% Coche dessinée (le glyphe ✓ n'existe pas dans les polices du document)
\newcommand{\popcheck}[1]{\tikz[baseline=-0.5ex]\draw[#1,line width=1.6pt,line cap=round,line join=round](-0.13,0.0)--(-0.04,-0.09)--(0.13,0.1);}
\providecommand{\coche}{\popcheck{popgreen}}
\providecommand{\resultat}[1]{\par\smallskip\noindent\fcolorbox{popgreen}{popgreenL}{\parbox{\dimexpr\linewidth-2\fboxsep-2\fboxrule\relax}{\textbf{Résultat.} #1}}\par\smallskip}
\newcommand{\popboxhead}[3]{\poppill{#1}{#2}{white}\ifx\relax#3\relax\else\hspace{6pt}{\popxbold\fontsize{10.6}{12}\selectfont\color{popink}#3}\fi\par\vspace{7pt}}

\newtcolorbox{aretenir}[1][]{popbase,colframe=popgreen,before upper={\popboxhead{À retenir}{popgreen}{#1}}}
% Texte en chinois (document de lecture, dialogue, extrait) : cadre
% d'étude. Un titre optionnel (source, auteur) s'affiche dans l'en-tête.
\newtcolorbox{textebox}[1][]{popbase,colframe=popblue,before upper={\popboxhead{文本 · Texte}{popblue}{#1}},after upper={}}
% Marques ✓ / ✗ (forme correcte / fautive) souvent utilisées par les modèles
\providecommand{\cross}{\textcolor{poppink}{\ensuremath{\times}}}
\providecommand{\xmark}{\cross}\providecommand{\crossmark}{\cross}\providecommand{\cmark}{\ensuremath{\checkmark}}
% Ligne de réponse / justification à compléter (inventée par les modèles)
\providecommand{\justification}[1]{\par\noindent\textit{Justification~:} \dotfill\par}
% \textipa (package tipa non chargé) : la police ne couvre pas l'API -> simple passage
\providecommand{\textipa}[1]{#1}
% Soulignement (ulem non chargé) : \uline, \ul
\providecommand{\uline}[1]{\underline{#1}}\providecommand{\ul}[1]{\underline{#1}}
% \glossaire{\item ...} inventée par les modèles : liste de glossaire
\providecommand{\glossaire}[1]{\begin{itemize}#1\end{itemize}}
% \alert (beamer) inventée par les modèles : texte mis en évidence
\providecommand{\alert}[1]{\textbf{\textcolor{poppink}{#1}}}
% variantes ASCII d'ordinaux écrites en macro
\providecommand{\iere}{\textsuperscript{re}}\providecommand{\ieme}{\textsuperscript{e}}\providecommand{\ere}{\textsuperscript{re}}\providecommand{\eme}{\textsuperscript{e}}\providecommand{\er}{\textsuperscript{er}}\providecommand{\ie}{\textsuperscript{e}}
% Ordinaux français écrits en macro par les modèles : 1\ière, 3\ième
\newcommand{\ière}{\textsuperscript{re}}\newcommand{\ième}{\textsuperscript{e}}
% \exemple (inventée par les modèles) : étiquette « Exemple : »
\providecommand{\exemple}{\textbf{Exemple~:}\ }
% \square utilisable aussi hors mode math (cases à cocher)
\AtBeginDocument{\let\squareMath\square\renewcommand{\square}{\ensuremath{\squareMath}}}
% Mot ou phrase en chinois à l'intérieur d'un paragraphe français
\newcommand{\zh}[1]{\ifmmode\text{#1}\else{#1}\fi}
\let\eng\zh \let\esp\zh \let\all\zh % alias tolérés
% flèches d'intonation inventées par les modèles
\providecommand{\textsearrow}{}\renewcommand{\textsearrow}{\ensuremath{\searrow}}\providecommand{\textnearrow}{}\renewcommand{\textnearrow}{\ensuremath{\nearrow}}
% \fr{..} : mot français (inventé par les modèles)
\providecommand{\fr}[1]{\textit{#1}}
% zhbox : encadré de texte chinois inventé par les modèles
\newenvironment{zhbox}{\begin{quote}}{\end{quote}}
% \courssubsub : titre de niveau 3 inventé par les modèles
\providecommand{\courssubsub}[1]{\par\medskip\noindent\textbf{#1}\par\smallskip}
% \vs : « versus » inventé par les modèles
\providecommand{\vs}{\textit{vs}\ }
% \blank : case à compléter inventée par les modèles
\providecommand{\blank}[1][]{\underline{\hspace{1.6em}}}
\newtcolorbox{exemplebox}[1][]{popbase,colframe=poporange,before upper={\popboxhead{Exemple}{poporange}{#1}}}
\newtcolorbox{definition}[1][]{popbase,colframe=popblue,before upper={\popboxhead{Définition}{popblue}{#1}}}
\newtcolorbox{propriete}[1][]{popbase,colframe=poppurple,before upper={\popboxhead{Propriété}{poppurple}{#1}}}
\newtcolorbox{methode}[1][]{popbase,colframe=popgold,before upper={\popboxhead{Méthode}{popgold}{#1}}}
\newtcolorbox{attention}[1][]{popbase,colframe=poppink,before upper={\popboxhead{Attention}{poppink}{#1}}}
\newtcolorbox{experience}[1][]{popbase,colframe=popblue,colback=popblueL,before upper={\popboxhead{Expérience}{popblue}{#1}}}
% « Le savais-tu ? » : carte violette pleine (contexte, culture, Cameroun)
\newtcolorbox{savaistu}[1][]{popbase,colback=poppurple,colframe=popink,
  fontupper=\popsemi\fontsize{10.4}{14.2}\selectfont\color{white},
  before upper={\poppill{Le savais-tu\,?}{white}{poppurple}\ifx\relax#1\relax\else\hspace{6pt}{\popxbold\color{white}#1}\fi\par\vspace{7pt}}}
% Exercice résolu : énoncé en haut, \tcblower puis la solution sur fond crème
\newcounter{popexo}\newcounter{popexr}
\newtcolorbox{exoresolu}[1][]{popbase,colframe=poporange,colbacklower=popcream,
  segmentation style={popink,line width=1.2pt,dashed},
  before upper={\stepcounter{popexr}\popboxhead{Exercice résolu \thepopexr}{poporange}{#1}},
  before lower={\poppill{Solution}{popink}{popcream}\par\vspace{6pt}}}
% Exercice (non résolu en place) — la correction est dans la partie Corrigés
\newtcolorbox{exercice}[1][]{popbase,colframe=popink,
  before upper={\stepcounter{popexo}\popboxhead{Exercice \thepopexo}{popink}{#1}}}
% Corrigé
\newtcolorbox{corrige}[1][]{popbase,colframe=popgreen,colback=popgreenL,
  before upper={\popboxhead{Corrigé}{popgreen}{#1}}}
% Objectifs / prérequis / bilan
\newtcolorbox{objectifs}{popbase,colframe=popink,colback=poporangeL,
  before upper={\popboxhead{Objectifs de la leçon}{popink}{}}}
\newtcolorbox{prerequis}{popbase,colframe=popink,colback=popblueL,
  before upper={\popboxhead{Prérequis}{popblue}{}}}
\newtcolorbox{bilan}{popbase,colframe=popink,colback=white,boxrule=2.8pt,
  before upper={\popboxhead{Fiche bilan}{poporange}{L'essentiel à connaître pour l'examen}}}
% Figure encadrée avec légende
\newtcolorbox{popfigure}[1][]{enhanced,breakable=false,colback=white,colframe=popline,boxrule=1.4pt,arc=8pt,
  left=10pt,right=10pt,top=10pt,bottom=8pt,before skip=10pt,after skip=12pt,halign=center,
  lower separated=false,halign lower=center,
  fontlower=\popsemi\fontsize{9}{11.5}\selectfont\color{popmuted},
  #1}
\newcommand{\legende}[1]{\par\vspace{6pt}{\popsemi\fontsize{9}{11.5}\selectfont\color{popmuted}#1}}

% ============================================================
% Signature OnBuch+
% ============================================================
\newcommand{\poplogoinline}[1]{{\poplogo\fontsize{#1}{#1}\selectfont\color{popink}OnBuch\color{poporange}+}}
\newcommand{\signatureonbuch}{%
  \begin{tcolorbox}[enhanced,colback=popink,colframe=popink,boxrule=2.2pt,arc=10pt,
      drop shadow=poporange,left=14pt,right=14pt,top=10pt,bottom=10pt,before skip=0pt,after skip=0pt]
    \tikz[baseline=-0.7ex]\node[circle,fill=popgreen,draw=popcream,line width=1.3pt,minimum size=22pt,inner sep=0]{\popcheck{white}};%
    \hspace{9pt}\begin{minipage}[c]{0.62\linewidth}
      {\popxbold\fontsize{12}{13}\selectfont\color{popcream}Signé\ \textbullet\ {\poplogo OnBuch\color{poporange}+}}\par\vspace{2pt}
      {\raggedright\popsemi\fontsize{8}{10.5}\selectfont\color{popline}\MakeUppercase{Équipe pédagogique OnBuch+}\\\MakeUppercase{Conforme au programme officiel MINESEC}\par}
    \end{minipage}\hfill
    {\poplogo\fontsize{22}{22}\selectfont\color{popcream}OnBuch\color{poporange}+}
  \end{tcolorbox}%
}

% ============================================================
% Page de garde
% #1 titre · #2 matière · #3 niveau · #4 module · #5 n° leçon · #6 durée ·
% #7 description / objectifs (texte ou itemize)
% ============================================================
\newcommand{\pagegarde}[7]{%
  \thispagestyle{empty}
  \newgeometry{a4paper,margin=1.7cm,top=1.6cm,bottom=1.4cm}%
  \begin{tikzpicture}[remember picture,overlay]
    \fill[poporange!18] ([xshift=-3.2cm,yshift=-3.6cm]current page.north east) circle (4.6cm);
    \fill[poppurple!14] ([xshift=2.4cm,yshift=7.6cm]current page.south west) circle (3.8cm);
  \end{tikzpicture}%
  {\poplogo\fontsize{30}{30}\selectfont\color{popink}OnBuch\color{poporange}+}\hfill
  \raisebox{0.6ex}{\poppill{Cours signé \textbullet\ Premium}{popink}{popcream}}\par
  \vspace{1pt}\popsquiggle{poppurple}\par
  \vspace{8pt}
  \begin{tcolorbox}[enhanced,colback=poporange,colframe=popink,boxrule=2.8pt,arc=13pt,
      drop shadow=popink,left=18pt,right=18pt,top=12pt,bottom=13pt,after skip=10pt]
    \poppill{#2 \textbullet\ #3}{popink}{popcream}\hspace{5pt}\poppill{#4}{white}{popink}\par\vspace{8pt}
    {\popdisplay\fontsize{14}{14}\selectfont\color{popink}LEÇON #5}\par\vspace{4pt}
    {\raggedright\hyphenpenalty=10000\exhyphenpenalty=10000\popdisplay\fontsize{25}{27}\selectfont\color{popcream}\MakeUppercase{#1}\par}
    \vspace{8pt}
    {\popxbold\fontsize{9.5}{9.5}\selectfont\color{popink}\MakeUppercase{Durée conseillée : #6}}
  \end{tcolorbox}
  \begin{tcolorbox}[enhanced,colback=white,colframe=popink,boxrule=2.2pt,arc=10pt,
      drop shadow=popink,left=16pt,right=16pt,top=10pt,bottom=10pt,after skip=8pt]
    \poppill{Dans cette leçon}{poppurple}{white}\par\vspace{5pt}
    \popsemi\fontsize{9.3}{12.6}\selectfont\color{popink}#7
  \end{tcolorbox}
  \noindent\begin{minipage}[t]{0.487\linewidth}
    \begin{tcolorbox}[enhanced,colback=popgreen,colframe=popink,boxrule=2.2pt,arc=10pt,
        drop shadow=popink,left=11pt,right=11pt,top=10pt,bottom=10pt,height=3.1cm,valign=center]
      \tcbox[colback=white,colframe=popink,boxrule=1.3pt,arc=5pt,boxsep=0pt,nobeforeafter,
        left=3pt,right=3pt,top=3pt,bottom=3pt,valign=center]{\qrcode[height=1.9cm]{https://play.google.com/store/apps/details?id=com.luvvix.onbuch&hl=fr}}%
      \hspace{8pt}\begin{minipage}[c]{0.5\linewidth}
        {\popdisplay\fontsize{11}{12.5}\selectfont\color{white}\MakeUppercase{Télécharge OnBuch}}\par\vspace{4pt}
        {\raggedright\popsemi\fontsize{8.4}{11}\selectfont\color{white}Gratuit \textbullet\ Google Play\\Tuteur IA, annales, concours, résultats\par}
      \end{minipage}
    \end{tcolorbox}
  \end{minipage}\hfill
  \begin{minipage}[t]{0.487\linewidth}
    \begin{tcolorbox}[enhanced,colback=poppurple,colframe=popink,boxrule=2.2pt,arc=10pt,
        drop shadow=popink,left=11pt,right=11pt,top=10pt,bottom=10pt,height=3.1cm,valign=center]
      \tikz[baseline=-0.7ex]\node[circle,fill=white,draw=popink,line width=1.3pt,minimum size=1.6cm,inner sep=0]{\popdisplay\fontsize{24}{24}\selectfont\color{poppurple}+};%
      \hspace{8pt}\begin{minipage}[c]{0.6\linewidth}
        {\popdisplay\fontsize{11}{12.5}\selectfont\color{white}\MakeUppercase{Passe à OnBuch+}}\par\vspace{4pt}
        {\raggedright\popsemi\fontsize{8.4}{11}\selectfont\color{white}Tous les cours signés, Tuteur IA illimité, fiches PDF — onglet OnBuch+ de l'app\par}
      \end{minipage}
    \end{tcolorbox}
  \end{minipage}
  \vspace{8pt}
  \popcontenu
  \vfill
  \signatureonbuch
  \restoregeometry
  \newpage
}

% Rangée « Ce cours contient » (page de garde)
\newcommand{\poptile}[3]{% #1 couleur #2 gros texte #3 légende
  \begin{tcolorbox}[enhanced,colback=#1,colframe=popink,boxrule=2pt,arc=9pt,shadow={2.5pt}{-2.5pt}{0pt}{popink},
      left=6pt,right=6pt,top=9pt,bottom=9pt,halign=center,valign=center,height=1.75cm,width=\linewidth,nobeforeafter]
    {\popdisplay\fontsize{12.5}{13}\selectfont\color{white}\MakeUppercase{#2}}\par\vspace{3pt}
    {\centering\popsemi\fontsize{7.8}{9.5}\selectfont\color{white}#3\par}
  \end{tcolorbox}}
\newcommand{\popcontenu}{%
  \par\noindent{\popxboldls\fontsize{7.5}{7.5}\selectfont\color{popmuted}CE COURS SIGNÉ CONTIENT}\par\vspace{7pt}
  \noindent\begin{minipage}[t]{0.235\linewidth}\poptile{popblue}{Cours}{complet, illustré, pas à pas}\end{minipage}\hfill
  \begin{minipage}[t]{0.235\linewidth}\poptile{poporange}{Méthodes}{et exercices résolus}\end{minipage}\hfill
  \begin{minipage}[t]{0.235\linewidth}\poptile{poppink}{Exercices}{type examen + corrigés}\end{minipage}\hfill
  \begin{minipage}[t]{0.235\linewidth}\poptile{popgreen}{Fiche bilan}{l'essentiel à retenir}\end{minipage}\par}

% Sommaire
\newtcolorbox{sommaire}{enhanced,colback=white,colframe=popink,boxrule=2.2pt,arc=10pt,
  drop shadow=popink,left=18pt,right=18pt,top=13pt,bottom=13pt,before skip=6pt,after skip=20pt,
  before upper={\poppill{Sommaire}{poppurple}{white}\par\vspace{9pt}\popsemi\fontsize{10.6}{14.5}\selectfont\color{popink}}}

% Signature de fin de cours — poussée tout en bas de la dernière page
\newcommand{\findecours}{%
  \par\vfill\nopagebreak
  \begin{center}\popsquiggle{poporange}\end{center}\vspace{4pt}
  \signatureonbuch
}
\AtBeginDocument{\def\llbracket{\mathopen{[\mkern-3.5mu[}}\def\rrbracket{\mathclose{]\mkern-3.5mu]}}} % intervalles d'entiers (glyphe ⟦ absent de la police)
% Glyphes absents de Fira Math : redéfinis à partir de glyphes disponibles
\AtBeginDocument{%
  \def\setminus{\mathbin{\backslash}}\let\smallsetminus\setminus
  \def\vdots{\mathord{\vbox{\baselineskip3.2pt\lineskiplimit0pt\kern4pt\hbox{.}\hbox{.}\hbox{.}}}}%
  \def\ddots{\mathinner{\mkern1mu\raise6.5pt\hbox{.}\mkern2mu\raise3.6pt\hbox{.}\mkern2mu\raise0.7pt\hbox{.}\mkern1mu}}%
  \def\triangle{\mathord{\scalebox{0.9}{$\symup{\Delta}$}}}%
  \def\nabla{\mathord{\raisebox{\depth}{\scalebox{1}[-1]{$\symup{\Delta}$}}}}%
  \def\checkmark{\popcheck{popdark}}%
  \def\star{\mathbin{\ast}}\def\bigstar{\mathord{\ast}}%
  \def\diamond{\mathbin{\scalebox{0.6}{$\Diamond$}}}%
  \def\bigcup{\mathop{\mathchoice{\raisebox{-0.3em}{\scalebox{1.7}{$\cup$}}}{\scalebox{1.25}{$\cup$}}{\cup}{\cup}}\displaylimits}%
  \def\bigcap{\mathop{\mathchoice{\raisebox{-0.3em}{\scalebox{1.7}{$\cap$}}}{\scalebox{1.25}{$\cap$}}{\cap}{\cap}}\displaylimits}%
  \def\lesseqqgtr{\mathrel{\lessgtr}}\def\bigoplus{\mathop{\oplus}\displaylimits}%
}
\providecommand{\ssi}{si et seulement si} % abréviation fréquente
\AtBeginDocument{\providecommand{\wideparen}{\overparen}} % arc de cercle

%% FIGFIX-BEGIN v5 — correctifs globaux des figures (docs/onbuch-generation/tools/figfix)
\usepackage{adjustbox}\usepackage{etoolbox}\tcbuselibrary{raster}
% toute figure TikZ/pgfplots plus large que la ligne est réduite à la largeur disponible
\BeforeBeginEnvironment{tikzpicture}{\begin{adjustbox}{max width=\linewidth}}
\AfterEndEnvironment{tikzpicture}{\end{adjustbox}}
% légendes pgfplots lisibles par-dessus les courbes
\pgfplotsset{every axis/.append style={legend style={fill=white,fill opacity=0.92,text opacity=1,draw=black!15,inner sep=3pt}}}
% étiquettes d'arêtes (syntaxe edge["..."]) avec fond blanc
\tikzset{every edge quotes/.append style={fill=white,inner sep=1.2pt}}
%% FIGFIX-END
\begin{document}

\pagegarde{Leçon 29 — Éléments socioculturels : différence entre la monnaie chinoise et la monnaie camerounaise}{Chinois}{Tle A}{Module 4 : Activités économiques et monde du travail}{ZH29}{2 h}{Cette leçon socioculturelle compare de façon factuelle et respectueuse la monnaie chinoise (le renminbi/yuan) et la monnaie camerounaise (le franc CFA) : histoire, émetteurs, billets, usages et paiement mobile. Elle fournit le vocabulaire chinois associé et prépare à un mini-exposé comparatif.
\vspace{4pt}
\begin{multicols}{2}\small
\begin{itemize}
\item À la fin de cette leçon, je sais nommer la monnaie chinoise (人民币, 元, 角, 分) et la monnaie camerounaise (非洲法郎) en chinois.
\item Je sais présenter factuellement qui émet chaque monnaie (中国人民银行 / 中非国家银行 BEAC).
\item Je sais décrire les billets et pièces des deux pays et ce qu'ils représentent.
\item Je sais comparer les usages : paiement en espèces au Cameroun, paiement mobile très répandu en Chine.
\item Je sais employer le vocabulaire socioculturel du thème (钱, 银行, 换钱, 现金, 手机支付).
\item Je sais exprimer un point de vue simple en chinois (我觉得..., 因为...).
\end{itemize}\end{multicols}}

\begin{sommaire}
\begin{multicols}{2}
\begin{enumerate}
\item Découverte : deux monnaies, deux réalités
\item Glossaire et écriture des caractères
\item Compréhension du texte
\item Faits comparés Cameroun / Chine
\item Exprimer son point de vue et débattre
\item Production guidée : mini-exposé comparatif
\item Activité d'intégration
\item Méthodes et exercices résolus
\item Exercices
\item Corrigés des exercices
\item Fiche bilan
\end{enumerate}
\end{multicols}
\end{sommaire}

\begin{objectifs}
\begin{itemize}
\item À la fin de cette leçon, je sais nommer la monnaie chinoise (人民币, 元, 角, 分) et la monnaie camerounaise (非洲法郎) en chinois.
\item Je sais présenter factuellement qui émet chaque monnaie (中国人民银行 / 中非国家银行 BEAC).
\item Je sais décrire les billets et pièces des deux pays et ce qu'ils représentent.
\item Je sais comparer les usages : paiement en espèces au Cameroun, paiement mobile très répandu en Chine.
\item Je sais employer le vocabulaire socioculturel du thème (钱, 银行, 换钱, 现金, 手机支付).
\item Je sais exprimer un point de vue simple en chinois (我觉得..., 因为...).
\item Je sais poser et répondre à des questions de réflexion sur les deux monnaies.
\item Je sais préparer et présenter un mini-exposé comparatif Cameroun/Chine.
\end{itemize}
\end{objectifs}

\begin{prerequis}
\begin{itemize}
\item Les nombres et les prix en chinois (块, 毛, 分) vus en Première
\item La structure d'opinion 我觉得 + phrase
\item Le vocabulaire des pays et nationalités (中国, 喀麦隆)
\item Les verbes courants : 有, 用, 买, 换
\item La comparaison simple avec 比 (révision utile)
\end{itemize}
\end{prerequis}

\newpage

\coursec{Découverte : deux monnaies, deux réalités}

À Yaoundé, Amina paie son taxi avec des billets de 1000 FCFA ornés des symboles de la CEMAC. À Pékin, sa correspondante Li Na règle ses achats avec son téléphone, via WeChat Pay ou Alipay, en yuans (\zh{人民币}). Quand Li Na vient au Cameroun, elle doit changer ses yuans contre des francs CFA. Qui émet ces monnaies ? Que représentent-elles ? Comment dit-on « argent », « monnaie », « banque », « taux de change » en chinois ? Cette leçon nous fait voyager dans le monde des monnaies chinoise et camerounaise.

\textbf{Problématique :} en quoi la monnaie chinoise et la monnaie camerounaise se ressemblent-elles et se distinguent-elles, et comment en parler en chinois pour comparer, échanger et argumenter ?

\courssub{Lecture du texte support}

Lisons d'abord un court texte qui présente les deux monnaies. Lis-le une première fois en silence, puis une deuxième fois à voix haute après le professeur.

\begin{textebox}[中国和喀麦隆的钱]
中国和喀麦隆的钱\\
\emph{Zhōngguó hé Kāmàilóng de qián}\\
« L'argent de la Chine et du Cameroun »\\[4pt]
中国用的钱叫人民币，单位是元。\\
\emph{Zhōngguó yòng de qián jiào rénmínbì, dānwèi shì yuán.}\\
« La monnaie utilisée en Chine s'appelle le renminbi, son unité est le yuan. »\\[4pt]
喀麦隆用的钱是非洲法郎。\\
\emph{Kāmàilóng yòng de qián shì Fēizhōu fǎláng.}\\
« Le Cameroun utilise le franc CFA. »\\[4pt]
中国人民银行发行人民币，中非国家银行发行非洲法郎。\\
\emph{Zhōngguó Rénmín Yínháng fāxíng rénmínbì, Zhōngfēi Guójiā Yínháng fāxíng Fēizhōu fǎláng.}\\
« La Banque populaire de Chine émet le renminbi, la Banque des États de l'Afrique centrale émet le franc CFA. »\\[4pt]
在中国，很多人用手机付钱；在喀麦隆，人们常常用现金。\\
\emph{Zài Zhōngguó, hěn duō rén yòng shǒujī fù qián ; zài Kāmàilóng, rénmen chángcháng yòng xiànjīn.}\\
« En Chine, beaucoup de gens paient avec leur téléphone ; au Cameroun, on utilise souvent les espèces. »
\end{textebox}

\courssub{Repérage global du texte}

Avant d'expliquer les mots nouveaux, vérifions la compréhension générale. Réponds aux questions suivantes, en français puis en chinois :

\begin{enumerate}
\item De quoi parle le texte ? (\emph{Indice : regarde le titre.})
\item Quels sont les deux pays cités ? Quelles sont les deux monnaies citées ?
\item Quelles sont les deux banques mentionnées ?
\item Quelle différence de pratique de paiement le texte signale-t-il entre la Chine et le Cameroun ?
\end{enumerate}

\begin{exoresolu}[Repérage global : réponses guidées]
Énoncé : réponds aux quatre questions de repérage ci-dessus en une phrase chinoise simple chacune.
\tcblower
Solution détaillée :
\begin{enumerate}
\item Le texte parle de l'argent : \zh{这个课文说钱。} \emph{Zhège kèwén shuō qián.} « Ce texte parle de l'argent. »
\item Les deux pays sont la Chine et le Cameroun : \zh{两个国家是中国和喀麦隆。} \emph{Liǎng ge guójiā shì Zhōngguó hé Kāmàilóng.} « Les deux pays sont la Chine et le Cameroun. »
\item Les deux banques : \zh{中国人民银行和中非国家银行。} \emph{Zhōngguó Rénmín Yínháng hé Zhōngfēi Guójiā Yínháng.} « La Banque populaire de Chine et la Banque des États de l'Afrique centrale. »
\item La différence de paiement : \zh{在中国，很多人用手机付钱；在喀麦隆，人们常常用现金。} \emph{Zài Zhōngguó, hěn duō rén yòng shǒujī fù qián ; zài Kāmàilóng, rénmen chángcháng yòng xiànjīn.} « En Chine, beaucoup paient par téléphone ; au Cameroun, on paie souvent en espèces. »
\end{enumerate}
\end{exoresolu}

\courssub{Observation : deux monnaies, deux réalités}

Observons les phrases clés du texte avant d'énoncer les faits.

\begin{exemplebox}[Observer avant de comprendre]
\begin{itemize}
\item \zh{中国用的钱叫人民币。} \emph{Zhōngguó yòng de qián jiào rénmínbì.} « La monnaie qu'utilise la Chine s'appelle le renminbi. »
\item \zh{单位是元。} \emph{Dānwèi shì yuán.} « L'unité est le yuan. »
\item \zh{中国人民银行发行人民币。} \emph{Zhōngguó Rénmín Yínháng fāxíng rénmínbì.} « La Banque populaire de Chine émet le renminbi. »
\item \zh{中非国家银行发行非洲法郎。} \emph{Zhōngfēi Guójiā Yínháng fāxíng Fēizhōu fǎláng.} « La Banque des États de l'Afrique centrale émet le franc CFA. »
\end{itemize}
On remarque la structure \zh{……用的钱叫……} \emph{……yòng de qián jiào……} « la monnaie que X utilise s'appelle … » : le verbe \zh{叫} (jiào, « s'appeler ») introduit le nom officiel de la monnaie, exactement comme « s'appeler » en français. La particule \zh{的} (de) nominalise le verbe \zh{用} (yòng, « utiliser ») : \zh{用的钱} = « la monnaie qu'on utilise / la monnaie utilisée ».
\end{exemplebox}

\begin{definition}[人民币 : la monnaie du peuple]
\zh{人民币} \emph{rénmínbì} signifie littéralement « monnaie du peuple » (\zh{人民} « le peuple » + \zh{币} « monnaie »). C'est le \cle{nom officiel} de la monnaie chinoise. Son \cle{unité de base} est le \zh{元} \emph{yuán}. On dit donc : \zh{一块钱} \emph{yí kuài qián} « un yuan » dans la langue courante (\zh{块} kuài est le classificateur familier de l'argent), et \zh{一元} \emph{yī yuán} dans la langue formelle, sur les billets et dans l'écrit comptable.
\end{definition}

\begin{propriete}[Présenter une monnaie : faits et structures]
Pour présenter une monnaie en chinois, on utilise trois structures simples :
\begin{itemize}
\item Nom de la monnaie : \cle{Pays + 用的钱叫 + nom} — \zh{喀麦隆用的钱叫非洲法郎。} \emph{Kāmàilóng yòng de qián jiào Fēizhōu fǎláng.} « La monnaie du Cameroun s'appelle le franc CFA. »
\item Unité : \cle{单位是 + unité} — \zh{单位是元。} \emph{Dānwèi shì yuán.} « L'unité est le yuan. »
\item Émission : \cle{Banque + 发行 + monnaie} — \zh{中国人民银行发行人民币。} \emph{Zhōngguó Rénmín Yínháng fāxíng rénmínbì.} « La Banque populaire de Chine émet le renminbi. » Le verbe \zh{发行} (fāxíng, « émettre, éditer ») s'emploie aussi pour les journaux (\zh{发行报纸}) et les timbres (\zh{发行邮票}).
\end{itemize}
\end{propriete}

\courssub{Faits présentés : la Chine et le Cameroun}

\begin{aretenir}[Le renminbi (人民币)]
\begin{itemize}
\item Nom officiel : \zh{人民币} \emph{rénmínbì} « monnaie du peuple ».
\item Unité : le \zh{元} \emph{yuán} (symbole ¥ ; abréviations internationales RMB et CNY). Sous-unités : 1 yuán = 10 \zh{角} (jiǎo) = 100 \zh{分} (fēn).
\item Émetteur : la Banque populaire de Chine, \zh{中国人民银行} \emph{Zhōngguó Rénmín Yínháng}, banque centrale du pays.
\item Usage : en Chine, le paiement par téléphone mobile (\zh{手机支付} shǒujī zhīfù, via WeChat Pay / Alipay) est hégémonique, même pour les tout petits achats (marchés, taxis, distributeurs automatiques).
\end{itemize}
\end{aretenir}

\begin{aretenir}[Le franc CFA (非洲法郎)]
\begin{itemize}
\item Nom : \zh{非洲法郎} \emph{Fēizhōu fǎláng} « franc d'Afrique », c'est-à-dire le franc CFA (Coopération Financière en Afrique Centrale).
\item Usage : il est utilisé au Cameroun et dans les 5 autres pays de la zone CEMAC (Communauté Économique et Monétaire de l'Afrique Centrale) : Tchad, RCA, Congo, Gabon, Guinée équatoriale.
\item Émetteur : la Banque des États de l'Afrique centrale (BEAC), \zh{中非国家银行} \emph{Zhōngfēi Guójiā Yínháng}, dont le siège est à Yaoundé.
\item Abréviations : FCFA en français courant ; XAF est le code ISO 4217 international.
\item Usage : au Cameroun, les paiements en espèces (\zh{现金} xiànjīn) restent très fréquents, bien que le paiement mobile (\zh{移动支付} yídòng zhīfù, ex. Mobile Money, Orange Money, Moov Money) se développe rapidement.
\end{itemize}
\end{aretenir}

\begin{attention}[人民币 ou 元 ? Ne confonds pas la monnaie et son unité]
Les francophones confondent souvent \zh{人民币} et \zh{元}. Retiens la comparaison : \zh{人民币} est le \emph{nom de la monnaie}, comme « franc CFA » ou « euro » ; \zh{元} est l'\emph{unité de compte}, comme « franc » ou « centime » (mais ici unité de base). On dit : « La monnaie chinoise s'appelle le renminbi » (\zh{中国的钱叫人民币。}) mais « Cela coûte dix yuans » (\zh{这个十块钱。} \emph{Zhège shí kuài qián.} ou \zh{这个十元。} \emph{Zhège shí yuán.}). \textbf{Erreur typique :} dire \zh{十人民币} — c'est faux, car on compte en \zh{元} / \zh{块}, pas en \zh{人民币}. De même, on ne dit pas « dix francs CFA » pour 10 000 FCFA : on dit « dix mille francs ».
\end{attention}

\begin{savaistu}[Que représentent les billets ?]
Les billets chinois (5e série, 2019) portent tous le portrait de Mao Zedong au recto. Au verso, des paysages célèbres : le site des Trois Gorges (10 yuans), le mont Tai (5 yuans), la salle de l'Assemblée du peuple (100 yuans), etc. Les billets de la BEAC, eux, représentent la faune (antilope, éléphant), l'agriculture (cacao, café, coton) et les infrastructures (barrage, pont) de l'Afrique centrale. Ainsi, un billet de banque est aussi une petite image du pays : il raconte son histoire, sa nature et ses réalisations. Regarde un billet de 1000 FCFA : que vois-tu ?
\end{savaistu}

\courssub{Vocabulaire de découverte}

Voici les mots clés de cette première partie. Le professeur les lit d'abord ; la classe répète ensuite en chœur, en faisant attention aux tons.

\begin{center}
\begin{tabularx}{\linewidth}{|X|X|X|X|}
\hline
\rowcolor{popblueL} Caractères & Pinyin & Nature & Traduction \\
\hline
钱 & qián & nom & argent, monnaie \\
\hline
人民币 & rénmínbì & nom & renminbi, « monnaie du peuple » \\
\hline
元 & yuán & nom / classif. & yuan (unité monétaire formelle) \\
\hline
块 & kuài & classif. & yuan (langue courante) \\
\hline
角 & jiǎo & nom / classif. & jiao (1/10 de yuan) \\
\hline
分 & fēn & nom / classif. & fen (1/100 de yuan) \\
\hline
非洲法郎 & Fēizhōu fǎláng & nom & franc CFA \\
\hline
银行 & yínháng & nom & banque \\
\hline
中国人民银行 & Zhōngguó Rénmín Yínháng & nom propre & Banque populaire de Chine \\
\hline
中非国家银行 & Zhōngfēi Guójiā Yínháng & nom propre & Banque des États de l'Afrique centrale (BEAC) \\
\hline
发行 & fāxíng & verbe & émettre (monnaie, timbre, journal) \\
\hline
单位 & dānwèi & nom & unité (de mesure, monétaire) \\
\hline
现金 & xiànjīn & nom & espèces, argent liquide \\
\hline
付钱 & fù qián & verbe + compl. & payer (litt. « donner de l'argent ») \\
\hline
手机 & shǒujī & nom & téléphone portable \\
\hline
换钱 & huàn qián & verbe + compl. & changer de l'argent \\
\hline
汇率 & huìlǜ & nom & taux de change \\
\hline
支付 & zhīfù & verbe / nom & payer / paiement \\
\hline
\end{tabularx}
\end{center}

\courssub{Écriture des caractères : ordre des traits et clés}

\begin{experience}[Apprendre à écrire : 钱、元、银、行]
Le professeur montre l'ordre des traits au tableau ; les élèves s'exercent sur papier quadrillé (纸格子纸), en respectant les règles : haut → bas, gauche → droite, horizontal avant vertical, trait central avant les ailes, cadre avant contenu, fermeture en dernier.
\begin{itemize}
\item \zh{钱} (qián, 10 traits) : clé \zh{钅} (jīn, « or/métal », 5 traits) à gauche + \zh{戈} (gē, « hallebarde ») à droite. \emph{Ordre :} traits de la clé or (point, horizontale, verticale, horizontale, point bas) puis traits de 戈 (horizontale, point, horizontale, verticale, horizontale crochue).
\item \zh{元} (yuán, 4 traits) : clé \zh{一} (yī, « un ») en haut + \zh{儿} (ér, « enfant/pied ») en bas. \emph{Ordre :} horizontale (一), puis les 3 traits de 儿 (horizontale courte, verticale crochue, horizontale longue).
\item \zh{银} (yín, 14 traits) : clé \zh{钅} (or/métal) + phonétique \zh{艮} (gèn). \emph{Ordre :} clé or (5 traits) puis 艮 (horizontale, verticale, horizontale, horizontale, verticale, horizontale).
\item \zh{行} (háng/xíng, 6 traits) : clé \zh{彳} (chì, « pas », 3 traits) à gauche + \zh{亍} (chù, « s'arrêter », 3 traits) à droite. \emph{Ordre :} 彳 (point, horizontale, verticale) puis 亍 (horizontale, verticale, horizontale).
\end{itemize}
\end{experience}

\courssub{Carte mentale : deux monnaies autour de 钱}

\begin{popfigure}
\begin{tikzpicture}[mindmap, grow cyclic,
  every node/.style={concept, align=center, font=\small},
  root concept/.append style={concept color=popblue, font=\large},
  level 1/.append style={level distance=3.6cm, sibling angle=180},
  level 2/.append style={level distance=2.6cm, sibling angle=60}]
\node[root concept, align=center]{钱\\ qián\\ (monnaie)}
  child[concept color=poporange] { node[align=center]{中国\\ Zhōngguó}
    child[concept color=poporangeL] { node[align=center]{人民币\\ rénmínbì\\ 元 yuán / 块 kuài} }
    child[concept color=poporangeL] { node[align=center]{中国人民银行\\ Banque populaire\\ de Chine} }
    child[concept color=poporangeL] { node[align=center]{Paiement mobile\\ 手机支付} }
  }
  child[concept color=popgreen] { node[align=center]{喀麦隆\\ Kāmàilóng}
    child[concept color=popgreenL] { node[align=center]{非洲法郎\\ Fēizhōu fǎláng\\ FCFA / XAF} }
    child[concept color=popgreenL] { node[align=center]{中非国家银行\\ BEAC (Yaoundé)} }
    child[concept color=popgreenL] { node[align=center]{Espèces / Mobile Money\\ 现金 / 移动支付} }
  };
\end{tikzpicture}
\legende{Carte mentale : les deux monnaies autour du mot 钱 (qián, « argent, monnaie »).}
\end{popfigure}

\courssub{Tableau comparatif à compléter au fil de la leçon}

Recopie ce tableau dans ton cahier et complète-le au fur et à mesure de la leçon. Les deux premières lignes sont déjà remplies pour te guider.

\begin{popfigure}
\begin{tikzpicture}
\draw[fill=popblue, draw=popblue] (0,0) rectangle (5,0.9);
\draw[fill=popblue, draw=popblue] (5,0) rectangle (10,0.9);
\node[white, font=\bfseries] at (2.5,0.45) {中国 Chine};
\node[white, font=\bfseries] at (7.5,0.45) {喀麦隆 Cameroun};
\draw[fill=poporangeL, draw=popline] (0,-0.9) rectangle (5,0);
\draw[fill=poporangeL, draw=popline] (5,-0.9) rectangle (10,0);
\node[align=center, font=\small] at (2.5,-0.45) {人民币 rénmínbì\\ (元 yuán, ¥ / RMB / CNY)};
\node[align=center, font=\small] at (7.5,-0.45) {非洲法郎 Fēizhōu fǎláng\\ (FCFA / XAF)};
\draw[fill=popgreenL, draw=popline] (0,-1.8) rectangle (5,-0.9);
\draw[fill=popgreenL, draw=popline] (5,-1.8) rectangle (10,-0.9);
\node[align=center, font=\small] at (2.5,-1.35) {中国人民银行\\ Banque populaire de Chine};
\node[align=center, font=\small] at (7.5,-1.35) {中非国家银行\\ BEAC (siège à Yaoundé)};
\draw[fill=popcream, draw=popline] (0,-2.7) rectangle (5,-1.8);
\draw[fill=popcream, draw=popline] (5,-2.7) rectangle (10,-1.8);
\node[align=center, font=\small, text=popmuted] at (2.5,-2.25) {Paiement dominant :\\ 手机支付 (WeChat/Alipay)};
\node[align=center, font=\small, text=popmuted] at (7.5,-2.25) {Paiement dominant :\\ 现金 + Mobile Money};
\draw[fill=popcream, draw=popline] (0,-3.6) rectangle (5,-2.7);
\draw[fill=popcream, draw=popline] (5,-3.6) rectangle (10,-2.7);
\node[align=center, font=\small, text=popmuted] at (2.5,-3.15) {Billets : Mao Zedong +\\ paysages (3 Gorges, etc.)};
\node[align=center, font=\small, text=popmuted] at (7.5,-3.15) {Billets : Faune, agriculture,\\ infrastructures CEMAC};
\draw[fill=popcream, draw=popline] (0,-4.5) rectangle (5,-3.6);
\draw[fill=popcream, draw=popline] (5,-4.5) rectangle (10,-3.6);
\node[align=center, font=\small, text=popmuted] at (2.5,-4.05) {Sous-unités : 角、分};
\node[align=center, font=\small, text=popmuted] at (7.5,-4.05) {Pas de sous-unité courante\\ (1 FCFA = unité de base)};
\end{tikzpicture}
\legende{Tableau comparatif Chine / Cameroun : cases complétées.}
\end{popfigure}

\courssub{Première écoute et répétition chorale}

\begin{experience}[Écoute et répétition des mots clés]
\begin{enumerate}
\item \textbf{Écoute silencieuse.} Le professeur lit le texte une première fois. Les élèves écoutent sans répéter, le doigt sur les mots.
\item \textbf{Répétition chorale.} La classe répète en chœur chaque mot clé après le professeur, deux fois : \zh{钱} qián, \zh{人民币} rénmínbì, \zh{元} yuán, \zh{银行} yínháng, \zh{非洲法郎} Fēizhōu fǎláng, \zh{发行} fāxíng, \zh{现金} xiànjīn, \zh{付钱} fù qián, \zh{汇率} huìlǜ, \zh{支付} zhīfù.
\item \textbf{Attention aux tons.}
      \begin{itemize}
      \item \zh{钱} qián (2\up{e} ton, montant), \zh{银} yín (2\up{e} ton), \zh{行} háng (2\up{e} ton) — attention : \zh{行} se prononce \emph{xíng} (2\up{e} ton) quand il signifie « marcher / aller / OK », mais \emph{háng} (2\up{e} ton) pour « banque / rang / métier ».
      \item \zh{法} fǎ (3\up{e} ton, descendant-montant), \zh{郎} láng (2\up{e} ton).
      \item \zh{汇} huì (4\up{e} ton, descendant), \zh{率} lǜ (4\up{e} ton).
      \item \zh{支} zhī (1\up{e} ton, haut), \zh{付} fù (4\up{e} ton).
      \end{itemize}
\item \textbf{Lecture en chaîne.} Chaque rangée lit une phrase du texte à voix haute ; les autres vérifient les tons.
\item \textbf{Mini-question.} Le professeur montre un billet (image) et demande : \zh{这是什么钱？} \emph{Zhè shì shénme qián ?} « Quelle monnaie est-ce ? » — Réponse attendue : \zh{这是人民币。} ou \zh{这是非洲法郎。}
\end{enumerate}
\end{experience}

\begin{exercice}[À toi de jouer : vrai ou faux ?]
Réponds par vrai (\zh{对} duì) ou faux (\zh{不对} bù duì), puis corrige les phrases fausses en chinois :
\begin{enumerate}
\item \zh{中国用的钱叫非洲法郎。}
\item \zh{中国人民银行发行人民币。}
\item \zh{在喀麦隆，人们常常用现金。}
\item \zh{人民币的单位是法郎。}
\item \zh{中国的钱的单位是元和角。}
\end{enumerate}
\end{exercice}

\begin{corrige}[Exercice « vrai ou faux »]
\begin{enumerate}
\item Faux. La monnaie de la Chine est le renminbi : \zh{中国用的钱叫人民币。}
\item Vrai. \zh{对。中国人民银行发行人民币。}
\item Vrai. \zh{对。在喀麦隆，人们常常用现金。}
\item Faux. L'unité du renminbi est le yuan : \zh{人民币的单位是元。}
\item Vrai. \zh{对。人民币的单位是元和角。} (On peut ajouter : \zh{还有分。})
\end{enumerate}
\end{corrige}

\courssub{Approfondissement : texte comparatif détaillé}

Pour aller plus loin dans la comparaison, lisons un texte un peu plus long qui introduit le vocabulaire du change et des habitudes de paiement.

\begin{textebox}[换钱和付钱 : 两个国家的不同]
换钱和付钱：两个国家的不同\\
\emph{Huàn qián hé fù qián : liǎng ge guójiā de bùtóng}\\
« Changer et payer : les différences entre deux pays »\\[4pt]
李娜是中国学生，阿米娜是喀麦隆学生。\\
\emph{Lǐ Nà shì Zhōngguó xuéshēng, Āmǐnà shì Kāmàilóng xuéshēng.}\\
« Li Na est étudiante chinoise, Amina est étudiante camerounaise. »\\[4pt]
李娜来喀麦隆旅游，她要换钱。\\
\emph{Lǐ Nà lái Kāmàilóng lǚyóu, tā yào huàn qián.}\\
« Li Na vient au Cameroun en voyage, elle doit changer de l'argent. »\\[4pt]
她去银行，问汇率。今天的汇率是 1 元 = 80 非洲法郎。\\
\emph{Tā qù yínháng, wèn huìlǜ. Jīntiān de huìlǜ shì yī yuán = bāshí Fēizhōu fǎláng.}\\
« Elle va à la banque, demande le taux de change. Le taux d'aujourd'hui est 1 yuan = 80 francs CFA. »\\[4pt]
李娜用手机付钱，很方便。阿米娜用现金，也习惯了。\\
\emph{Lǐ Nà yòng shǒujī fù qián, hěn fāngbiàn. Āmǐnà yòng xiànjīn, yě xíguàn le.}\\
« Li Na paie par téléphone, c'est très pratique. Amina paie en espèces, elle y est habituée aussi. »\\[4pt]
两个国家的钱不一样，但都方便生活。\\
\emph{Liǎng ge guójiā de qián bù yīyàng, dàn dōu fāngbiàn shēnghuó.}\\
« Les monnaies des deux pays sont différentes, mais toutes deux facilitent la vie. »
\end{textebox}

\begin{exercice}[Compréhension du texte approfondi]
Réponds en chinois (phrase complète) :
\begin{enumerate}
\item 李娜是谁？她要做什么？
\item 今天的汇率是多少？(用中文说)
\item 李娜怎么付钱？阿米娜怎么付钱？
\item 最后一句怎么说？翻译成法语。
\end{enumerate}
\end{exercice}

\begin{corrige}[Exercice compréhension approfondie]
\begin{enumerate}
\item \zh{李娜是中国学生，她要换钱。} (Lǐ Nà shì Zhōngguó xuéshēng, tā yào huàn qián.)
\item \zh{今天的汇率是一元等于八十非洲法郎。} (Jīntiān de huìlǜ shì yī yuán děngyú bāshí Fēizhōu fǎláng.)
\item \zh{李娜用手机付钱，阿米娜用现金付钱。} (Lǐ Nà yòng shǒujī fù qián, Āmǐnà yòng xiànjīn fù qián.)
\item « Les monnaies des deux pays sont différentes, mais toutes deux facilitent la vie. »
\end{enumerate}
\end{corrige}

\courssub{Expression orale : débat et échange d'opinions}

\begin{experience}[Débat : Espèces ou paiement mobile ?]
\textbf{Consigne :} En binôme, préparez un dialogue de 1 minute. L'un joue un Chinois (habitué au paiement mobile), l'autre un Camerounais (habitué aux espèces / Mobile Money). Utilisez les structures et le vocabulaire vus.
\textbf{Mots-outils :} \zh{我觉得……} (wǒ juéde, « je pense que »), \zh{因为……所以……} (yīnwèi… suǒyǐ, « parce que… donc »), \zh{比较} (bǐjiào, « relativement »), \zh{方便} (fāngbiàn, « pratique »), \zh{安全} (ānquán, « sûr/sécurisé »), \zh{习惯} (xíguàn, « avoir l'habitude »).
\textbf{Exemple d'amorce :}
A: \zh{我觉得用手机付钱比较方便，因为不用带现金。} (Je trouve le paiement mobile plus pratique, car pas besoin de porter du cash.)
B: \zh{可是在喀麦隆，有时候网络不好，用现金更安全。} (Mais au Cameroun, parfois le réseau est mauvais, l'espèce est plus sûre.)
\textbf{Déroulement :} 5 min préparation → 3-4 binômes passent devant la classe → correction collective des tons et structures.
\end{experience}

\courssub{Méthode : préparer un mini-exposé comparatif}

\begin{methode}[Comment structurer un mini-exposé de 2 minutes en chinois]
\begin{enumerate}
\item \textbf{Introduction (1 phrase) :} Présenter le thème. \zh{今天我比较中国和喀麦隆的钱。} (Aujourd'hui je compare les monnaies chinoise et camerounaise.)
\item \textbf{Point 1 : Nom et unité (2 phrases) :} \zh{中国用人民币，单位是元。喀麦隆用非洲法郎，单位是法郎。}
\item \textbf{Point 2 : Émission (1 phrase) :} \zh{中国人民银行发行人民币，中非国家银行发行非洲法郎。}
\item \textbf{Point 3 : Paiement (2 phrases) :} \zh{在中国，大家常用手机支付。在喀麦隆，大家常用现金和 Mobile Money。}
\item \textbf{Point 4 : Ton avis personnel (1 phrase) :} \zh{我觉得手机支付很方便，但现金也很重要。} (Je trouve le paiement mobile pratique, mais l'espèce est aussi importante.)
\item \textbf{Conclusion (1 phrase) :} \zh{谢谢大家。}
\end{enumerate}
\textbf{Astuce :} Écris d'abord le plan en français, puis traduis phrase par phrase en chinois en réutilisant \textbf{exactement} les structures du cours (\zh{用的钱叫}, \zh{单位是}, \zh{发行}, \zh{用……付钱}). Vérifie les tons sur ton brouillon.
\end{methode}

\courssub{Production guidée : mini-exposé comparatif}

\begin{exercice}[Mini-exposé écrit : « Les deux monnaies »]
Rédige un court texte de 80 à 100 caractères chinois (environ 8-10 phrases courtes) pour présenter les différences et similitudes entre le renminbi et le franc CFA. Tu dois utiliser \textbf{au moins 10 mots} du vocabulaire de la leçon (surligne-les dans ta copie).
\textbf{Contraintes :}
\begin{itemize}
\item Utilise les structures : \zh{……用的钱叫……}, \zh{单位是……}, \zh{……发行……}, \zh{在……用……付钱}.
\item Inclus le vocabulaire : \zh{汇率} (taux de change), \zh{方便} (pratique), \zh{现金} (espèces), \zh{手机支付} (paiement mobile).
\item Soigne l'écriture des caractères (ordre des traits, proportion).
\end{itemize}
\end{exercice}

\begin{corrige}[Exercice mini-exposé : proposition de rédaction]
\textbf{Modèle (environ 95 caractères) :}
\zh{中国和喀麦隆的钱不一样。中国用的钱叫人民币，单位是元，也叫块。中国人民银行发行人民币。喀麦隆用的钱叫非洲法郎，单位是法郎。中非国家银行发行非洲法郎。今天汇率是一元等于八十法郎。在中国，人们常用手机支付，很方便。在喀麦隆，人们常用现金，也用 Mobile Money。我觉得两种钱都方便生活。}

\textbf{Traduction :} « Les monnaies de Chine et du Cameroun sont différentes. La Chine utilise le renminbi, unité le yuan (aussi appelé kuai). La Banque populaire de Chine émet le renminbi. Le Cameroun utilise le franc CFA, unité le franc. La BEAC émet le franc CFA. Le taux d'aujourd'hui est 1 yuan = 80 francs. En Chine, on utilise souvent le paiement mobile, c'est pratique. Au Cameroun, on utilise souvent les espèces, aussi le Mobile Money. Je trouve que les deux monnaies facilitent la vie. »

\textbf{Mots du vocabulaire utilisés (12) :} \zh{钱、人民币、元、块、中国人民银行、发行、非洲法郎、法郎、中非国家银行、汇率、手机支付、现金、方便、生活} (14 au total).
\end{corrige}

\begin{aretenir}[L'essentiel de la leçon 29 — Bilan]
\begin{itemize}
\item \textbf{Vocabulaire clé :} \zh{钱、人民币、元、块、角、分、非洲法郎、银行、中国人民银行、中非国家银行、发行、单位、现金、付钱、手机、换钱、汇率、支付、方便、习惯}.
\item \textbf{Structures grammaticales :}
      \begin{itemize}
      \item \zh{国家 + 用的钱叫 + 货币名} (Nommer la monnaie)
      \item \zh{单位是 + 单位名} (Donner l'unité)
      \item \zh{银行名 + 发行 + 货币名} (Dire qui émet)
      \item \zh{在 + 地方 + 用 + 方式 + 付钱} (Décrire le mode de paiement)
      \item \zh{今天汇率是 + 数字} (Annoncer le taux de change)
      \end{itemize}
\item \textbf{Faits socioculturels :}
      \begin{itemize}
      \item Chine : Renminbi (RMB/CNY), unité Yuan (¥), Banque populaire de Chine, paiement mobile dominant (WeChat/Alipay), billets à l'effigie de Mao Zedong.
      \item Cameroun : Franc CFA (FCFA/XAF), unité Franc, BEAC (siège Yaoundé), espèces + Mobile Money en développement, billets faune/agriculture/infrastructures CEMAC.
      \end{itemize}
\item \textbf{Point de vigilance :} Ne pas confondre \zh{人民币} (nom de la monnaie) et \zh{元/块} (unité de compte). Dire \zh{十块钱} ou \zh{十元}, jamais \zh{十人民币}.
\end{itemize}
\end{aretenir}

\coursec{Leçon 29 — Éléments socioculturels : différence entre la monnaie chinoise et la monnaie camerounaise}

\courssub{Découverte : deux monnaies, deux réalités}

\begin{textebox}[Dialogue : Au bureau de change de l'aéroport de Pékin]
\zh{— 你好，我想换钱。} \emph{Nǐ hǎo, wǒ xiǎng huàn qián.}\\
\zh{— 你好，换多少？} \emph{Nǐ hǎo, huàn duōshao?}\\
\zh{— 我有非洲法郎，想换人民币。} \emph{Wǒ yǒu Fēizhōu fǎláng, xiǎng huàn rénmínbì.}\\
\zh{— 好的，现在的汇率是 1 元 人民币 兑换 80 法郎 左右。} \emph{Hǎo de, xiànzài de huìlǜ shì 1 yuán rénmínbì duìhuàn 80 fǎláng zuǒyòu.}\\
\zh{— 那我换 50 000 法郎。} \emph{Nà wǒ huàn 50 000 fǎláng.}\\
\zh{— 没问题，请拿好你的 人民币 和 兑换单。} \emph{Méi wèntí, qǐng ná hǎo nǐ de rénmínbì hé duìhuàn dān.}\\
\zh{— 谢谢！} \emph{Xièxie!}
\end{textebox}

\begin{definition}[Mots nouveaux du dialogue]
\begin{itemize}
\item \zh{汇率} \emph{huìlǜ} (n.) : taux de change.
\item \zh{兑换} \emph{duìhuàn} (v.) : convertir, échanger (devises).
\item \zh{左右} \emph{zuǒyòu} (loc.) : environ, à peu près (placé après le nombre).
\item \zh{兑换单} \emph{duìhuàn dān} (n.) : bordereau de change, reçu de change.
\item \zh{拿好} \emph{ná hǎo} (v. + compl.) : bien prendre, bien garder.
\end{itemize}
\end{definition}

\begin{experience}[Activité de mise en route : « Mon porte-monnaie »]
\begin{enumerate}
\item En binôme, sortez une pièce ou un billet (CFA ou yuan si vous en avez) ou montrez une image sur téléphone.
\item Décrivez-le en chinois simple : \zh{这是...元/角/分} \emph{Zhè shì... yuán/jiǎo/fēn}, \zh{这是纸币/硬币} \emph{Zhè shì zhǐbì/yìngbì}.
\item Comparez : couleur, personnage/édifice représenté, taille, matière. Utilisez \zh{不一样} \emph{bù yīyàng} (différent) et \zh{一样} \emph{yīyàng} (pareil).
\item Mise en commun au tableau : le professeur note le vocabulaire clé qui émerge (\zh{颜色} couleur, \zh{图案} motif, \zh{大小} taille).
\end{enumerate}
\end{experience}

\courssub{Glossaire et écriture des caractères}

Dans la section précédente, nous avons découvert deux monnaies, deux réalités : le \cle{renminbi} chinois et le \cle{franc CFA} camerounais. Nous avons vu des billets, des pièces, des banques et des paiements par téléphone. Pour parler de tout cela en chinois, il faut d'abord maîtriser un petit vocabulaire de base, puis apprendre à écrire correctement les caractères clés. C'est l'objectif de cette section : lire, comprendre, prononcer et écrire les mots de la monnaie.

\courssub{Le glossaire de la monnaie}

\begin{definition}[Vocabulaire : l'argent et la monnaie]
Voici les quinze mots essentiels de la leçon. Observez bien la colonne « Nature » : en chinois comme en français, chaque mot a une fonction (nom, verbe, adjectif, spécificatif/classificateur).
\end{definition}

\begin{center}
\begin{tabularx}{\linewidth}{|X|X|X|X|}
\hline
\rowcolor{popblueL} Caractères & Pinyin & Nature & Traduction \\
\hline
钱 & qián & n. & argent, monnaie \\
\hline
人民币 & rénmínbì & n. & renminbi (monnaie chinoise) \\
\hline
元 & yuán & n./cl. & yuan (unité monétaire principale) \\
\hline
角 & jiǎo & cl. & jiao (1/10 de yuan) \\
\hline
分 & fēn & cl. & fen (1/100 de yuan) \\
\hline
非洲法郎 & Fēizhōu fǎláng & n. & franc CFA \\
\hline
银行 & yínháng & n. & banque \\
\hline
发行 & fāxíng & v. & émettre (monnaie, billet) \\
\hline
现金 & xiànjīn & n. & espèces, argent liquide \\
\hline
硬币 & yìngbì & n. & pièce de monnaie \\
\hline
纸币 & zhǐbì & n. & billet de banque \\
\hline
换钱 & huàn qián & v. + compl. & changer de l'argent \\
\hline
手机支付 & shǒujī zhīfù & n. & paiement mobile \\
\hline
用 & yòng & v. & utiliser \\
\hline
贵 & guì & adj. & cher \\
\hline
\end{tabularx}
\end{center}

\begin{exemplebox}[Les mots en situation]
\zh{我想换钱。} \emph{Wǒ xiǎng huàn qián.} « Je voudrais changer de l'argent. »\\
\zh{这个多少钱？} \emph{Zhège duōshao qián?} « Combien ça coûte ? »\\
\zh{人民币是中国的钱。} \emph{Rénmínbì shì Zhōngguó de qián.} « Le renminbi est la monnaie de la Chine. »\\
\zh{在银行可以换钱。} \emph{Zài yínháng kěyǐ huàn qián.} « On peut changer de l'argent à la banque. »\\
\zh{现在中国人常常用手机支付。} \emph{Xiànzài Zhōngguórén chángcháng yòng shǒujī zhīfù.} « Aujourd'hui, les Chinois utilisent souvent le paiement mobile. »\\
\zh{这个太贵了！} \emph{Zhège tài guì le!} « C'est trop cher ! »
\end{exemplebox}

\begin{attention}[角 \emph{jiǎo} ou 毛 \emph{máo} ? — Écrit vs Oral]
Le caractère \zh{角} se prononce \emph{jiǎo} dans la langue écrite et formelle : 1 \zh{角} = 10 \zh{分}, donc 10 \zh{角} = 1 \zh{元}. Mais à l'oral, dans les marchés et les magasins, les Chinois disent presque toujours \zh{毛} \emph{máo} à la place de \zh{角}. Ainsi « 0,5 yuan » se dit \zh{五毛} \emph{wǔ máo} et non \zh{五角} dans la conversation quotidienne. De même, \zh{元} \emph{yuán} devient souvent \zh{块} \emph{kuài} à l'oral : \zh{五块钱} \emph{wǔ kuài qián} = « cinq yuans ». \textbf{Règle :} Écrit = \zh{元}/\zh{角} ; Oral = \zh{块}/\zh{毛}.
\end{attention}

\courssub{Les clés (radicaux) : lire le sens caché des caractères}

\begin{propriete}[Les radicaux de la leçon]
Un \cle{radical} (clé) est un composant graphique qui donne souvent un indice de sens. Quatre radicaux sont utiles dans cette leçon :
\begin{itemize}
\item \zh{钅} (métal), forme abrégée de \zh{金} : on le trouve dans \zh{钱} (argent) et \zh{银} (argent métal, argent monétaire). Les anciennes monnaies étaient en métal !
\item \zh{巾} (tissu, morceau d'étoffe) : on le trouve dans \zh{币} (monnaie). Autrefois, des pièces d'étoffe servaient de moyen d'échange en Chine.
\item \zh{氵} (eau) : on le trouve dans \zh{法} (loi, droit), comme dans \zh{法郎} (franc).
\item \zh{阝} (monticule/ville) : cette clé apparaît dans des caractères comme \zh{院} (cour, institution) ou \zh{部} (département) ; elle aide à reconnaître les noms d'institutions (ex: \zh{银行} ne la contient pas mais \zh{部} oui).
\end{itemize}
\end{propriete}

\begin{popfigure}
\begin{tikzpicture}[grow=down, level distance=1.6cm, sibling distance=3.2cm,
  every node/.style={draw, rounded corners, fill=popblueL, align=center, font=\small},
  edge from parent/.style={draw=popblue, thick, -{Stealth}}]
\node[fill=poporangeL] {Radicaux de la leçon}
  child { node {\zh{钅} métal}
    child { node[align=center]{\zh{钱} qián\\argent} }
    child { node[align=center]{\zh{银} yín\\argent métal} }
  }
  child { node {\zh{巾} tissu}
    child { node[align=center]{\zh{币} bì\\monnaie} }
  }
  child { node {\zh{氵} eau}
    child { node[align=center]{\zh{法} fǎ\\loi, France} }
  };
\end{tikzpicture}
\legende{Mini-arbre des radicaux : \zh{钅} donne \zh{钱} et \zh{银} ; \zh{巾} donne \zh{币} ; \zh{氵} donne \zh{法}.}
\end{popfigure}

\begin{methode}[Mémoriser \zh{钱} grâce au radical]
\begin{enumerate}
\item Regarde la partie gauche du caractère : \zh{钅}, c'est le radical du \textbf{métal}.
\item Pense aux pièces de monnaie : elles sont en métal.
\item Fais le lien : métal + monnaie = « argent ». Le caractère \zh{钱} \emph{qián} signifie « argent ».
\item Applique la même stratégie à \zh{银} \emph{yín} : métal + \zh{艮} (indice phonétique) = « argent métal », d'où \zh{银行} « banque » (maison de l'argent).
\end{enumerate}
\end{methode}

\begin{savaistu}[\zh{银行}, la « maison de l'argent »]
Le mot \zh{银行} \emph{yínháng} « banque » se décompose littéralement : \zh{银} \emph{yín} = « argent (métal) » et \zh{行} \emph{háng} = « maison, établissement commercial, firme ». Une banque est donc, en chinois, la « maison de l'argent ». Attention : \zh{行} a une autre lecture, \emph{xíng}, qui signifie « aller, marcher » (comme dans \zh{不行} \emph{bù xíng} « ça ne va pas ») ou « rangée ». Le contexte décide de la prononciation !
\end{savaistu}

\courssub{L'ordre des traits : écrire correctement}

\begin{propriete}[Règles générales d'écriture]
En chinois, on écrit les traits dans un ordre fixe : de \textbf{haut en bas}, de \textbf{gauche à droite}, l'\textbf{horizontal avant le vertical}, l'extérieur avant l'intérieur, et on ferme un cadre en dernier. Respecter cet ordre rend l'écriture plus rapide, plus belle et plus facile à mémoriser.
\end{propriete}

\begin{exemplebox}[Quatre caractères clés, trait par trait]
\textbf{钱} \emph{qián} (10 traits) : d'abord le radical \zh{钅} à gauche (5 traits : \zh{丿}, \zh{一}, \zh{一}, \zh{一}, \zh{(trait brisé)}), puis la partie droite \zh{戋} (5 traits : \zh{一}, \zh{一}, \zh{㇂}, \zh{丿}, \zh{丶}).\\[2pt]
\textbf{银} \emph{yín} (11 traits) : le radical \zh{钅} (5 traits), puis la partie droite \zh{艮} (6 traits : \zh{一}, \zh{丨}, \zh{一}, \zh{一}, \zh{丿}, \zh{丶}).\\[2pt]
\textbf{币} \emph{bì} (4 traits) : \zh{丿} (le trait tombant en haut), puis \zh{丨}, puis \zh{(trait brisé)} (horizontal + crochet), puis le dernier \zh{丨} vertical au centre.\\[2pt]
\textbf{元} \emph{yuán} (4 traits) : d'abord \zh{二} (deux horizontaux, le court en haut, le long en dessous), puis \zh{儿} (\zh{丿} à gauche, puis le crochet \zh{乚} à droite).
\end{exemplebox}

\begin{popfigure}
\begin{tikzpicture}[scale=1.05, font=\small]
% Grille 钱
\draw[popline, dashed] (0,0) grid (2,2);
\draw[popline] (0,0) rectangle (2,2);
\node[popdark] at (1,1) {\Huge 钱};
\node[below, align=center] at (1,-0.3) {\zh{钱} qián : 10 traits\\gauche \zh{钅} (5) puis droite \zh{戋} (5)};
% Grille 币
\begin{scope}[xshift=4cm]
\draw[popline, dashed] (0,0) grid (2,2);
\draw[popline] (0,0) rectangle (2,2);
\draw[poporange, very thick, -{Stealth}] (1.3,1.9) .. controls (1.0,1.7) .. (0.7,1.45);
\node[poporange] at (0.5,1.85) {1};
\draw[poporange, very thick, -{Stealth}] (0.7,1.35) -- (0.7,0.25);
\node[poporange] at (0.45,1.15) {2};
\draw[poporange, very thick, -{Stealth}] (0.85,1.35) -- (1.45,1.35) -- (1.45,0.25);
\node[poporange] at (1.62,1.15) {3};
\draw[poporange, very thick, -{Stealth}] (1.05,1.25) -- (1.05,0.15);
\node[poporange] at (1.28,0.6) {4};
\node[below, align=center] at (1,-0.3) {\zh{币} bì : 4 traits\\丿, 丨, (trait brisé), 丨};
\end{scope}
\end{tikzpicture}
\legende{Écriture de \zh{钱} (composant par composant) et ordre numéroté des 4 traits de \zh{币}.}
\end{popfigure}

\begin{attention}[Caractères proches : ne pas confondre !]
\begin{itemize}
\item \zh{钱} \emph{qián} (argent) et \zh{线} \emph{xiàn} (fil, ligne) : même partie droite \zh{戋}, mais le radical change tout : \zh{钅} métal $\rightarrow$ argent ; \zh{纟} soie $\rightarrow$ fil.
\item \zh{元} \emph{yuán} (yuan) et \zh{无} \emph{wú} (ne pas avoir) : très ressemblants ! La différence : le dernier trait de \zh{元} est un crochet simple \zh{乚} (fermé), celui de \zh{无} est un crochet relevé \zh{乚} plus ouvert, et \zh{无} n'a pas le petit \zh{丿} à gauche en bas (pas de composant \zh{儿} complet). Écrivez-les côte à côte pour bien voir.
\end{itemize}
\end{attention}

\begin{exoresolu}[Reconnaître et compléter]
Énoncé : (1) Quel radical commun trouve-t-on dans \zh{钱} et \zh{银} ? Que signifie-t-il ? (2) Complète avec le bon caractère (\zh{元} ou \zh{角}) : 这本书十 \_\_\_ 五 \_\_\_ 。 (3) Traduis : \zh{我想在银行换钱。}
\tcblower
Solution détaillée : (1) Le radical commun est \zh{钅}, la clé du \textbf{métal} ; il rappelle que les monnaies anciennes étaient métalliques. (2) \zh{这本书十元五角。} \emph{Zhè běn shū shí yuán wǔ jiǎo.} « Ce livre coûte dix yuans cinq jiao (10,5 yuan). » On met l'unité principale \zh{元} avant l'unité secondaire \zh{角}, comme en français « dix euros cinquante ». (3) \zh{我想在银行换钱。} \emph{Wǒ xiǎng zài yínháng huàn qián.} « Je voudrais changer de l'argent à la banque. » Structure : Sujet + 想 + 在 + lieu + verbe.
\end{exoresolu}

\courssub{Copie guidée des caractères}

\begin{experience}[Atelier d'écriture sur ardoise ou cahier]
\begin{enumerate}
\item Le professeur écrit au tableau \zh{钱}, \zh{银}, \zh{币}, \zh{元} en grand, trait par trait, en annonçant chaque trait à voix haute.
\item Les élèves tracent chaque caractère dans les airs avec le doigt, en respectant l'ordre des traits.
\item Sur l'ardoise ou le cahier, copie chaque caractère \textbf{5 fois} : \zh{钱} $\times$5, \zh{银} $\times$5, \zh{币} $\times$5, \zh{元} $\times$5.
\item Écris ensuite sans modèle les deux mots : \zh{人民币} (renminbi) et \zh{银行} (banque).
\item Vérification en binôme : ton voisin contrôle l'ordre des traits et la place des radicaux.
\end{enumerate}
\end{experience}

\begin{aretenir}[L'essentiel de la section]
\begin{itemize}
\item Le vocabulaire de la monnaie : \zh{钱} argent, \zh{人民币} renminbi, \zh{元}/\zh{角}/\zh{分} les trois unités, \zh{非洲法郎} franc CFA, \zh{银行} banque, \zh{现金} liquide, \zh{硬币} pièce, \zh{纸币} billet, \zh{换钱} changer de l'argent, \zh{手机支付} paiement mobile.
\item Les radicaux donnent le sens : \zh{钅} métal (\zh{钱}, \zh{银}), \zh{巾} tissu (\zh{币}), \zh{氵} eau (\zh{法}).
\item L'ordre des traits : haut $\rightarrow$ bas, gauche $\rightarrow$ droite ; \zh{钱} 10 traits, \zh{银} 11 traits, \zh{币} 4 traits, \zh{元} 4 traits.
\item À l'oral : \zh{块} remplace \zh{元} et \zh{毛} remplace \zh{角}.
\end{itemize}
\end{aretenir}

\courssub{Compréhension du texte : La monnaie en Chine et au Cameroun}

\begin{textebox}[Texte : Évolution et usage des monnaies]
\zh{中国的钱叫人民币，单位是元、角、分。中国人民银行发行纸币和硬币。现在，中国人很少用现金，大多数人用手机支付，比如微信支付和支付宝。扫一扫二维码，就能付钱，非常方便。}\\
\zh{喀麦隆的钱叫非洲法郎，单位是法郎。中央银行发行纸币和硬币。在喀麦隆，人们常用现金，但手机支付也越来越多，比如 Orange Money 和 MTN MoMo. 两个国家的货币都有自己的历史和特点。}
\end{textebox}

\begin{exemplebox}[Traduction et notes]
\emph{Zhōngguó de qián jiào Rénmínbì, dānwèi shì yuán, jiǎo, fēn. Zhōngguó Rénmín Yínháng fāxíng zhǐbì hé yìngbì. Xiànzài, Zhōngguórén hěnshǎo yòng xiànjīn, dàduōshù rén yòng shǒujī zhīfù, bǐrú Wēixìn zhīfù hé Zhīfùbǎo. Sǎo yī sǎo èrwéimǎ, jiù néng fù qián, fēicháng fāngbiàn.}\\
« La monnaie de Chine s'appelle le Renminbi, les unités sont yuan, jiao, fen. La Banque populaire de Chine émet des billets et des pièces. Aujourd'hui, les Chinois utilisent peu l'espèce, la plupart utilisent le paiement mobile, comme WeChat Pay et Alipay. On scanne le QR code et on paie, c'est très pratique. »\\[4pt]
\emph{Kāmàilóng de qián jiào Fēizhōu fǎláng, dānwèi shì fǎláng. Zhōngyāng Yínháng fāxíng zhǐbì hé yìngbì. Zài Kāmàilóng, rénmen cháng yòng xiànjīn, dàn shǒujī zhīfù yě yuè lái yuè duō, bǐrú Orange Money hé MTN MoMo. Liǎng gè guójiā de huòbì dōu yǒu zìjǐ de lìshǐ hé tèdiǎn.}\\
« La monnaie du Cameroun s'appelle le Franc CFA, l'unité est le franc. La Banque centrale émet des billets et des pièces. Au Cameroun, les gens utilisent souvent l'espèce, mais le paiement mobile augmente aussi, comme Orange Money et MTN MoMo. Les deux pays ont leur propre histoire et leurs caractéristiques. »
\end{exemplebox}

\begin{exercice}[Questions de compréhension]
\begin{enumerate}
\item Quelle est l'unité principale du Renminbi ? Et celle du Franc CFA ?
\item Quel est le nom de la banque centrale chinoise ? Et celui de la banque centrale des États d'Afrique centrale (BEAC) ?
\item Citez deux applications de paiement mobile en Chine et deux au Cameroun.
\item Selon le texte, quel est le moyen de paiement dominant en Chine aujourd'hui ? Et au Cameroun ?
\item Traduisez en français : \zh{扫一扫二维码，就能付钱，非常方便。}
\end{enumerate}
\end{exercice}

\begin{corrige}[Exercice Compréhension]
\begin{enumerate}
\item Chine : \zh{元} (yuán) ; Cameroun : \zh{法郎} (fǎláng).
\item Chine : \zh{中国人民银行} (Banque populaire de Chine) ; Cameroun/BEAC : \zh{中央银行} (Banque centrale / BEAC).
\item Chine : \zh{微信支付} (WeChat Pay), \zh{支付宝} (Alipay) ; Cameroun : Orange Money, MTN MoMo.
\item Chine : \zh{手机支付} (paiement mobile) dominant, \zh{很少用现金} (peu d'espèces) ; Cameroun : \zh{常用现金} (espèces courantes), mais mobile \zh{越来越多} (en hausse).
\item « On scanne le QR code et on peut payer, c'est très pratique. »
\end{enumerate}
\end{corrige}

\courssub{Faits comparés Cameroun / Chine}

\begin{propriete}[Tableau comparatif : Renminbi vs Franc CFA]
\begin{center}
\begin{tabularx}{\linewidth}{|X|X|X|}
\hline
\rowcolor{popblueL} \textbf{Critère} & \textbf{Chine (Renminbi / 人民币)} & \textbf{Cameroun (Franc CFA / 非洲法郎)} \\
\hline
Unité principale & \zh{元} (yuán) & \zh{法郎} (fǎláng) \\
\hline
Sous-unités & 1 元 = 10 角 = 100 分 & 1 Franc = 100 centimes (théorique, pièces rares) \\
\hline
Émission & \zh{中国人民银行} (Banque populaire de Chine) & \zh{中央银行} (BEAC -- Banque des États de l'Afrique Centrale) \\
\hline
Billets courants & 1, 5, 10, 20, 50, 100 元 (5e série 2019) & 500, 1000, 2000, 5000, 10000 FCFA \\
\hline
Pièces courantes & 1 元, 5 角, 1 角 (1 分 rare) & 1, 2, 5, 10, 25, 50, 100, 500 FCFA \\
\hline
Parité / Change & Taux flottant géré (≈ 1 € = 7,8 CNY) & Parité fixe avec l'Euro : 1 € = 655,957 FCFA \\
\hline
Paiement mobile & \textbf{Très dominant} (WeChat Pay, Alipay) : QR code, NFC & \textbf{En forte croissance} (Orange Money, MTN MoMo) : USSD, App \\
\hline
Usage espèces & \zh{很少} (très rare en ville) & \zh{很常见} (très courant, dominant) \\
\hline
Symbole culturel & \zh{红包} (hóngbāo) : enveloppes rouges Nouvel An & Enveloppes blanches/colorées : fêtes, mariages, dot \\
\hline
\end{tabularx}
\end{center}
\end{propriete}

\begin{savaistu}[Le Franc CFA : particularités]
Le \zh{非洲法郎} (Franc de la Coopération Financière en Afrique) est utilisé par 6 pays de la zone CEMAC (Cameroun, Tchad, RCA, Congo, Gabon, Guinée équatoriale). Sa parité fixe avec l'Euro (anciennement Franc français) garantit la stabilité mais limite la politique monétaire nationale. Les billets portent la mention d'un pays (ex: \zh{C} pour Cameroun, \zh{E} pour Tchad) mais ont cours légal dans toute la zone.
\end{savaistu}

\begin{savaistu}[Le Renminbi : internationalisation]
Le \zh{人民币} (RMB) est la monnaie officielle de la RPC depuis 1948. Depuis 2016, il fait partie du panier des Droits de Tirage Spéciaux (DTS) du FMI aux côtés du Dollar, Euro, Yen, Livre sterling. La Chine promeut son usage dans le commerce international (ex: paiement du pétrole, accords de swap avec la BEAC).
\end{savaistu}

\courssub{Exprimer son point de vue et débattre}

\begin{propriete}[Structures pour comparer et donner son avis]
\begin{center}
\begin{tabularx}{\linewidth}{|X|X|X|}
\hline
\rowcolor{popblueL} \textbf{Structure} & \textbf{Exemple (Caractères + Pinyin)} & \textbf{Traduction} \\
\hline
A \zh{比} B \zh{形容词} & \zh{中国的手机支付比喀麦隆方便。} \emph{Zhōngguó de shǒujī zhīfù bǐ Kāmàilóng fāngbiàn.} & Le paiement mobile en Chine est plus pratique qu'au Cameroun. \\
\hline
A \zh{跟/和} B \zh{不一样} & \zh{人民币跟非洲法郎不一样。} \emph{Rénmínbì gēn Fēizhōu fǎláng bù yīyàng.} & Le Renminbi et le Franc CFA sont différents. \\
\hline
\zh{我觉得} / \zh{我认为} + phrase & \zh{我觉得无现金社会很好。} \emph{Wǒ juéde wú xiànjīn shèhuì hěn hǎo.} & Je pense que la société sans espèces est bien. \\
\hline
\zh{虽然}... \zh{但是}... & \zh{虽然现金少了，但是安全多了。} \emph{Suīrán xiànjīn shǎo le, dànshì ānquán duō le.} & Bien qu'il y ait moins d'espèces, c'est beaucoup plus sûr. \\
\hline
\zh{优点} / \zh{缺点} & \zh{手机支付的优点是快，缺点是要网。} \emph{Shǒujī zhīfù de yōudiǎn shì kuài, quēdiǎn shì yào wǎng.} & L'avantage du paiement mobile est la rapidité, l'inconvénient est qu'il faut Internet. \\
\hline
\end{tabularx}
\end{center}
\end{propriete}

\begin{experience}[Débat guidé : « Vers une société sans espèces ? »]
\begin{enumerate}
\item \textbf{Préparation (5 min) :} En groupes de 4. Chaque groupe tire un rôle : « Partisans du 100\% mobile (Chine) », « Défenseurs de l'espèce (Cameroun actuel) », « Observateurs neutres », « Banquiers ».
\item \textbf{Arguments clés à préparer en chinois (aide au tableau) :}
\begin{itemize}
\item \zh{方便} (pratique), \zh{快} (rapide), \zh{安全} (sûr), \zh{环保} (écolo -- pas de papier).
\item \zh{不用网} (pas besoin de réseau), \zh{保护隐私} (protège la vie privée), \zh{老人用得会} (personnes âgées savent utiliser).
\end{itemize}
\item \textbf{Débat (10 min) :} Tour de table. Chaque élève doit prononcer au moins 2 phrases complètes avec \zh{比}, \zh{虽然...但是...} ou \zh{我觉得}.
\item \textbf{Synthèse (5 min) :} Le professeur note au tableau les arguments gagnants en caractères + pinyin.
\end{enumerate}
\end{experience}

\courssub{Production guidée : Mini-exposé comparatif}

\begin{methode}[Rédiger un mini-exposé comparatif (80--100 caractères)]
\begin{enumerate}
\item \textbf{Introduction (1 phrase) :} Présenter les deux monnaies. \zh{中国用人民币，喀麦隆用非洲法郎。}
\item \textbf{Point commun (1 phrase) :} \zh{两个国家都有纸币和硬币。} / \zh{都在推广手机支付。}
\item \textbf{Différence majeure (1--2 phrases) :} Utiliser \zh{比} ou \zh{不一样}. \zh{中国人很少用现金，喀麦隆人常用现金。} \zh{中国的手机支付比喀麦隆普及。}
\item \textbf{Avis personnel (1 phrase) :} \zh{我觉得手机支付更方便，但现金也很重要。}
\item \textbf{Conclusion (1 phrase) :} \zh{了解两种货币，对交流很有帮助。}
\end{enumerate}
\end{methode}

\begin{exemplebox}[Modèle d'exposé écrit]
\zh{中国用人民币，喀麦隆用非洲法郎。两个国家都有纸币和硬币，也都在推广手机支付。但是，中国的手机支付比喀麦隆普及，中国人很少用现金，喀麦隆人常用现金。我觉得手机支付更方便、更快，但是现金不需要网络，对老人更友好。了解两种货币的差别，对中非交流很有帮助。}\\[4pt]
\emph{Zhōngguó yòng Rénmínbì, Kāmàilóng yòng Fēizhōu fǎláng. Liǎng gè guójiā dōu yǒu zhǐbì hé yìngbì, yě dōu zài tuīguǎng shǒujī zhīfù. Dànshì, Zhōngguó de shǒujī zhīfù bǐ Kāmàilóng pǔjí, Zhōngguórén hěnshǎo yòng xiànjīn, Kāmàilóng rén cháng yòng xiànjīn. Wǒ juéde shǒujī zhīfù gèng fāngbiàn, gèng kuài, dànshì xiànjīn bù xūyào wǎngluò, duì lǎorén gèng yǒuhǎo. Liǎojiě liǎng zhǒng huòbì de chābié, duì Zhōng-Fēi jiāoliú hěn yǒu bāngzhù.}
\end{exemplebox}

\begin{exercice}[Production écrite : Votre mini-exposé]
Rédigez votre propre exposé comparatif (80--100 caractères chinois) en suivant la méthode ci-dessus. Vous devez utiliser \textbf{au moins} : \zh{比}, \zh{虽然...但是...}, \zh{觉得/认为}, \zh{优点/缺点} (ou \zh{方便/不方便}). Écrivez les caractères, le pinyin et la traduction française.
\end{exercice}

\begin{corrige}[Exercice Production écrite -- Proposition de correction]
\zh{中国和喀麦隆的货币不一样。中国用人民币，单位是元；喀麦隆用非洲法郎，单位是法郎。虽然两个国家都有手机支付，但是中国比喀麦隆普及得多。在中国，付钱很方便，只要扫码；在喀麦隆，现金还是很重要。我觉得手机支付的优点是快、安全，缺点是没网不能用。我希望喀麦隆的手机支付以后越来越方便。}\\[4pt]
\emph{Zhōngguó hé Kāmàilóng de huòbì bù yīyàng. Zhōngguó yòng Rénmínbì, dānwèi shì yuán; Kāmàilóng yòng Fēizhōu fǎláng, dānwèi shì fǎláng. Suīrán liǎng gè guójiā dōu yǒu shǒujī zhīfù, dànshì Zhōngguó bǐ Kāmàilóng pǔjí de duō. Zài Zhōngguó, fù qián hěn fāngbiàn, zhǐyào sǎomǎ; zài Kāmàilóng, xiànjīn hái shì hěn zhòngyào. Wǒ juéde shǒujī zhīfù de yōudiǎn shì kuài, ānquán, quēdiǎn shì méi wǎng bù néng yòng. Wǒ xīwàng Kāmàilóng de shǒujī zhīfù yǐhòu yuè lái yuè fāngbiàn.}
\end{corrige}

\begin{aretenir}[Bilan de la leçon 29]
\begin{itemize}
\item \textbf{Vocabulaire clé :} \zh{人民币, 非洲法郎, 元/角/分, 银行, 发行, 现金, 硬币, 纸币, 换钱, 手机支付, 汇率, 兑换, 红包}.
\item \textbf{Grammaire :} Comparatif \zh{A 比 B Adj}, Différence \zh{A 跟 B 不一样}, Concession \zh{虽然...但是...}, Opinion \zh{我觉得/我认为}.
\item \textbf{Socioculturel :} Chine = société quasi sans espèces (QR code omniprésent), enveloppes rouges \zh{红包}. Cameroun = espèces dominantes, mobile money (USSD/App) en essor, parité fixe Euro, enveloppes blanches/cérémonies.
\item \textbf{Écriture :} Radicaux \zh{钅} (métal), \zh{巾} (tissu), \zh{氵} (eau). Caractères \zh{钱, 银, 币, 元, 角, 分}.
\item \textbf{Oral :} \zh{块} (kuài) = \zh{元}, \zh{毛} (máo) = \zh{角}.
\end{itemize}
\end{aretenir}

\coursec{Compréhension du texte}

Dans la section précédente, nous avons découvert le glossaire du thème (\zh{人民币}, \zh{非洲法郎}, \zh{银行}, \zh{现金}, \zh{手机支付}…) et nous avons appris à écrire les caractères clés en respectant l'ordre des traits. Nous sommes maintenant prêts à revenir sur le texte support de la leçon et à le comprendre en profondeur : phrase par phrase, puis globalement, puis en détail. Objectif : être capable de répondre, en chinois, à la question centrale de la leçon — \zh{中国人用什么钱？喀麦隆人用什么钱？}

\courssub{Relecture guidée du texte, phrase par phrase}

\begin{methode}[Comment relire un texte chinois en quatre passages]
\begin{enumerate}
\item \textbf{Lecture silencieuse} : repérer les mots déjà connus grâce au glossaire.
\item \textbf{Lecture à voix haute} : le professeur lit, les élèves répètent en veillant aux tons.
\item \textbf{Traduction collective} : chaque phrase est traduite à tour de rôle ; on s'aide des mots connus avant de regarder le glossaire.
\item \textbf{Vérification} : on compare sa traduction avec la version française de référence.
\end{enumerate}
\end{methode}

Reprenons le texte support de la leçon et découpons-le en quatre phrases de travail. Pour chaque phrase, procédons ainsi : lecture, observation des mots connus, traduction.

\textbf{Phrase 1.} \zh{在中国，人们用人民币。} \emph{Zài Zhōngguó, rénmen yòng rénmínbì.} « En Chine, on utilise le renminbi. »

Observation : la structure est \cle{在 + lieu + ，sujet + 用 + monnaie}. Le mot \zh{用} (yòng, « utiliser ») est le verbe central de toute la leçon.

\textbf{Phrase 2.} \zh{人民币是中国人民银行发行的。} \emph{Rénmínbì shì Zhōngguó Rénmín Yínháng fāxíng de.} « Le renminbi est émis par la Banque populaire de Chine. »

Observation : la structure \cle{是 …… 发行的} (shì … fāxíng de) exprime le passif : « est émis par ». C'est une tournure très fréquente pour présenter une information factuelle.

\textbf{Phrase 3.} \zh{在喀麦隆，人们用非洲法郎，就是中非法郎。} \emph{Zài Kāmàilóng, rénmen yòng Fēizhōu fǎláng, jiùshì Zhōngfēi fǎláng.} « Au Cameroun, on utilise le franc CFA, c'est-à-dire le franc de la BEAC. »

Observation : \zh{就是} (jiùshì) introduit une précision, comme « c'est-à-dire » en français.

\textbf{Phrase 4.} \zh{在中国，很多人用手机付钱；在喀麦隆，很多人还是用现金。} \emph{Zài Zhōngguó, hěn duō rén yòng shǒujī fù qián; zài Kāmàilóng, hěn duō rén háishì yòng xiànjīn.} « En Chine, beaucoup de gens paient avec leur téléphone portable ; au Cameroun, beaucoup de gens paient encore en espèces. »

Observation : \zh{还是} (háishì) signifie ici « encore, toujours » (adverbe), et non « ou bien » comme dans les questions alternatives. Le point-virgule chinois \zh{；} relie deux faits comparés.

\begin{exemplebox}[Exemples traduits pour consolider]
\zh{在中国，很多人用手机付钱。} \emph{Zài Zhōngguó, hěn duō rén yòng shǒujī fù qián.} « En Chine, beaucoup de gens paient avec leur téléphone portable. »\\[2pt]
\zh{在杜阿拉的市场，人们常用现金。} \emph{Zài Dù'ālā de shìchǎng, rénmen cháng yòng xiànjīn.} « Au marché de Douala, on utilise souvent les espèces. »\\[2pt]
\zh{非洲法郎是中非国家银行发行的。} \emph{Fēizhōu fǎláng shì Zhōngfēi guójiā yínháng fāxíng de.} « Le franc CFA est émis par la banque centrale des États d'Afrique centrale. »
\end{exemplebox}

\courssub{Questions de compréhension globale}

Après la relecture, on vérifie d'abord la compréhension \textbf{globale} : de quoi parle le texte ? Deux questions suffisent.

\begin{enumerate}
\item \textbf{Quelle monnaie utilise-t-on en Chine ?} — On utilise le renminbi (\zh{人民币}), dont l'unité de base est le yuan (\zh{元}).
\item \textbf{Quelle monnaie utilise-t-on au Cameroun ?} — On utilise le franc CFA (\zh{非洲法郎}), émis dans la région par la banque centrale commune (\zh{中非国家银行}).
\end{enumerate}

\begin{methode}[Répondre à une question en chinois : reprendre les mots de la question]
En chinois, la réponse courte reprend les mots de la question, comme en français on répond « oui, j'utilise… ». Observez la transformation :
\begin{itemize}
\item Question : \zh{中国人用什么钱？} \emph{Zhōngguórén yòng shénme qián ?} « Quelle monnaie les Chinois utilisent-ils ? »
\item Réponse : \zh{他们用人民币。} \emph{Tāmen yòng rénmínbì.} « Ils utilisent le renminbi. »
\end{itemize}
Le mot interrogatif \zh{什么} (shénme, « quoi ») est simplement \textbf{remplacé} par la réponse, à la même place dans la phrase. L'ordre des mots ne change pas : Sujet + 用 + monnaie. Même méthode pour le Cameroun :
\begin{itemize}
\item \zh{喀麦隆人用什么钱？} \emph{Kāmàilóngrén yòng shénme qián ?} « Quelle monnaie les Camerounais utilisent-ils ? »
\item \zh{他们用非洲法郎。} \emph{Tāmen yòng Fēizhōu fǎláng.} « Ils utilisent le franc CFA. »
\end{itemize}
\end{methode}

\courssub{Questions de détail}

Passons maintenant aux questions de \textbf{détail}, qui demandent de retrouver une information précise dans le texte.

\begin{textebox}[Questions types sur le texte]
1. \zh{人民币是哪个银行发行的？} \emph{Rénmínbì shì nǎge yínháng fāxíng de ?} « Par quelle banque le renminbi est-il émis ? »\\
2. \zh{喀麦隆人常用现金还是手机支付？} \emph{Kāmàilóngrén cháng yòng xiànjīn háishì shǒujī zhīfù ?} « Les Camerounais utilisent-ils souvent les espèces ou le paiement mobile ? »\\
3. \zh{你觉得手机支付方便吗？} \emph{Nǐ juéde shǒujī zhīfù fāngbiàn ma ?} « Trouves-tu que le paiement mobile est pratique ? »
\end{textebox}

\begin{corrige}[Réponses attendues]
\begin{enumerate}
\item \zh{人民币是中国人民银行发行的。} « Le renminbi est émis par la Banque populaire de Chine. »
\item \zh{他们常用现金。} « Ils utilisent souvent les espèces. » (Question alternative avec \zh{还是} : on choisit l'un des deux termes, on ne répond pas « oui ».)
\item Réponse libre, modèle : \zh{我觉得手机支付很方便。} \emph{Wǒ juéde shǒujī zhīfù hěn fāngbiàn.} « Je trouve que le paiement mobile est très pratique. »
\end{enumerate}
\end{corrige}

Récapitulons les informations de détail dans un tableau :

\begin{center}
\begin{tabularx}{\linewidth}{|X|X|X|}
\hline
\rowcolor{popblueL} Question & Chine & Cameroun \\
\hline
Quelle monnaie ? & 人民币 (rénmínbì), le renminbi & 非洲法郎 (Fēizhōu fǎláng), le franc CFA \\
\hline
Qui émet ? & 中国人民银行, la Banque populaire de Chine & 中非国家银行 (BEAC), la banque centrale de la région \\
\hline
Comment paie-t-on souvent ? & 手机支付 (shǒujī zhīfù), le paiement mobile & 现金 (xiànjīn), les espèces (et de plus en plus le mobile money) \\
\hline
\end{tabularx}
\end{center}

\begin{attention}[付钱, littéralement « verser l'argent »]
\zh{付钱} (fù qián) signifie « payer » : c'est un verbe + son complément figé, littéralement « verser l'argent ». Deux pièges pour les francophones :
\begin{itemize}
\item Ne dites pas *\zh{给钱了} pour « payer » dans ce contexte : \zh{给钱} (gěi qián) signifie « donner de l'argent » (à quelqu'un, comme un cadeau ou une aide), pas « régler un achat ».
\item N'insérez pas \zh{了} entre le verbe et son complément dans une habitude : on dit \zh{很多人用手机付钱} et non *\zh{付了钱} pour décrire une pratique générale. \zh{了} marque une action ponctuelle accomplie : \zh{我昨天付了钱。} \emph{Wǒ zuótiān fù le qián.} « J'ai payé hier. »
\end{itemize}
\end{attention}

\courssub{Organigramme de compréhension}

La figure suivante résume le chemin de compréhension du texte : trois questions clés permettent d'en extraire toutes les informations.

\begin{popfigure}
\begin{tikzpicture}[
  node distance=0.9cm,
  box/.style={draw=popblue, fill=popblueL, rounded corners, align=center, minimum width=3.2cm, minimum height=0.9cm, font=\small},
  qbox/.style={draw=poporange, fill=poporangeL, rounded corners, align=center, minimum width=3.6cm, minimum height=0.9cm, font=\small},
  abox/.style={draw=popgreen, fill=popgreenL, rounded corners, align=center, minimum width=3.6cm, minimum height=0.9cm, font=\small},
  arr/.style={-{Stealth[length=2.5mm]}, thick, popdark}
]
\node[box, align=center] (t){Texte : 两种钱\\ \emph{Deux monnaies}};
\node[qbox, below left=1.1cm and 2.6cm of t, align=center] (q1){Qui ?\\ 谁发行的？};
\node[qbox, below=1.1cm of t, align=center] (q2){Quoi ?\\ 用什么钱？};
\node[qbox, below right=1.1cm and 2.6cm of t, align=center] (q3){Comment ?\\ 怎么付钱？};
\node[abox, below=0.9cm of q1, align=center] (a1){中国人民银行\\ 中非国家银行};
\node[abox, below=0.9cm of q2, align=center] (a2){人民币\\ 非洲法郎};
\node[abox, below=0.9cm of q3, align=center] (a3){手机支付\\ 现金};
\draw[arr] (t) -- (q1);
\draw[arr] (t) -- (q2);
\draw[arr] (t) -- (q3);
\draw[arr] (q1) -- (a1);
\draw[arr] (q2) -- (a2);
\draw[arr] (q3) -- (a3);
\end{tikzpicture}
\legende{Organigramme de compréhension : du texte aux trois questions clés (émetteurs, monnaies, moyens de paiement).}
\end{popfigure}

\courssub{Vérification : vrai ou faux}

\begin{exoresolu}[Vrai ou faux ?]
Énoncé. Lisez les affirmations suivantes et dites si elles sont vraies ou fausses, en justifiant par une phrase du texte.
\begin{enumerate}
\item « Le franc CFA est émis par la Banque populaire de Chine. »
\item « En Chine, beaucoup de gens paient avec leur téléphone portable. »
\item « Au Cameroun, personne n'utilise les espèces. »
\end{enumerate}
\tcblower
Solution détaillée.
\begin{enumerate}
\item \textbf{Faux.} Le texte dit : \zh{人民币是中国人民银行发行的。} C'est le \emph{renminbi} qui est émis par la Banque populaire de Chine. Le franc CFA est émis par la banque centrale de la région (la BEAC, \zh{中非国家银行}) pour le Cameroun. Attention à ne pas croiser les deux colonnes du tableau !
\item \textbf{Vrai.} Le texte dit : \zh{在中国，很多人用手机付钱。} « En Chine, beaucoup de gens paient avec leur téléphone portable. »
\item \textbf{Faux.} Le texte dit : \zh{在喀麦隆，很多人还是用现金。} « Au Cameroun, beaucoup de gens paient encore en espèces. » Le mot \zh{很多} (beaucoup) contredit directement l'affirmation « personne ».
\end{enumerate}
\end{exoresolu}

\begin{exercice}[À vous de jouer]
Répondez en chinois, en phrases complètes, en reprenant les mots de la question :
\begin{enumerate}
\item \zh{中国人用什么钱？}
\item \zh{喀麦隆人用什么钱？}
\item \zh{人民币是哪个银行发行的？}
\item \zh{你常用现金还是手机支付？}
\end{enumerate}
\end{exercice}

\begin{corrige}[Exercice « À vous de jouer »]
\begin{enumerate}
\item \zh{他们用人民币。} \emph{Tāmen yòng rénmínbì.} « Ils utilisent le renminbi. »
\item \zh{他们用非洲法郎。} \emph{Tāmen yòng Fēizhōu fǎláng.} « Ils utilisent le franc CFA. »
\item \zh{人民币是中国人民银行发行的。} \emph{Rénmínbì shì Zhōngguó Rénmín Yínháng fāxíng de.} « Le renminbi est émis par la Banque populaire de Chine. »
\item Réponse personnelle, modèles : \zh{我常用现金。} \emph{Wǒ cháng yòng xiànjīn.} « J'utilise souvent les espèces. » / \zh{我常用手机支付。} \emph{Wǒ cháng yòng shǒujī zhīfù.} « J'utilise souvent le paiement mobile. »
\end{enumerate}
\end{corrige}

\courssub{Reformulation : résumer le texte en deux phrases}

Dernière étape de la compréhension : chaque élève doit être capable de \textbf{résumer le texte} en deux phrases simples en chinois, sans regarder le texte. La méthode consiste à garder les deux structures maîtresses de la leçon : \cle{在 + lieu + ，人们用 …} et \cle{…… 是 …… 发行的}.

\begin{exemplebox}[Modèles de résumé en deux phrases]
Modèle A (monnaies) :\\
\zh{在中国，人们用人民币；在喀麦隆，人们用非洲法郎。} \emph{Zài Zhōngguó, rénmen yòng rénmínbì; zài Kāmàilóng, rénmen yòng Fēizhōu fǎláng.} « En Chine, on utilise le renminbi ; au Cameroun, on utilise le franc CFA. »\\[4pt]
Modèle B (paiement) :\\
\zh{中国人常用手机付钱，喀麦隆人常用现金。} \emph{Zhōngguórén cháng yòng shǒujī fù qián, Kāmàilóngrén cháng yòng xiànjīn.} « Les Chinois paient souvent par téléphone portable, les Camerounais paient souvent en espèces. »
\end{exemplebox}

\begin{experience}[Résumé à voix haute]
Par groupes de deux : l'élève A résume le texte en deux phrases chinoises (modèle A ou B, ou les deux), l'élève B vérifie avec l'organigramme de compréhension (Qui ? Quoi ? Comment ?) et signale toute information manquante ou fausse. Puis on échange les rôles. Le professeur fait ensuite produire deux ou trois résumés au tableau et la classe les corrige collectivement.
\end{experience}

\begin{aretenir}[L'essentiel de la section]
\begin{itemize}
\item Pour répondre à une question chinoise, on \textbf{reprend les mots de la question} et on remplace le mot interrogatif (\zh{什么}, \zh{哪个}, \zh{谁}) par la réponse, sans changer l'ordre des mots.
\item \zh{是 …… 发行的} exprime le passif : « est émis par ».
\item \zh{付钱} (fù qián) = « payer » ; ne pas confondre avec \zh{给钱} (gěi qián) = « donner de l'argent ».
\item \zh{还是} : « ou bien » dans une question alternative ; « encore, toujours » comme adverbe dans une phrase affirmative.
\end{itemize}
\end{aretenir}

\coursec{Faits comparés Cameroun / Chine}

Après avoir lu et compris le texte de la section précédente, nous savons maintenant \emph{ce que disent} les deux monnaies. Il reste à organiser ces informations, à les comparer point par point, et surtout à apprendre à en parler en chinois avec des mots justes et une attitude respectueuse. C'est l'objectif de cette section : passer de la compréhension à la comparaison argumentée.

\courssub{Observer avant de comparer : deux phrases modèles}

Commençons par observer deux phrases simples qui résument à elles seules la différence d'usage entre les deux pays.

\begin{exemplebox}[Deux phrases à observer]
\zh{在喀麦隆，人们常常用现金。}\\
\emph{Zài Kāmàilóng, rénmen chángcháng yòng xiànjīn.}\\
« Au Cameroun, on utilise souvent les espèces. »

\medskip

\zh{在中国，手机支付很方便。}\\
\emph{Zài Zhōngguó, shǒujī zhīfù hěn fāngbiàn.}\\
« En Chine, le paiement mobile est très pratique. »
\end{exemplebox}

Observons la construction commune : \zh{在} + lieu, \zh{，} puis le constat. C'est la structure typique pour présenter un fait de société en chinois :

\begin{center}
\begin{tabularx}{\linewidth}{|X|X|X|}
\hline
\rowcolor{popblueL} Structure & Exemple & Traduction \\
\hline
在 + lieu + ， + constat & 在中国，手机支付很方便。 & En Chine, le paiement mobile est pratique. \\
\hline
Sujet + 常常 + verbe & 人们常常用现金。 & Les gens utilisent souvent les espèces. \\
\hline
Sujet + 很 + adjectif & 手机支付很方便。 & Le paiement mobile est très pratique. \\
\hline
\end{tabularx}
\end{center}

\begin{attention}[Le mot \zh{很} n'est pas toujours « très »]
Devant un adjectif seul, le chinois exige presque toujours \zh{很} (hěn), qui sert alors de simple liaison et ne veut pas forcément dire « très ». \zh{手机支付很方便。} peut se traduire « le paiement mobile est pratique ». Erreur fréquente des francophones : écrire \zh{手机支付方便。} qui, sans négation ni contexte particulier, sonne incomplet ou comparatif.
\end{attention}

\courssub{Fait comparé n° 1 : qui émet la monnaie ?}

\begin{definition}[Émetteur de la monnaie]
L'\cle{émetteur} est l'institution qui fabrique et met en circulation la monnaie d'un pays ou d'une zone.
\begin{itemize}
\item \textbf{Chine} : la Banque populaire de Chine, \zh{中国人民银行} (Zhōngguó Rénmín Yínháng), banque centrale chinoise, émet le renminbi.
\item \textbf{Cameroun} : la Banque des États de l'Afrique centrale, \zh{中非国家银行} (Zhōngfēi Guójiā Yínháng), la BEAC, émet le franc CFA pour les six pays de la CEMAC.
\end{itemize}
\end{definition}

La différence est fondamentale : la Chine possède \emph{sa propre} banque centrale nationale, tandis que le Cameroun \emph{partage} une banque centrale commune avec cinq autres pays. C'est ce qu'on exprime ainsi :

\zh{人民币是中国的钱，非洲法郎是好几个国家一起用的钱。}\\
\emph{Rénmínbì shì Zhōngguó de qián, Fēizhōu Fǎláng shì hǎojǐge guójiā yìqǐyòng de qián.}\\
« Le renminbi est la monnaie de la Chine ; le franc CFA est une monnaie utilisée en commun par plusieurs pays. »

Notez la structure \zh{好几个} (hǎojǐge) : « plusieurs, un bon nombre de », très utile à l'oral. Et \zh{一起用} (yìqǐyòng) : « utiliser ensemble ».

\courssub{Fait comparé n° 2 : les unités}

\begin{definition}[Les unités monétaires]
\begin{itemize}
\item \textbf{Chine} : le renminbi se divise en \zh{元} (yuán, unité principale), \zh{角} (jiǎo, un dixième de yuán) et \zh{分} (fēn, un centième de yuán). À l'oral, on dit souvent \zh{块} (kuài) pour \zh{元} et \zh{毛} (máo) pour \zh{角}.
\item \textbf{Cameroun} : le franc CFA (\zh{非洲法郎}, Fēizhōu Fǎláng) n'a pas de subdivision utilisée au quotidien ; on paie avec des billets de 500, 1 000, 2 000, 5 000 et 10 000 francs, et des pièces.
\end{itemize}
\end{definition}

\begin{center}
\begin{tabularx}{\linewidth}{|X|X|X|X|}
\hline
\rowcolor{popblueL} Caractères & Pinyin & Nature & Traduction \\
\hline
元 / 块 & yuán / kuài & nom / classificateur & yuan (unité chinoise) \\
\hline
角 / 毛 & jiǎo / máo & nom & jiao (1/10 de yuan) \\
\hline
分 & fēn & nom & fen (1/100 de yuan) \\
\hline
非洲法郎 & Fēizhōu Fǎláng & nom & franc CFA \\
\hline
现金 & xiànjīn & nom & espèces, argent liquide \\
\hline
钞票 & chāopiào & nom & billet de banque \\
\hline
硬币 & yìngbì & nom & pièce de monnaie \\
\hline
\end{tabularx}
\end{center}

Exemple de prix en chinois :

\zh{这个面包三块五。}\\
\emph{Zhège miànbāo sān kuài wǔ.}\\
« Ce pain coûte trois yuan cinq (3,5 yuan). »

\begin{attention}[Ne traduisez pas « franc » par \zh{法郎} seul]
\zh{法郎} (fǎláng) désigne le franc en général (franc français, franc suisse...). Pour le franc CFA, dites \zh{非洲法郎} (« franc africain »). De même, ne dites pas \zh{中国法郎} : la monnaie chinoise est \zh{人民币} (rénmínbì, « monnaie du peuple »), dont l'unité est \zh{元}.
\end{attention}

\courssub{Fait comparé n° 3 : comment paie-t-on au quotidien ?}

\begin{definition}[\zh{手机支付} — shǒujī zhīfù]
Le \cle{paiement par téléphone portable} : régler ses achats avec son téléphone, via des applications comme \zh{支付宝} (Zhīfùbǎo, Alipay) ou \zh{微信支付} (Wēixìn zhīfù, WeChat Pay) en Chine, ou le mobile money (MTN Mobile Money, Orange Money) au Cameroun.
\end{definition}

Voici le contraste d'usage, à présenter de façon factuelle :

\begin{itemize}
\item \textbf{Au Cameroun}, l'argent liquide (\zh{现金}) domine encore les achats quotidiens : marchés, transport, petits commerces. Mais le mobile money connaît un essor rapide : on envoie de l'argent à sa famille, on paie des factures par téléphone.
\item \textbf{En Chine}, le paiement mobile est très répandu, y compris pour les tout petits achats : un fruit au marché, un ticket de bus. Beaucoup de Chinois sortent sans portefeuille, seulement avec leur téléphone.
\end{itemize}

\zh{在中国，买水果也可以用手机支付。}\\
\emph{Zài Zhōngguó, mǎi shuǐguǒ yě kěyǐ yòng shǒujī zhīfù.}\\
« En Chine, même pour acheter des fruits, on peut payer par téléphone. »

\zh{在喀麦隆，很多人用手机付钱，也用现金。}\\
\emph{Zài Kāmàilóng, hěn duō rén yòng shǒujī fù qián, yě yòng xiànjīn.}\\
« Au Cameroun, beaucoup de gens paient par téléphone, et utilisent aussi les espèces. »

\courssub{Fait comparé n° 4 : changer de l'argent}

Quand on voyage entre les deux pays, il faut \cle{changer de l'argent} : \zh{换钱} (huàn qián). Le \cle{taux de change}, \zh{汇率} (huìlǜ), indique combien vaut une monnaie par rapport à l'autre, et il varie avec le temps.

\zh{去中国以前，我要换钱。}\\
\emph{Qù Zhōngguó yǐqián, wǒ yào huàn qián.}\\
« Avant d'aller en Chine, je dois changer de l'argent. »

\zh{今天的汇率怎么样？}\\
\emph{Jīntiān de huìlǜ zěnmeyàng?}\\
« Comment est le taux de change aujourd'hui ? »

\begin{popfigure}
\begin{tikzpicture}[node distance=1.6cm, every node/.style={font=\small}]
\node[draw=popblue, fill=popblueL, rounded corners, thick, minimum width=3.2cm, minimum height=1.1cm, align=center] (rmb) {人民币\\ \emph{Rénmínbì (yuan)}};
\node[draw=poporange, fill=poporangeL, rounded corners, thick, minimum width=3.2cm, minimum height=1.1cm, align=center, right=4.2cm of rmb] (cfa) {非洲法郎\\ \emph{franc CFA}};
\node[draw=popgreen, fill=popgreenL, rounded corners, thick, align=center] (huan) at ($(rmb)!0.5!(cfa)$) {换钱\\ \emph{huàn qián}};
\draw[-{Stealth}, very thick, popdark] (rmb) -- (huan);
\draw[-{Stealth}, very thick, popdark] (huan) -- (cfa);
\draw[{Stealth}-{Stealth}, very thick, poppurple, bend left=25] (rmb.north) to node[above, align=center]{汇率 taux de change\\ (il varie)} (cfa.north);
\end{tikzpicture}
\legende{Changer de l'argent : le taux de change relie les deux monnaies.}
\end{popfigure}

\courssub{Le grand tableau comparatif}

\begin{popfigure}
\begin{tikzpicture}[every node/.style={font=\small}]
\node[draw=popdark, fill=popblue, text=white, rounded corners=2pt, minimum width=3.4cm, minimum height=0.9cm, align=center] (t1) {Critère};
\node[draw=popdark, fill=popblue, text=white, rounded corners=2pt, minimum width=5.2cm, minimum height=0.9cm, align=center, right=0.15cm of t1] (t2) {中国 Chine};
\node[draw=popdark, fill=poporange, text=white, rounded corners=2pt, minimum width=5.2cm, minimum height=0.9cm, align=center, right=0.15cm of t2] (t3) {喀麦隆 Cameroun};
\node[draw=popline, fill=popcream, minimum width=3.4cm, minimum height=0.85cm, align=center, below=0.12cm of t1] (a1) {Monnaie};
\node[draw=popline, fill=white, minimum width=5.2cm, minimum height=0.85cm, align=center, below=0.12cm of t2] (a2) {人民币 (yuan)};
\node[draw=popline, fill=white, minimum width=5.2cm, minimum height=0.85cm, align=center, below=0.12cm of t3] (a3) {非洲法郎 (franc CFA)};
\node[draw=popline, fill=popcream, minimum width=3.4cm, minimum height=0.85cm, align=center, below=0.12cm of a1] (b1) {Émetteur};
\node[draw=popline, fill=white, minimum width=5.2cm, minimum height=0.85cm, align=center, below=0.12cm of a2] (b2) {中国人民银行\\ (banque nationale)};
\node[draw=popline, fill=white, minimum width=5.2cm, minimum height=0.85cm, align=center, below=0.12cm of a3] (b3) {中非国家银行 BEAC\\ (banque commune, 6 pays)};
\node[draw=popline, fill=popcream, minimum width=3.4cm, minimum height=0.85cm, align=center, below=0.12cm of b1] (c1) {Unités};
\node[draw=popline, fill=white, minimum width=5.2cm, minimum height=0.85cm, align=center, below=0.12cm of b2] (c2) {元 / 角 / 分};
\node[draw=popline, fill=white, minimum width=5.2cm, minimum height=0.85cm, align=center, below=0.12cm of b3] (c3) {billets de 500 à 10 000,\\ pièces};
\node[draw=popline, fill=popcream, minimum width=3.4cm, minimum height=0.85cm, align=center, below=0.12cm of c1] (d1) {Paiement courant};
\node[draw=popline, fill=white, minimum width=5.2cm, minimum height=0.85cm, align=center, below=0.12cm of c2] (d2) {手机支付 très répandu\\ (Alipay, WeChat Pay)};
\node[draw=popline, fill=white, minimum width=5.2cm, minimum height=0.85cm, align=center, below=0.12cm of c3] (d3) {现金 dominant,\\ mobile money en essor};
\node[draw=popline, fill=popcream, minimum width=3.4cm, minimum height=0.85cm, align=center, below=0.12cm of d1] (e1) {Particularité};
\node[draw=popline, fill=white, minimum width=5.2cm, minimum height=0.85cm, align=center, below=0.12cm of d2] (e2) {monnaie propre\\ à un seul pays};
\node[draw=popline, fill=white, minimum width=5.2cm, minimum height=0.85cm, align=center, below=0.12cm of d3] (e3) {monnaie partagée\\ par plusieurs pays};
\end{tikzpicture}
\legende{Tableau comparatif des deux monnaies.}
\end{popfigure}

\begin{savaistu}[Deux francs CFA différents]
Le franc CFA de la zone CEMAC — Cameroun, Tchad, Gabon, Congo, Guinée équatoriale, Centrafrique — est émis par la BEAC. Il existe un \emph{autre} franc CFA, celui de la zone UEMOA (Afrique de l'Ouest : Sénégal, Côte d'Ivoire, Mali...), émis par une autre banque centrale. Les deux ont la même valeur mais ne circulent pas dans les mêmes pays. En chinois, on peut dire \zh{中部非洲法郎} (franc CFA d'Afrique centrale) et \zh{西部非洲法郎} (franc CFA d'Afrique de l'Ouest).
\end{savaistu}

\courssub{Texte de synthèse : deux monnaies, deux réalités}

\begin{textebox}[两个国家，两种钱]
中国有自己的钱，叫人民币，是中国人民银行发的。人民币有元、角、分。现在在中国，人们买东西常常不用现金，用手机支付，支付宝和微信支付都很方便，买一杯茶也可以用手机付钱。\\
喀麦隆和五个国家一起用非洲法郎，是中非国家银行发的。在喀麦隆，人们常常用现金，但是现在用手机付钱也越来越多。\\
要去中国旅行的喀麦隆人要先换钱，因为人民币和非洲法郎不一样。汇率每天可能不一样。\\
这两种钱没有好坏，只是不一样：每个国家的钱都有自己的历史和用处。\\
\emph{Zhōngguó yǒu zìjǐ de qián, jiào rénmínbì, shì Zhōngguó Rénmín Yínháng fā de. Rénmínbì yǒu yuán, jiǎo, fēn. Xiànzài zài Zhōngguó, rénmen mǎi dōngxi chángcháng bú yòng xiànjīn, yòng shǒujī zhīfù, Zhīfùbǎo hé Wēixìn zhīfù dōu hěn fāngbiàn, mǎi yì bēi chá yě kěyǐ yòng shǒujī fù qián. Kāmàilóng hé wǔ ge guójiā yìqǐ yòng Fēizhōu Fǎláng, shì Zhōngfēi Guójiā Yínháng fā de. Zài Kāmàilóng, rénmen chángcháng yòng xiànjīn, dànshì xiànzài yòng shǒujī fù qián yě yuèláiyuè duō. Yào qù Zhōngguó lǚxíng de Kāmàilóng rén yào xiān huàn qián, yīnwèi rénmínbì hé Fēizhōu Fǎláng bù yíyàng. Huìlǜ měi tiān kěnéng bù yíyàng. Zhè liǎng zhǒng qián méiyǒu hǎo huài, zhǐshì bù yíyàng : měi ge guójiā de qián dōu yǒu zìjǐ de lìshǐ hé yòngchù.}\\
« La Chine a sa propre monnaie, appelée renminbi, émise par la Banque populaire de Chine. Le renminbi comprend le yuan, le jiao et le fen. Aujourd'hui en Chine, pour faire leurs achats, les gens n'utilisent souvent pas d'espèces : ils paient par téléphone ; Alipay et WeChat Pay sont très pratiques, et même pour acheter un thé, on peut payer avec son téléphone. Le Cameroun utilise, avec cinq autres pays, le franc CFA, émis par la BEAC. Au Cameroun, on utilise souvent les espèces, mais aujourd'hui le paiement par téléphone se développe de plus en plus. Le Camerounais qui veut voyager en Chine doit d'abord changer de l'argent, car le renminbi et le franc CFA sont différents. Le taux de change peut varier chaque jour. Ces deux monnaies ne sont ni bonnes ni mauvaises, elles sont simplement différentes : la monnaie de chaque pays a sa propre histoire et ses propres usages. »
\end{textebox}

\textbf{Glossaire du texte}

\begin{center}
\begin{tabularx}{\linewidth}{|X|X|X|X|}
\hline
\rowcolor{popblueL} Caractères & Pinyin & Nature & Traduction \\
\hline
发 & fā & verbe & émettre, envoyer \\
\hline
付钱 & fù qián & verbe + compl. & payer \\
\hline
越来越 & yuèláiyuè & adverbe & de plus en plus \\
\hline
换钱 & huàn qián & verbe + compl. & changer de l'argent \\
\hline
汇率 & huìlǜ & nom & taux de change \\
\hline
不一样 & bù yíyàng & expression & différent, pas pareil \\
\hline
用处 & yòngchù & nom & usage, utilité \\
\hline
\end{tabularx}
\end{center}

\textbf{Questions de compréhension}

\begin{enumerate}
\item \zh{人民币是谁发的？} Qui émet le renminbi ?
\item \zh{喀麦隆人用现金还是用手机支付？} Au Cameroun, utilise-t-on les espèces ou le paiement mobile ?
\item \zh{为什么去中国以前要换钱？} Pourquoi faut-il changer de l'argent avant d'aller en Chine ?
\end{enumerate}

\courssub{Parler des différences sans juger}

\begin{propriete}[Dire « différent », pas « meilleur »]
Pour comparer deux réalités socioculturelles avec respect, le chinois utilise \zh{不一样} (bù yíyàng, « pas pareil ») ou \zh{不同} (bùtóng, « différent »). Structure : \textbf{A 和 B 不一样。} (A hé B bù yíyàng. — « A et B sont différents. »). Évitez de traduire « meilleur / moins bon » par \zh{更好}/\zh{更坏} dans ce contexte : chaque système monétaire répond à l'histoire et aux besoins de son pays.
\end{propriete}

\zh{人民币和非洲法郎不一样，但是都很好用。}\\
\emph{Rénmínbì hé Fēizhōu Fǎláng bù yíyàng, dànshì dōu hěn hǎoyòng.}\\
« Le renminbi et le franc CFA sont différents, mais tous deux sont pratiques. »

\begin{attention}[Le piège du jugement de valeur]
En français, on glisse facilement vers « plus moderne », « moins développé ». En chinois comme en français, un exposé socioculturel exige des faits : dites \zh{不一样} (différent), \zh{每个国家有自己的历史} (chaque pays a sa propre histoire). Le paiement mobile dominant en Chine et les espèces dominantes au Cameroun sont deux réponses adaptées à deux contextes, ni « mieux » ni « moins bien ».
\end{attention}

\begin{exoresolu}[Comparer en chinois]
Énoncé : complétez avec \zh{不一样}, \zh{换钱} ou \zh{手机支付}, puis traduisez.\\
1. 人民币和非洲法郎\_\_\_\_\_\_。\\
2. 去中国旅行以前，我要\_\_\_\_\_\_。\\
3. 在中国，\_\_\_\_\_\_很方便。
\tcblower
Solution détaillée :\\
1. 人民币和非洲法郎\zh{不一样}。 \emph{Rénmínbì hé Fēizhōu Fǎláng bù yíyàng.} « Le renminbi et le franc CFA sont différents. » — Structure A 和 B 不一样。\\
2. 去中国旅行以前，我要\zh{换钱}。 \emph{Qù Zhōngguó lǚxíng yǐqián, wǒ yào huàn qián.} « Avant de voyager en Chine, je dois changer de l'argent. » — \zh{要} exprime ici la nécessité.\\
3. 在中国，\zh{手机支付}很方便。 \emph{Zài Zhōngguó, shǒujī zhīfù hěn fāngbiàn.} « En Chine, le paiement mobile est très pratique. » — Sujet + \zh{很} + adjectif.
\end{exoresolu}

\begin{aretenir}[L'essentiel de la section]
\begin{itemize}
\item Émetteurs : \zh{中国人民银行} (Chine, nationale) ; \zh{中非国家银行} BEAC (commune à 6 pays de la CEMAC).
\item Unités : \zh{元、角、分} ; franc CFA en billets et pièces.
\item Usages : \zh{现金} dominant au Cameroun, \zh{手机支付} très répandu en Chine.
\item Voyage : \zh{换钱} et \zh{汇率} (le taux varie).
\item Attitude : dire \zh{不一样}, sans jugement de valeur.
\end{itemize}
\end{aretenir}

\coursec{Exprimer son point de vue et débattre}

Après avoir comparé les faits objectifs sur le renminbi et le franc CFA dans la section précédente, nous allons maintenant apprendre à \textbf{formuler une opinion personnelle}, à \textbf{la justifier} et à \textbf{débattre} avec un camarade. En chinois, exprimer son point de vue suit des structures très régulières qu'il faut maîtriser pour l'oral (exposé, dialogue) et pour l'écrit (rédaction d'opinion), deux compétences exigées à l'examen.

\courssub{Observer : deux élèves comparent les moyens de paiement}

\begin{textebox}[Dialogue : Espèces ou mobile ?]
\zh{马林：你喜欢用现金还是手机支付？} \\
Mǎ Lín : Nǐ xǐhuan yòng xiànjīn háishi shǒujī zhīfù? \\
\textit{Marc : Tu préfères utiliser l'argent liquide ou le paiement mobile ?}

\zh{陈思：我觉得手机支付比现金方便，因为不用找零钱。} \\
Chén Sī : Wǒ juéde shǒujī zhīfù bǐ xiànjīn fāngbiàn, yīnwèi bù yòng zhǎo língqián. \\
\textit{Chen : Je trouve le paiement mobile plus pratique que l'argent liquide, parce qu'on n'a pas besoin de rendre la monnaie.}

\zh{马林：我不同意。我觉得现金更安全，因为手机没电了就不能付钱。} \\
Mǎ Lín : Wǒ bù tóngyì. Wǒ juéde xiànjīn gèng ānquán, yīnwèi shǒujī méi diàn le jiù bù néng fù qián. \\
\textit{Marc : Je ne suis pas d'accord. Je trouve l'argent liquide plus sûr, parce que si le téléphone n'a plus de batterie, on ne peut pas payer.}

\zh{陈思：你说得对。但是在中国，手机支付已经非常普及了。} \\
Chén Sī : Nǐ shuō de duì. Dànshì zài Zhōngguó, shǒujī zhīfù yǐjīng fēicháng pǔjí le. \\
\textit{Chen : Tu as raison. Mais en Chine, le paiement mobile est déjà très répandu.}
\end{textebox}

\begin{center}
\begin{tabularx}{\linewidth}{|X|X|X|X|}
\hline
\rowcolor{popblueL} \textbf{Caractères} & \textbf{Pinyin} & \textbf{Nature} & \textbf{Traduction} \\
\hline
现金 & xiànjīn & nom & argent liquide, espèces \\
\hline
手机支付 & shǒujī zhīfù & nom & paiement mobile \\
\hline
方便 & fāngbiàn & adj. & pratique, commode \\
\hline
找零钱 & zhǎo língqián & v. + obj. & rendre la monnaie (petite monnaie) \\
\hline
安全 & ānquán & adj. & sûr, en sécurité \\
\hline
没电 & méi diàn & v. + obj. & ne plus avoir de batterie \\
\hline
普及 & pǔjí & verbe / adj. & se répandre, se généraliser ; répandu \\
\hline
同意 & tóngyì & verbe & être d'accord \\
\hline
你说得对 & nǐ shuō de duì & expression & tu as raison (litt. « tu parles correctement ») \\
\hline
\end{tabularx}
\end{center}

\courssub{Structure fondamentale : 我觉得 + [opinion]}

\begin{propriete}[Exprimer son opinion : 我觉得 (wǒ juéde)]
La structure de base pour dire « Je pense que / Je trouve que / À mon avis » est :
\begin{center}
\textbf{Sujet + 觉得 + Phrase d'opinion}
\end{center}
\begin{itemize}
    \item \zh{觉得} (juéde) est un \textbf{verbe disyllabique inséparable} signifiant « trouver, avoir l'impression que, penser que ». Le deuxième caractère \zh{得} se prononce ici en \textbf{ton neutre} (de, court et léger).
    \item Contrairement à \zh{认为} (rènwéi — « considérer que », plus formel et plus écrit), \zh{觉得} est le standard à l'oral et dans l'écrit courant (HSK 3-4).
    \item La phrase d'opinion qui suit contient souvent un adjectif qualificatif (généralement précédé de \zh{很} hěn) ou une phrase complète.
\end{itemize}
\end{propriete}

\begin{exemplebox}[Exemples avec 我觉得]
\begin{itemize}
    \item \zh{我觉得人民币很漂亮。} Wǒ juéde rénmínbì hěn piàoliang. \textit{Je trouve le renminbi très joli.}
    \item \zh{我觉得在喀麦隆用现金更普遍。} Wǒ juéde zài Kāmàilóng yòng xiànjīn gèng pǔbiàn. \textit{Je trouve qu'au Cameroun, l'utilisation des espèces est plus répandue.}
    \item \zh{我觉得学汉字很有意思。} Wǒ juéde xué hànzì hěn yǒu yìsi. \textit{Je trouve l'apprentissage des caractères très intéressant.}
\end{itemize}
\end{exemplebox}

\begin{attention}[Piège majeur pour francophones : l'oubli ou la mauvaise écriture de \zh{得}]
\begin{itemize}
    \item \textbf{Erreurs fréquentes} : écrire \zh{我觉} (wǒ jué — verbe amputé) ou \zh{我觉的} (avec le \zh{的} de possession, qui n'a rien à faire ici).
    \item \textbf{Règle} : le verbe est \zh{觉得}, un mot \textbf{fixe de deux caractères} : on ne peut ni le raccourcir ni changer le second caractère. Attention à ne pas confondre ce \zh{得} (de, ton neutre) avec le \zh{得} structural des compléments de degré (\zh{说得对}) ni avec \zh{的}.
    \item \textbf{Prononciation} : \zh{jué} (2\up{e} ton, montant) + \zh{de} (ton neutre, bref). Ne dites pas « jou-é-dé » à la française.
\end{itemize}
\end{attention}

\courssub{Justifier son opinion : 因为 ... (yīnwèi)}

En français, on dit « Je pense que X \textbf{parce que} Y ». En chinois, la conjonction \zh{因为} (yīnwèi, « parce que ») introduit la cause, placée le plus souvent \textbf{après} l'opinion à l'oral (la structure complète \zh{因为... 所以...}, « puisque... donc... », vue en Première, est plus formelle).

\begin{propriete}[Justification : 我觉得 ..., 因为 ...]
\textbf{Structure :} \zh{我觉得 + [Opinion] ，因为 + [Raison]} \\
La virgule chinoise (，) sépare les deux parties.
\begin{itemize}
    \item \zh{因为} (yīnwèi) introduit la cause / la raison.
    \item La raison est une phrase complète (Sujet + Verbe / Adjectif), même si le sujet peut être sous-entendu quand il est évident.
\end{itemize}
\end{propriete}

\begin{exemplebox}[Opinion + Justification]
\begin{itemize}
    \item \zh{我觉得手机支付很快，因为不用排队。} Wǒ juéde shǒujī zhīfù hěn kuài, yīnwèi bù yòng páiduì. \textit{Je trouve le paiement mobile très rapide, parce qu'on n'a pas à faire la queue.}
    \item \zh{我觉得硬币不方便，因为太重了。} Wǒ juéde yìngbì bù fāngbiàn, yīnwèi tài zhòng le. \textit{Je trouve les pièces peu pratiques, parce qu'elles sont trop lourdes.}
    \item \zh{我觉得人民币很有文化，因为纸币上有不同的图案。} Wǒ juéde rénmínbì hěn yǒu wénhuà, yīnwèi zhǐbì shàng yǒu bùtóng de tú'àn. \textit{Je trouve que le renminbi est très culturel, parce qu'il y a différents motifs sur les billets.}
\end{itemize}
\end{exemplebox}

\begin{popfigure}
\begin{tikzpicture}[
    node distance=1.2cm and 1.5cm,
    box/.style={draw=popblue, thick, rounded corners, fill=popblueL, text width=3.5cm, align=center, minimum height=1.8cm},
    arrow/.style={-Stealth, thick, popdark}
]
\node[box, align=center] (sujet){\textbf{Sujet}\\\zh{我} (wǒ)};
\node[box, right=of sujet, align=center] (juede){\textbf{Verbe d'opinion}\\\zh{觉得} (juéde)};
\node[box, right=of juede, align=center] (opinion){\textbf{Opinion (phrase)}\\\zh{手机支付很方便}\\(shǒujī zhīfù hěn fāngbiàn)};
\node[box, below=of juede, xshift=1.5cm, fill=poporangeL, draw=poporange, align=center] (yinwei){\textbf{Conjonction}\\\zh{因为} (yīnwèi)};
\node[box, right=of yinwei, fill=poporangeL, draw=poporange, align=center] (raison){\textbf{Raison (phrase)}\\\zh{不用找零钱}\\(bù yòng zhǎo língqián)};

\draw[arrow] (sujet) -- (juede);
\draw[arrow] (juede) -- (opinion);
\draw[arrow] (juede) -- (yinwei);
\draw[arrow] (yinwei) -- (raison);
\end{tikzpicture}
\legende{Schéma de la structure : 我觉得 + Opinion ，因为 + Raison}
\end{popfigure}

\courssub{Réviser la comparaison avec 比 (bǐ)}

Pour nuancer son opinion, on compare souvent deux éléments. Rappel de la structure \zh{比} (bǐ) vue en Première et révisée en Terminale (Module 3) :

\begin{propriete}[Comparatif de supériorité : A 比 B + Adj.]
\textbf{Structure :} \zh{A + 比 + B + Adjectif (sans 很)} \\
Infériorité : \zh{A 没有 B + Adj.} (méiyǒu) — « A n'est pas aussi ... que B ». \\
Égalité : \zh{A 跟 B 一样 + Adj.} (gēn ... yīyàng) — « A est aussi ... que B ».
\end{propriete}

\begin{center}
\begin{tabularx}{\linewidth}{|X|X|X|}
\hline
\rowcolor{popblueL} \textbf{Structure} & \textbf{Exemple (Caractères + Pinyin)} & \textbf{Traduction} \\
\hline
A 比 B + Adj. & \zh{手机支付比现金方便。} (Shǒujī zhīfù bǐ xiànjīn fāngbiàn.) & Le mobile est plus pratique que les espèces. \\
\hline
A 比 B + 更 + Adj. & \zh{人民币比法郎更有意思。} (Rénmínbì bǐ fǎláng gèng yǒu yìsi.) & Le renminbi est encore plus intéressant que le franc. \\
\hline
A 没有 B + Adj. & \zh{现金没有手机支付快。} (Xiànjīn méiyǒu shǒujī zhīfù kuài.) & Les espèces ne sont pas aussi rapides que le mobile. \\
\hline
A 跟 B 一样 + Adj. & \zh{现金跟手机支付一样安全。} (Xiànjīn gēn shǒujī zhīfù yīyàng ānquán.) & Les espèces sont aussi sûres que le mobile. \\
\hline
\end{tabularx}
\end{center}

\begin{attention}[Erreur classique : \zh{很} après \zh{比}]
\begin{itemize}
    \item \textbf{Incorrect} : \zh{手机支付比现金很方便。}
    \item \textbf{Correct} : \zh{手机支付比现金方便。} (l'adjectif seul suffit ; \zh{很} est interdit dans le comparatif avec \zh{比}).
    \item Pour insister : \zh{更} (gèng, « encore plus ») avant l'adjectif, ou \zh{得多} (de duō, « beaucoup plus ») après l'adjectif : \zh{手机支付比现金方便得多。} (Shǒujī zhīfù bǐ xiànjīn fāngbiàn de duō.) \textit{Le mobile est beaucoup plus pratique que les espèces.}
\end{itemize}
\end{attention}

\courssub{Outils du débat : accord, désaccord, nuance}

Pour interagir, il faut des formules de réaction rapides et polies.

\begin{definition}[Vocabulaire de l'interaction orale]
\begin{center}
\begin{tabularx}{\linewidth}{|X|X|X|X|}
\hline
\rowcolor{popgreenL} \textbf{Caractères} & \textbf{Pinyin} & \textbf{Nature} & \textbf{Traduction / Contexte} \\
\hline
我同意 & wǒ tóngyì & expression & je suis d'accord \\
\hline
我不同意 & wǒ bù tóngyì & expression & je ne suis pas d'accord \\
\hline
我不太同意 & wǒ bù tài tóngyì & expression & je ne suis pas tout à fait d'accord (poli) \\
\hline
你说得对 & nǐ shuō de duì & expression & tu as raison \\
\hline
我也这么觉得 & wǒ yě zhème juéde & expression & je pense la même chose / moi aussi \\
\hline
不过 / 但是 & bùguò / dànshì & conj. & cependant / mais (introduit une nuance) \\
\hline
一方面... 另一方面... & yī fāngmiàn... lìng yī fāngmiàn... & structure & d'un côté... de l'autre côté... (HSK 4) \\
\hline
\end{tabularx}
\end{center}
\end{definition}

\begin{attention}[Politesse du désaccord]
En chinois comme en français, on évite de dire frontalement « tu as tort ». La formule \zh{你说得不对} (nǐ shuō de bù duì) existe mais sonne très direct, voire impoli entre camarades. Préférez :
\begin{itemize}
    \item \zh{我不太同意。} Wǒ bù tài tóngyì. \textit{Je ne suis pas tout à fait d'accord.}
    \item \zh{你说得有道理，不过...} Nǐ shuō de yǒu dàolǐ, bùguò... \textit{Ce que tu dis est sensé, cependant...}
\end{itemize}
\end{attention}

\begin{exemplebox}[Mini-échanges types]
\textbf{A :} \zh{我觉得现金更安全。} (Je trouve les espèces plus sûres.) \\
\textbf{B :} \zh{我同意，因为不用联网。} (D'accord, parce qu'on n'a pas besoin d'internet.) \\[4pt]
\textbf{A :} \zh{手机支付比现金快。} (Le mobile est plus rapide que les espèces.) \\
\textbf{B :} \zh{我不太同意。我觉得现金也很快，因为不需要打开软件。} (Pas tout à fait d'accord. Je trouve les espèces rapides aussi, parce qu'on n'a pas besoin d'ouvrir une application.) \\[4pt]
\textbf{A :} \zh{在中国，不带现金也没问题。} (En Chine, ne pas avoir d'espèces ne pose pas de problème.) \\
\textbf{B :} \zh{你说得对。但是在喀麦隆的农村，网络不太好。} (Tu as raison. Mais dans les zones rurales du Cameroun, le réseau n'est pas très bon.)
\end{exemplebox}

\courssub{Questions de réflexion pour la classe (expression orale guidée)}

Le professeur pose ces questions. Les élèves répondent en utilisant \zh{我觉得...，因为...} et \zh{比}.

\begin{experience}[Activité orale : « Enquête de classe »]
\textbf{Consigne :} en binômes, posez les 3 questions et répondez-y. Notez la réponse de votre partenaire pour un rapport oral devant la classe.
\begin{enumerate}
    \item \zh{问题一：你喜欢用现金还是手机支付？为什么？} (Q1 : Tu préfères les espèces ou le mobile ? Pourquoi ?)
    \item \zh{问题二：你觉得人民币和法郎，哪个更方便？} (Q2 : Selon toi, entre le renminbi et le franc, lequel est plus pratique ?)
    \item \zh{问题三：如果你去中国旅游，你怎么准备钱？} (Q3 : Si tu voyages en Chine, comment prépares-tu ton argent ?)
\end{enumerate}
\textbf{Aide lexicale pour la Q3 :} \zh{换钱} (huàn qián, changer de l'argent / de la devise), \zh{银行卡} (yínháng kǎ, carte bancaire), \zh{把银行卡绑定到微信或支付宝} (bǎ yínháng kǎ bǎngdìng dào Wēixìn huò Zhīfùbǎo, lier sa carte bancaire à WeChat ou Alipay).
\end{experience}

\courssub{Mini-débat guidé : « Espèces ou paiement mobile pour le Cameroun de demain ? »}

\begin{methode}[Organisation du mini-débat (15 min)]
\begin{enumerate}
    \item \textbf{Préparation (5 min) :} en groupes de 4. Deux élèves défendent \zh{现金} (les espèces), deux défendent \zh{手机支付} (le mobile). Chacun prépare 3 arguments avec \zh{我觉得...，因为...} et \zh{比}.
    \item \textbf{Débat (7 min) :} tour de table, un argument par élève. Les autres réagissent avec \zh{我同意 / 我不太同意 / 你说得对 / 不过...}.
    \item \textbf{Synthèse (3 min) :} un rapporteur par groupe résume en 2 phrases : \zh{我们组觉得...，因为...} (Notre groupe pense que..., parce que...).
\end{enumerate}
\end{methode}

\begin{exemplebox}[Arguments types pour le débat]
\textbf{Camp « 现金 » (espèces) :}
\begin{itemize}
    \item \zh{我觉得现金更安全，因为没有网络风险。} (Plus sûr, pas de risque lié au réseau.)
    \item \zh{现金比手机支付更普遍，因为农村没有网。} (Plus universel, car il n'y a pas de réseau au village.)
    \item \zh{我觉得用现金更容易控制预算，因为钱看得见。} (Plus facile de contrôler son budget, parce qu'on voit l'argent.)
\end{itemize}
\textbf{Camp « 手机支付 » (mobile) :}
\begin{itemize}
    \item \zh{我觉得手机支付比现金方便，因为不用带很厚的钱包。} (Plus pratique, pas besoin de porter un portefeuille épais.)
    \item \zh{手机支付更快，因为扫一扫就好了。} (Plus rapide, on scanne et c'est fini.)
    \item \zh{我觉得手机支付很方便查账，因为每一笔都有记录。} (Pratique pour vérifier ses comptes, car chaque paiement est enregistré.)
\end{itemize}
\end{exemplebox}

\courssub{Production guidée : mini-exposé comparatif (écrit / oral)}

C'est l'évaluation formative de la leçon. L'élève doit produire un paragraphe court (3 à 5 phrases) structuré.

\begin{methode}[Réussir son mini-exposé comparatif (3 phrases minimum)]
\begin{enumerate}
    \item \textbf{Présenter les deux sujets :} \zh{中国用人民币，喀麦隆用法郎。} / \zh{现在有两种支付方式：现金和手机支付。}
    \item \textbf{Citer une différence factuelle (vocabulaire socio) :} \zh{在中国，手机支付很普及；在喀麦隆，大家用现金比较多。} / \zh{人民币的单位有元、角、分，法郎的单位是法郎。}
    \item \textbf{Donner VOTRE opinion justifiée :} \zh{我觉得手机支付比现金方便，因为很快。} / \zh{我觉得现金比较好，因为不需要电。}
    \item \textbf{(Niveau +) Nuancer :} \zh{不过，现金也很重要。} / \zh{一方面...，另一方面...}
\end{enumerate}
\end{methode}

\begin{exoresolu}[Mini-exposé : « Mon avis sur les deux monnaies »]
\textbf{Consigne :} rédigez un court paragraphe (environ 60 à 80 caractères) comparant le renminbi et le franc CFA et donnant votre préférence. Utilisez \zh{觉得}、\zh{因为}、\zh{比}.

\tcblower
\textbf{Proposition de rédaction :}

\zh{中国用人民币，喀麦隆用法郎。我觉得人民币比较有意思，因为纸币上有少数民族的图案，很有文化。但是，我觉得手机支付比现金方便，因为不用找零钱。如果我去中国，我会用微信支付，因为很快。}

\textbf{Pinyin :}

Zhōngguó yòng rénmínbì, Kāmàilóng yòng fǎláng. Wǒ juéde rénmínbì bǐjiào yǒu yìsi, yīnwèi zhǐbì shàng yǒu shǎoshù mínzú de tú'àn, hěn yǒu wénhuà. Dànshì, wǒ juéde shǒujī zhīfù bǐ xiànjīn fāngbiàn, yīnwèi bù yòng zhǎo língqián. Rúguǒ wǒ qù Zhōngguó, wǒ huì yòng Wēixìn zhīfù, yīnwèi hěn kuài.

\textbf{Traduction :}

La Chine utilise le renminbi, le Cameroun utilise le franc. Je trouve le renminbi assez intéressant, parce que les billets portent des motifs représentant les minorités ethniques : c'est très culturel. Mais je trouve le paiement mobile plus pratique que les espèces, parce qu'on n'a pas besoin de rendre la monnaie. Si je vais en Chine, j'utiliserai WeChat Pay, parce que c'est très rapide.

\textbf{Analyse structurelle :}
\begin{itemize}
    \item Phrase 1 : contexte factuel (éléments socioculturels).
    \item Phrase 2 : opinion 1 (\zh{觉得} + \zh{因为}) sur le design culturel des billets.
    \item Phrase 3 : opinion 2 (\zh{觉得} + \zh{比} + \zh{因为}) sur l'usage pratique.
    \item Phrase 4 : projection personnelle (\zh{如果...，...会...}) avec justification.
\end{itemize}
\end{exoresolu}

\courssub{Exercices d'entraînement}

\begin{exercice}[Entraînement 1 : Complétez avec 觉得 / 因为 / 比 / 更 / 没有 / 如果 / 会]
Complétez les phrases en caractères (pinyin accepté si le caractère est oublié).
\begin{enumerate}
    \item 我 \_\_\_\_\_\_ 手机支付 \_\_\_\_\_\_ 现金方便，\_\_\_\_\_\_ 付钱很快。
    \item 你 \_\_\_\_\_\_ 法郎 \_\_\_\_\_\_ 人民币 \_\_\_\_\_\_ 漂亮吗？
    \item 我不同意，\_\_\_\_\_\_ 现金 \_\_\_\_\_\_ 手机支付安全。
    \item \_\_\_\_\_\_ 你去中国，你 \_\_\_\_\_\_ 准备现金还是手机支付？
\end{enumerate}
\end{exercice}

\begin{corrige}[Exercice 1]
\begin{enumerate}
    \item 我 \zh{觉得} 手机支付 \zh{比} 现金方便，\zh{因为} 付钱很快。 (Wǒ juéde shǒujī zhīfù bǐ xiànjīn fāngbiàn, yīnwèi fù qián hěn kuài.)
    \item 你 \zh{觉得} 法郎 \zh{比} 人民币 \zh{更} 漂亮吗？ (Nǐ juéde fǎláng bǐ rénmínbì gèng piàoliang ma?) — \textit{Rappel : pas de \zh{很} après \zh{比} ; on utilise \zh{更} pour insister.}
    \item 我不同意，\zh{因为} 现金 \zh{比} 手机支付安全。 (Wǒ bù tóngyì, yīnwèi xiànjīn bǐ shǒujī zhīfù ānquán.)
    \item \zh{如果} 你去中国，你 \zh{会} 准备现金还是手机支付？ (Rúguǒ nǐ qù Zhōngguó, nǐ huì zhǔnbèi xiànjīn háishi shǒujī zhīfù?) — \textit{Structure conditionnelle : \zh{如果...，...会...} (si..., alors...).}
\end{enumerate}
\end{corrige}

\begin{exercice}[Entraînement 2 : Traduction thème (français → chinois)]
Traduisez en caractères chinois + pinyin.
\begin{enumerate}
    \item Je trouve que le paiement mobile est plus pratique que l'argent liquide.
    \item Je ne suis pas d'accord. Je pense que l'argent liquide est plus sûr, parce qu'on n'a pas besoin d'internet.
    \item Tu as raison. Mais en Chine, tout le monde utilise le téléphone pour payer.
    \item Pour le Cameroun de demain, je préfère le paiement mobile, parce que c'est moderne.
\end{enumerate}
\end{exercice}

\begin{corrige}[Exercice 2]
\begin{enumerate}
    \item \zh{我觉得手机支付比现金方便。} (Wǒ juéde shǒujī zhīfù bǐ xiànjīn fāngbiàn.)
    \item \zh{我不同意。我觉得现金更安全，因为不需要网络。} (Wǒ bù tóngyì. Wǒ juéde xiànjīn gèng ānquán, yīnwèi bù xūyào wǎngluò.)
    \item \zh{你说得对。但是在中国，大家都用手机付钱。} (Nǐ shuō de duì. Dànshì zài Zhōngguó, dàjiā dōu yòng shǒujī fù qián.)
    \item \zh{对于喀麦隆的明天，我觉得手机支付比较好，因为它很现代。} (Duìyú Kāmàilóng de míngtiān, wǒ juéde shǒujī zhīfù bǐjiào hǎo, yīnwèi tā hěn xiàndài.) — \textit{Variante acceptée : \zh{为了喀麦隆的明天，我更喜欢手机支付，因为它很现代。}}
\end{enumerate}
\end{corrige}

\begin{savaistu}[Le « Hongbao » 红包 : l'enveloppe rouge numérique]
En Chine, la tradition de l'enveloppe rouge (\zh{红包} hóngbāo), offerte au Nouvel An ou lors des mariages, est massivement passée sur WeChat (\zh{微信红包} Wēixìn hóngbāo) depuis 2014.
\begin{itemize}
    \item \textbf{Fonctionnement :} on envoie une somme d'argent (fixe ou aléatoire) dans une conversation de groupe. Les membres « ouvrent » (\zh{抢} qiǎng, litt. « saisir ») l'enveloppe ; les plus rapides reçoivent les plus gros montants.
    \item \textbf{Impact socio-économique :} ce jeu a incité les personnes âgées à adopter WeChat Pay (pour recevoir les hóngbāo de leurs petits-enfants) et a familiarisé toute la population aux transferts d'argent instantanés entre particuliers.
    \item \textbf{Au Cameroun :} le Mobile Money (MTN MoMo, Orange Money) domine les transferts entre particuliers, mais l'aspect ludique et festif du \zh{红包} chinois n'a pas d'équivalent exact : l'usage camerounais du Mobile Money est plutôt utilitaire (retrait, paiement de factures, transfert familial).
\end{itemize}
C'est un bel exemple de \textbf{co-évolution entre culture et technologie} : une tradition sociale ancienne a servi de moteur à l'adoption du paiement numérique en Chine.
\end{savaistu}

\coursec{Production guidée : mini-exposé comparatif}

Après avoir appris à exprimer votre point de vue avec \zh{我觉得……因为……} et à débattre en classe, il est temps de rassembler tout ce que vous avez découvert dans cette leçon : le vocabulaire des monnaies, les faits comparés Cameroun / Chine et les structures d'opinion. Cette dernière étape vous conduit à produire un \cle{mini-exposé comparatif} : un court discours oral, préparé par écrit, présenté devant la classe. C'est exactement le type de production attendu à l'évaluation de fin d'année en Terminale.

\begin{aretenir}[Vocabulaire clé de la leçon (rappel)]
\begin{center}
\begin{tabularx}{\linewidth}{|X|X|X|X|}
\hline
\rowcolor{popblueL} Caractères & Pinyin & Nature & Traduction \\
\hline
钱 & qián & nom & argent, monnaie \\
人民币 & rénmínbì & nom propre & le renminbi (monnaie chinoise) \\
非洲法郎 & Fēizhōu fǎláng & nom propre & le franc CFA \\
中国人民银行 & Zhōngguó Rénmín Yínháng & nom propre & Banque populaire de Chine \\
中非国家银行 & Zhōngfēi Guójiā Yínháng & nom propre & Banque des États de l'Afrique centrale (BEAC) \\
发行 & fāxíng & verbe & émettre (monnaie) \\
手机支付 & shǒujī zhīfù & nom & paiement mobile \\
现金 & xiànjīn & nom & espèces, argent liquide \\
普遍 & pǔbiàn & adj. & répandu, généralisé \\
重要 & zhòngyào & adj. & important \\
\hline
\end{tabularx}
\end{center}
\end{aretenir}

\courssub{Observer d'abord : un exposé modèle}

Avant de produire, observons. Voici un exposé complet, tel qu'un élève de Terminale pourrait le présenter. Lisez-le, écoutez le professeur le lire, puis repérez les trois parties qui le composent.

\begin{textebox}[Exposé modèle : deux monnaies, deux réalités]
中国的钱叫人民币，喀麦隆的钱叫非洲法郎。它们不一样。\\
Zhōngguó de qián jiào rénmínbì, Kāmàilóng de qián jiào Fēizhōu fǎláng. Tāmen bù yíyàng.\\
« La monnaie de la Chine s'appelle le renminbi, celle du Cameroun le franc CFA. Elles sont différentes. »\\
中国人民银行发行人民币，中非国家银行发行非洲法郎。\\
Zhōngguó Rénmín Yínháng fāxíng rénmínbì, Zhōngfēi Guójiā Yínháng fāxíng Fēizhōu fǎláng.\\
« La Banque populaire de Chine émet le renminbi, la Banque des États de l'Afrique centrale émet le franc CFA. »\\
在中国，手机支付很普遍；在喀麦隆，现金很重要。\\
Zài Zhōngguó, shǒujī zhīfù hěn pǔbiàn; zài Kāmàilóng, xiànjīn hěn zhòngyào.\\
« En Chine, le paiement mobile est très répandu ; au Cameroun, les espèces sont importantes. »\\
我觉得两种钱都有自己的好处，因为人们可以用不同的方法买东西。\\
Wǒ juéde liǎng zhǒng qián dōu yǒu zìjǐ de hǎochù, yīnwèi rénmen kěyǐ yòng bùtóng de fāngfǎ mǎi dōngxi.\\
« Je pense que les deux monnaies ont leurs avantages, car on peut utiliser des méthodes différentes pour acheter des choses. »
\end{textebox}

Observez la construction : l'élève \textbf{présente} d'abord les deux monnaies, puis il \textbf{compare} (qui émet ? comment paie-t-on ?), et il \textbf{conclut} par son opinion personnelle justifiée. C'est la trame que vous allez suivre.

\courssub{La trame de l'exposé : forme et emploi}

\begin{propriete}[Structure de l'exposé comparatif en trois parties]
Un mini-exposé comparatif suit toujours le même plan :
\begin{enumerate}
\item \cle{Introduction} : présenter les deux objets comparés. Structure : \zh{A的钱是/叫……，B的钱是/叫……} (la monnaie de A est/s'appelle…, celle de B est/s'appelle…).
\item \cle{Comparaison} : deux ou trois faits, introduits par \zh{在中国，……；在喀麦隆，……} (en Chine… ; au Cameroun…). Le point-virgule chinois \zh{；} relie les deux faits parallèles.
\item \cle{Opinion} : \zh{我觉得……，因为……} (je pense que…, parce que…). L'opinion vient d'abord, la justification ensuite — jamais l'inverse en chinois.
\end{enumerate}
\end{propriete}

\begin{attention}[L'ordre opinion + justification]
En français, on peut dire « parce que…, je pense que… ». En chinois, dans cette structure d'exposé, on place \zh{我觉得} (je pense) \textbf{avant} \zh{因为} (parce que) : \zh{我觉得现金很方便，因为每个人都可以用它。} (\emph{Wǒ juéde xiànjīn hěn fāngbiàn, yīnwèi měi ge rén dōu kěyǐ yòng tā.}) « Je pense que les espèces sont pratiques, parce que tout le monde peut les utiliser. » Ne commencez pas votre conclusion par \zh{因为}.
\end{attention}

\begin{center}
\begin{tabularx}{\linewidth}{|X|X|X|}
\hline
\rowcolor{popblueL} Partie & Structure et exemple & Traduction \\
\hline
Introduction & 中国的钱是人民币，喀麦隆的钱是非洲法郎。\newline \emph{Zhōngguó de qián shì rénmínbì, Kāmàilóng de qián shì Fēizhōu fǎláng.} & La monnaie de la Chine est le renminbi, celle du Cameroun le franc CFA. \\
\hline
Émetteur & 中国人民银行发行人民币。\newline \emph{Zhōngguó Rénmín Yínháng fāxíng rénmínbì.} & La Banque populaire de Chine émet le renminbi. \\
\hline
Usage & 在中国，很多人用手机支付；在喀麦隆，人们常用现金。\newline \emph{Zài Zhōngguó, hěn duō rén yòng shǒujī zhīfù; zài Kāmàilóng, rénmen cháng yòng xiànjīn.} & En Chine, beaucoup de gens paient par téléphone ; au Cameroun, on utilise souvent les espèces. \\
\hline
Opinion & 我觉得手机支付很快，因为它很方便。\newline \emph{Wǒ juéde shǒujī zhīfù hěn kuài, yīnwèi tā hěn fāngbiàn.} & Je pense que le paiement mobile est rapide, parce qu'il est pratique. \\
\hline
\end{tabularx}
\end{center}

\begin{popfigure}
\begin{tikzpicture}[
node distance=0.9cm,
etape/.style={rectangle, rounded corners=4pt, draw=popblue, fill=popblueL, align=center, text width=5.2cm, minimum height=1.1cm, font=\small},
fleche/.style={-{Stealth[length=3mm]}, thick, poporange}
]
\node[etape, align=center] (intro){\textbf{Introduction}\\ 中国的钱……喀麦隆的钱……\\ \emph{présenter les deux monnaies}};
\node[etape, below=of intro, align=center] (comp){\textbf{Comparaison}\\ 发行 / 手机支付 / 现金\\ \emph{émetteur et usage}};
\node[etape, below=of comp, align=center] (op){\textbf{Opinion}\\ 我觉得……因为……\\ \emph{point de vue justifié}};
\draw[fleche] (intro) -- (comp);
\draw[fleche] (comp) -- (op);
\end{tikzpicture}
\legende{Organigramme de l'exposé : introduction, comparaison, opinion.}
\end{popfigure}

\courssub{Le déroulement en trois étapes}

\begin{methode}[Réussir son mini-exposé : les trois étapes]
\begin{enumerate}
\item \textbf{Étape 1 — Écrire (10 minutes, individuel).} Complétez la trame par écrit dans votre cahier : remplacez les pointillés par vos propres idées. Écrivez les caractères, puis le pinyin au crayon pour préparer la prononciation.
\item \textbf{Étape 2 — Répéter (en binôme).} Lisez votre exposé à votre voisin. Le voisin écoute et corrige les \cle{tons} : vérifiez surtout \zh{钱} qián (2\up{e} ton), \zh{发行} fāxíng (1\up{er} puis 2\up{e} ton), \zh{普遍} pǔbiàn (3\up{e} puis 4\up{e} ton). Échangez ensuite les rôles.
\item \textbf{Étape 3 — Présenter (devant la classe).} Les volontaires présentent leur exposé en 1 à 2 minutes. Parlez \textbf{lentement}, regardez l'auditoire (pas seulement votre feuille), et utilisez le tableau comparatif du cours comme support visuel : montrez la ligne dont vous parlez.
\end{enumerate}
\end{methode}

\begin{exemplebox}[Critères de réussite affichés au tableau]
Votre exposé est réussi si :
\begin{itemize}
\item le vocabulaire est correct : \zh{钱} (argent/monnaie), \zh{人民币} (renminbi), \zh{非洲法郎} (franc CFA) sont employés à bon escient ;
\item la structure \zh{我觉得……因为……} est présente et complète (opinion + justification) ;
\item les tons sont prononcés correctement, en particulier sur les mots du thème ;
\item le débit est calme et l'auditoire est regardé.
\end{itemize}
\end{exemplebox}

\courssub{S'entraîner avant de présenter}

\begin{exoresolu}[Compléter la trame]
Énoncé : complétez la trame suivante avec les mots de la leçon, puis traduisez votre phrase d'opinion.\\
\zh{中国的钱是＿＿＿，喀麦隆的钱是＿＿＿。在中国，很多人用＿＿＿支付；在喀麦隆，人们常用＿＿＿。我觉得＿＿＿，因为＿＿＿。}
\tcblower
Solution détaillée :\\
\zh{中国的钱是人民币，喀麦隆的钱是非洲法郎。在中国，很多人用手机支付；在喀麦隆，人们常用现金。我觉得现金很重要，因为不是每个人都有手机。}\\
\emph{Zhōngguó de qián shì rénmínbì, Kāmàilóng de qián shì Fēizhōu fǎláng. Zài Zhōngguó, hěn duō rén yòng shǒujī zhīfù; zài Kāmàilóng, rénmen cháng yòng xiànjīn. Wǒ juéde xiànjīn hěn zhòngyào, yīnwèi bú shì měi ge rén dōu yǒu shǒujī.}\\
« La monnaie de la Chine est le renminbi, celle du Cameroun le franc CFA. En Chine, beaucoup de gens paient par téléphone ; au Cameroun, on utilise souvent les espèces. Je pense que les espèces sont importantes, parce que tout le monde n'a pas de téléphone portable. »\\
Méthode : les quatre premiers blancs sont des faits (recopiez le cours) ; les deux derniers sont personnels — toute opinion correcte en grammaire est acceptée.
\end{exoresolu}

\begin{attention}[Enrichissement pour l'examen : 普遍 et 重要]
Deux adjectifs utiles pour enrichir votre exposé et briller à l'évaluation : \zh{普遍} pǔbiàn « répandu, courant » (\zh{手机支付很普遍。} « le paiement mobile est très répandu ») et \zh{重要} zhòngyào « important » (\zh{现金很重要。} « les espèces sont importantes »). Attention aux tons : pǔbiàn = 3\up{e} + 4\up{e} ton ; zhòngyào = 4\up{e} + 4\up{e} ton. Ne confondez pas \zh{重要} zhòngyào (important) et \zh{重} zhòng (lourd).
\end{attention}

\begin{exercice}[À vous de parler]
\begin{enumerate}
\item Complétez la trame de l'exposé avec votre propre opinion, puis écrivez le pinyin sous chaque phrase.
\item En binôme, répétez votre exposé ; notez deux erreurs de tons corrigées par votre voisin.
\item Présentez votre exposé devant la classe en 1 à 2 minutes, en montrant le tableau comparatif comme support.
\end{enumerate}
\end{exercice}

\begin{corrige}[Exercice n°1]
Éléments de correction : la trame complétée doit contenir \zh{人民币} et \zh{非洲法郎} dans l'introduction, un fait sur l'émetteur (\zh{中国人民银行发行人民币}) ou sur l'usage (\zh{手机支付} / \zh{现金}) dans la comparaison, et une phrase complète \zh{我觉得……因为……} dans la conclusion. Exemple d'opinion acceptée : \zh{我觉得人民币很有意思，因为它的名字是«人民的钱»。} (\emph{Wǒ juéde rénmínbì hěn yǒuyìsi, yīnwèi tā de míngzi shì «rénmín de qián».}) « Je pense que le renminbi est intéressant, parce que son nom signifie "la monnaie du peuple". » Le professeur vérifie les tons de \zh{钱} qián, \zh{发行} fāxíng et \zh{觉得} juéde.
\end{corrige}

\courssub{Trace écrite}

Recopiez dans votre cahier le tableau comparatif suivant, qui résume toute la leçon et sert de support visuel à l'exposé.

\begin{center}
\begin{tabularx}{\linewidth}{|X|X|X|}
\hline
\rowcolor{popblueL} Point comparé & Chine 中国 & Cameroun 喀麦隆 \\
\hline
Monnaie & 人民币 rénmínbì (le renminbi) & 非洲法郎 Fēizhōu fǎláng (le franc CFA) \\
\hline
Émetteur & 中国人民银行 Zhōngguó Rénmín Yínháng & 中非国家银行 Zhōngfēi Guójiā Yínháng \\
\hline
Usage courant & 手机支付 shǒujī zhīfù (paiement mobile, très répandu 普遍) & 现金 xiànjīn (espèces, très importantes 重要) \\
\hline
Mon opinion & \multicolumn{2}{X|}{我觉得……因为…… (Wǒ juéde… yīnwèi…)} \\
\hline
\end{tabularx}
\end{center}

\begin{savaistu}[Comparer avec respect : une compétence évaluée]
Savoir comparer deux cultures de façon factuelle et respectueuse — dire « les deux monnaies ont leurs avantages » plutôt que juger — est une compétence explicitement attendue à l'évaluation de fin d'année en Terminale. En chinois, la phrase \zh{它们不一样，但是都很好。} (\emph{Tāmen bù yíyàng, dànshì dōu hěn hǎo.}) « Elles sont différentes, mais toutes les deux sont bien » est une excellente conclusion d'exposé : elle montre ouverture d'esprit et maîtrise de la langue.
\end{savaistu}

\newpage

\coursec{Activité d'intégration}

\coursec{Leçon 29 — Éléments socioculturels : différence entre la monnaie chinoise et la monnaie camerounaise}

\courssub{Objectifs de la leçon}
À l'issue de cette leçon, tu seras capable de :
\begin{itemize}
\item nommer et écrire les unités monétaires de la Chine (\zh{人民币}、\zh{元}、\zh{角}、\zh{分}) et du Cameroun (\zh{非洲法郎}) ;
\item décrire les usages contrastés : prédominance des \zh{现金} (espèces) au Cameroun vs généralisation du \zh{手机支付} (paiement mobile) en Chine ;
\item employer les structures grammaticales de comparaison (\zh{比}), d'opinion (\zh{我觉得}) et de justification (\zh{因为……所以……}) dans un contexte économique ;
\item réaliser un jeu de rôle oral (touriste, agent de change, journaliste) et rédiger un mini-exposé écrit structuré de 5 phrases ;
\item adopter une posture \cle{factuelle et respectueuse} pour comparer deux réalités socioculturelles sans hiérarchisation de valeur.
\end{itemize}

\courssub{Découverte : deux monnaies, deux réalités}

\begin{textebox}[Panneau du bureau de change — Affiche à l'entrée de l'aéroport Yaoundé-Nsimalen]
换钱处 — Bureau de change\\
Huàn qián chù\\
« Bureau de change »\\[0.4em]
人民币 $\leftrightarrow$ 非洲法郎\\
Rénmínbì $\leftrightarrow$ Fēizhōu fǎláng\\
« Renminbi $\leftrightarrow$ Franc CFA »\\[0.4em]
欢迎！请问，您要换多少钱？\\
Huānyíng! Qǐngwèn, nín yào huàn duōshao qián?\\
« Bienvenue ! Combien souhaitez-vous changer ? »\\[0.4em]
提示：在喀麦隆常用现金；在中国普及手机支付。\\
Tíshì: Zài Kāmàilóng cháng yòng xiànjīn; zài Zhōngguó pǔjí shǒujī zhīfù.\\
« Info : au Cameroun on utilise souvent les espèces ; en Chine le paiement mobile est généralisé. »
\end{textebox}

\courssub{Glossaire et écriture des caractères}

\begin{center}
\begin{tabularx}{\linewidth}{|X|X|X|X|}
\hline
\rowcolor{popblueL} Caractères & Pinyin & Nature & Traduction \\
\hline
人民币 & rénmínbì & n. & Renminbi (monnaie officielle chinoise) \\
\hline
元 & yuán & classif./n. & Yuan (unité principale, formel) \\
\hline
块 & kuài & classif. & Yuan (unité courante, oral) \\
\hline
角 & jiǎo & classif./n. & Jiao (1/10 de yuan, formel) \\
\hline
毛 & máo & classif. & Jiao (1/10 de yuan, oral) \\
\hline
分 & fēn & classif./n. & Fen (1/100 de yuan) \\
\hline
非洲法郎 & Fēizhōu fǎláng & n. & Franc CFA (monnaie zone CEMAC) \\
\hline
换钱 & huàn qián & v. + compl. & Changer de l'argent \\
\hline
兑换 & duìhuàn & v. & Convertir (devises) \\
\hline
汇率 & huìlǜ & n. & Taux de change \\
\hline
现金 & xiànjīn & n. & Espèces, argent liquide \\
\hline
手机支付 & shǒujī zhīfù & n. & Paiement par téléphone (mobile) \\
\hline
微信支付 / 支付宝 & Wēixìn zhīfù / Zhīfùbǎo & n. & WeChat Pay / Alipay \\
\hline
银行 & yínháng & n. & Banque \\
\hline
方便 & fāngbiàn & adj. & Pratique, commode \\
\hline
普及 & pǔjí & v. & Être généralisé, se répandre \\
\hline
习惯 & xíguàn & n. & Habitude, coutume \\
\hline
不一样 & bù yíyàng & adj. & Différent, pas pareil \\
\hline
\end{tabularx}
\end{center}

\begin{definition}[Écriture des caractères clés — Ordre des traits et radicaux]
\textbf{币 (bì) :} Radicale \zh{巾} (jīn, tissu) à gauche. Ordre : haut → bas (le haut de \zh{巾}), puis \zh{巾} (3 traits), puis \zh{戈} (gē, hallebarde) à droite (4 traits : horizontale, point, horizontale crochue, point). \textit{Sens : monnaie, billet.}\\
\textbf{元 (yuán) :} Radicale \zh{一} (yī, un) ou \zh{儿} (ér) selon dictionnaires. Ordre : horizontale supérieure, puis \zh{二} (deux horizontales), puis trait final descendant \zh{㇏} (nà). \textit{Sens : origine, unité monétaire principale.}\\
\textbf{法 (fǎ) :} Radicale \zh{氵} (shuǐ, eau) à gauche (3 traits : deux points, trait montant). Partie droite \zh{去} (qù) : \zh{土} (tǔ, terre) + \zh{厶} (sī, privé). \textit{Sens : loi, méthode ; ici transcription phonétique « franc ».}\\
\textbf{郎 (láng) :} Radicale \zh{阝} (fù, monticule) à droite (3 traits). Partie gauche \zh{良} (liáng) : \zh{艹} (cǎo, herbe) + \zh{良} (genoux + un). \textit{Sens : jeune homme, monsieur ; transcription « long » dans franc.}\\
\textbf{支 (zhī) :} Radicale \zh{手} (shǒu, main) \zh{扌} à gauche (3 traits). Partie droite \zh{枝} (zhī) simplifiée : \zh{十} + \zh{又}. \textit{Sens : soutenir, branche ; classificateur pour objets longs, verbes « payer » (支付).}\\
\textbf{付 (fù) :} Radicale \zh{亻} (rén, homme) à gauche (2 traits). Partie droite \zh{寸} (cùn, pouce/main). \textit{Sens : payer, remettre.}
\end{definition}

\begin{aretenir}[Unités monétaires chinoises : formel vs oral]
En chinois, on distingue le registre formel (écrit, bancaire) et le registre oral (quotidien) :
\begin{center}
\begin{tabularx}{\linewidth}{|X|X|X|}
\hline
\rowcolor{poporangeL} Valeur & Formel (écrit/banque) & Oral (quotidien) \\
\hline
1 yuan & \zh{一元} (yī yuán) & \zh{一块} (yī kuài) / \zh{一块钱} (yī kuài qián) \\
\hline
0,1 yuan (1/10) & \zh{一角} (yī jiǎo) & \zh{一毛} (yī máo) / \zh{一毛钱} (yī máo qián) \\
\hline
0,01 yuan (1/100) & \zh{一分} (yī fēn) & \zh{一分} (yī fēn) / \zh{一分钱} (yī fēn qián) \\
\hline
\end{tabularx}
\end{center}
\emph{Exemple :} 12,35 元 s'écrit \zh{十二元三角五分} (shí'èr yuán sān jiǎo wǔ fēn) mais se dit \zh{十二块三毛五} (shí'èr kuài sān máo wǔ). Le franc CFA ne se subdivise plus en centimes dans l'usage courant.
\end{aretenir}

\courssub{Compréhension du texte (Panneau \& Contexte)}

\begin{exercice}[Questions de compréhension]
Réponds en français (ou en chinois si possible) d'après le panneau \zh{换钱处} :
\begin{enumerate}
\item Quelles sont les deux monnaies affichées pour l'échange ?
\item Quelle indication socioculturelle est donnée aux clients sur les usages de paiement ?
\item Traduire en chinois : « Je veux changer 1000 yuans en francs CFA. » (Indice : \zh{我想换...}).
\item Pourquoi l'affiche mentionne-t-elle \zh{手机支付} alors qu'on est dans un bureau de change \emph{physique} ?
\end{enumerate}
\end{exercice}

\begin{corrige}[Exercice Compréhension]
\begin{enumerate}
\item \zh{人民币} (Renminbi / Yuan) et \zh{非洲法郎} (Franc CFA).
\item « Au Cameroun on utilise souvent les espèces (\zh{现金}) ; en Chine le paiement mobile (\zh{手机支付}) est généralisé (\zh{普及}). »
\item \zh{我想换一千元人民币换非洲法郎。} (Wǒ xiǎng huàn yīqiān yuán rénmínbì huàn Fēizhōu fǎláng.)
\item Pour sensibiliser le touriste chinois (habitué au tout-mobile) qu'au Cameroun, il lui \textbf{faut} des espèces (\zh{现金}) pour la plupart des transactions (marchés, taxis, petits commerces).
\end{enumerate}
\end{corrige}

\courssub{Faits comparés Cameroun / Chine}

\begin{center}
\begin{tabularx}{\linewidth}{|X|X|X|}
\hline
\rowcolor{popgreenL} Dimension & Cameroun (Zone CEMAC) & Chine (RPC) \\
\hline
Monnaie officielle & Franc CFA (\zh{非洲法郎} Fēizhōu fǎláng) & Renminbi (\zh{人民币} Rénmínbì) \\
\hline
Unité principale & Franc (pas de subdivision courante) & Yuan (\zh{元} yuán / \zh{块} kuài) \\
\hline
Sous-unités & Centimes (théoriques, disparus) & Jiao (\zh{角} jiǎo / \zh{毛} máo), Fen (\zh{分} fēn) \\
\hline
Billets courants & 500, 1000, 2000, 5000, 10000 FCFA & 1, 5, 10, 20, 50, 100 yuans \\
\hline
Paiement dominant & \zh{现金} (Espèces) quasi exclusif hors banques/grandes surfaces & \zh{手机支付} (Mobile : WeChat Pay, Alipay) > 90\% transactions urbaines \\
\hline
Rôle de la banque & Retrait/salaire, peu de paiement direct & Compte lié aux apps, virements instantanés P2P \\
\hline
Pourboire / Petit commerce & Espèces obligatoires & Code QR affiché, paiement 0,01 yuan possible \\
\hline
Taux de change (repère 2024) & 1 EUR $\approx$ 655,957 FCFA (fixe) & 1 EUR $\approx$ 7,8 CNY (flottant géré) \\
\hline
\end{tabularx}
\end{center}

\begin{savaistu}[Le Franc CFA et le Renminbi : deux histoires monétaires]
\textbf{Franc CFA (Communauté Financière Africaine) :} Créé en 1945. Au Cameroun (zone CEMAC), il est émis par la BEAC (Banque des États de l'Afrique Centrale), siège à Yaoundé. Parité fixe avec l'Euro (1 EUR = 655,957 FCFA) garantie par le Trésor français. Billets communs aux 6 pays CEMAC (Cameroun, Tchad, RCA, Gabon, Guinée équatoriale, Congo) distingués par une lettre de série (ex: \textbf{U} = Cameroun).\\
\textbf{Renminbi (Monnaie du Peuple) :} Lancée en 1948 par la Banque Populaire de Chine (中央银行 Zhōngyāng Yínháng). Unité : Yuan. Depuis 2005, taux de change flottant géré (panier de devises). Internationalisation progressive (inclusion dans le panier DTS du FMI en 2016). Billets à l'effigie de Mao Zedong (depuis 1999, série 2019 actuelle).
\end{savaistu}

\begin{popfigure}
\begin{tikzpicture}[
    node distance=1.5cm and 2cm,
    box/.style={draw=popblue, thick, rounded corners, fill=popblueL, text width=3.5cm, align=center, font=\small},
    arrow/.style={-Stealth, thick, popdark}
]
\node[box, align=center] (cm){Cameroun\\ \zh{非洲法郎}\\ Espèces roi};
\node[box, right=of cm, align=center] (ch){Chine\\ \zh{人民币}\\ Mobile first};
\node[box, below=of cm, align=center] (usage){Contexte : Marchés, taxis,\\ secteur informel dominant};
\node[box, below=of ch, align=center] (tech){Contexte : Urbanisation rapide,\\ écosystème WeChat/Alipay};
\draw[arrow] (cm) -- node[above, font=\tiny] {Taux fixe vs Euro} (ch);
\draw[arrow] (usage) -- node[above, font=\tiny] {Habitudes culturelles} (tech);
\draw[arrow, poporange] (cm.south) -- ++(0,-0.5) -| (usage.north);
\draw[arrow, poporange] (ch.south) -- ++(0,-0.5) -| (tech.north);
\end{tikzpicture}
\legende{Comparaison schématique : monnaies et écosystèmes de paiement}
\end{popfigure}

\courssub{Point de grammaire : structures pour comparer et argumenter}

\begin{propriete}[Structure comparative avec \cle{比} (bǐ)]
\textbf{Schéma :} Sujet A + \zh{比} + Sujet B + Adjectif (sv) \textbf{(sans \zh{的}, sans \zh{更} obligatoire)}.\\
\textbf{Règle :} A possède le degré supérieur de la qualité exprimée par l'adjectif. Négation : \zh{没有} (méiyǒu) + B + \zh{那么/这么} + Adj. (A n'est pas aussi... que B).\\
\textbf{Exemples :}
\begin{itemize}
\item \zh{中国的手机支付比喀麦隆普及。} (Zhōngguó de shǒujī zhīfù bǐ Kāmàilóng pǔjí.) « Le paiement mobile est plus généralisé en Chine qu'au Cameroun. »
\item \zh{喀麦隆的现金使用比中国多。} (Kāmàilóng de xiànjīn shǐyòng bǐ Zhōngguó duō.) « L'usage des espèces est plus important au Cameroun qu'en Chine. »
\item \zh{非洲法郎没有人民币那么国际化。} (Fēizhōu fǎláng méiyǒu rénmínbì nàme guójìhuà.) « Le franc CFA n'est pas aussi internationalisé que le renminbi. »
\end{itemize}
\emph{Attention :} L'adjectif ne prend pas de marque de comparatif (pas de \zh{更} gèng obligatoire, pas de \zh{-er}).
\end{propriete}

\begin{propriete}[Expression de l'opinion : \cle{我觉得} (wǒ juéde)]
\textbf{Schéma :} \zh{我觉得} + Proposition complète (Sujet + Verbe/Adj...).\\
\textbf{Emploi :} Exprimer un sentiment, une impression, un jugement subjectif. Plus doux que \zh{我认为} (wǒ rènwéi, « je considère que »).\\
\textbf{Exemples :}
\begin{itemize}
\item \zh{我觉得手机支付很方便。} (Wǒ juéde shǒujī zhīfù hěn fāngbiàn.) « Je trouve le paiement mobile très pratique. »
\item \zh{我觉得喀麦隆的钱颜色很漂亮。} (Wǒ juéde Kāmàilóng de qián yánsè hěn piàoliang.) « Je trouve que l'argent camerounais a de jolies couleurs. »
\end{itemize}
\end{propriete}

\begin{propriete}[Justification cause-conséquence : \cle{因为……所以……} (yīnwèi... suǒyǐ...)]
\textbf{Schéma :} \zh{因为} + Cause (phrase), \zh{所以} + Conséquence (phrase).\\
\textbf{Ordre :} Cause $\rightarrow$ Conséquence (logique déductive). En français : « Parce que..., donc... » ou « Comme..., ... ».\\
\textbf{Variante orale :} \zh{因为}... (cause), (conséquence) — \zh{所以} souvent omis à l'oral.\\
\textbf{Exemples :}
\begin{itemize}
\item \zh{因为市场上不收手机支付，所以我必须带现金。} (Yīnwèi shìchǎng shang bù shōu shǒujī zhīfù, suǒyǐ wǒ bìxū dài xiànjīn.) « Comme le marché n'accepte pas le paiement mobile, je dois emporter des espèces. »
\item \zh{因为汇率固定，所以换算很容易。} (Yīnwèi huìlǜ gùdìng, suǒyǐ huànsuàn hěn róngyì.) « Parce que le taux est fixe, la conversion est facile. »
\end{itemize}
\end{propriete}

\begin{propriete}[Place du complément de lieu \cle{在} (zài)]
\textbf{Règle absolue :} Le complément de lieu (ou de temps) se place \textbf{AVANT} le verbe (ou l'adjectif verbal).\\
\textbf{Schéma :} Sujet + \zh{在} + Lieu + Verbe + Complément.\\
\textbf{Correct :} \zh{在喀麦隆，人们用现金。} (Zài Kāmàilóng, rénmen yòng xiànjīn.)\\
\textbf{Incorrect :} \zh{人们用现金在喀麦隆。} (Influence français « Les gens utilisent des espèces au Cameroun »).
\end{propriete}

\begin{attention}[Pièges fréquents pour francophones]
\begin{itemize}
\item \textbf{Confusion \zh{元} / \zh{块} :} N'écris pas \zh{块} dans un texte formel (examen, lettre officielle) ; ne dis pas \zh{元} (yuán) dans la rue pour acheter du pain (dis \zh{块} kuài).
\item \textbf{Traduction de « argent » :} \zh{钱} (qián) = argent (général, espèces, monnaie). \zh{人民币} = Renminbi (la monnaie en tant que système). Ne dis pas \zh{中国的钱叫元} mais \zh{中国的货币叫人民币，单位是元} (La monnaie de Chine s'appelle Renminbi, l'unité est le Yuan).
\item \textbf{Tons de \zh{非洲法郎} :} \zh{非} (fēi, 1), \zh{洲} (zhōu, 1), \zh{法} (fǎ, 3), \zh{郎} (láng, 2 \textbf{montant}). Erreur fréquente : dire \zh{làng} (4e ton, vague).
\item \textbf{Supériorité culturelle interdite :} Évite \zh{中国比喀麦隆先进} (La Chine est plus avancée). Préfère \zh{中国和喀麦隆的支付习惯不一样} (Les habitudes de paiement sont différentes) + faits (\zh{比} + adj. factuel : \zh{普及}, \zh{多}, \zh{少}).
\end{itemize}
\end{attention}

\courssub{Exprimer son point de vue et débattre : le Bureau de change interculturel}

\begin{experience}[Mise en situation : « Bureau de change interculturel » — Jeu de rôles (3 pers.)]
\textbf{Contexte :} Semaine des langues, lycée de Yaoundé. Salle transformée en bureau de change aéroport Yaoundé-Nsimalen.
\textbf{Rôles :}
\begin{enumerate}
\item \textbf{Li Na (Touriste chinoise)} : Arrive avec des yuans en espèces (\zh{现金}), veut des FCFA. Pose des questions sur l'usage de l'argent ici.
\item \textbf{M. Mbarga (Agent de change camerounais)} : Accueille, explique le taux, donne des conseils pratiques (petits billets, marchés, taxis).
\item \textbf{Journaliste CRTV} : Interviewe les deux pour un reportage « L'argent ici et là-bas ».
\end{enumerate}
\textbf{Déroulement (25 min total) :}
\begin{enumerate}
\item \textbf{Préparation groupe (10 min) :} Choisir les rôles. Sur le glossaire, surligner \textbf{8 mots} obligatoires à placer dans le sketch (ex: \zh{人民币、元、非洲法郎、现金、手机支付、换钱、汇率、方便、银行、习惯}).
\item \textbf{Rédaction canevas (10 min) :} 10-12 répliques. Contraintes obligatoires :
    \begin{itemize}
    \item Touriste ouvre : \zh{我想换钱。人民币换非洲法郎。}
    \item Agent explique usage local : \zh{在喀麦隆，人们常用现金。}
    \item Journaliste pose \textbf{2 questions} avec \zh{为什么}.
    \item Sketch contient : (a) 1 \zh{比} ; (b) 1 \zh{我觉得} ; (c) 1 \zh{因为……所以……}.
    \end{itemize}
\item \textbf{Répétition / Présentation (5 min) :} Focus tons (\zh{rén-mín-bì}, \zh{fǎ-láng}). Présentation devant la classe. Public : grille d'écoute (vocabulaire entendu, faits exacts, structures repérées).
\end{enumerate}
\end{experience}

\begin{textebox}[Carnet du journaliste — Questions types pour l'interview]
1. 您觉得手机支付方便吗？为什么？\\
Nín juéde shǒujī zhīfù fāngbiàn ma? Wèishénme?\\
« Trouvez-vous le paiement mobile pratique ? Pourquoi ? »\\[0.4em]
2. 在中国，人们还用现金吗？\\
Zài Zhōngguó, rénmen hái yòng xiànjīn ma?\\
« En Chine, utilise-t-on encore des espèces ? »\\[0.4em]
3. 您觉得喀麦隆和中国的钱有什么不一样？\\
Nín juéde Kāmàilóng hé Zhōngguó de qián yǒu shénme bù yíyàng?\\
« Quelles différences entre l'argent camerounais et chinois ? »\\[0.4em]
4. 对游客来说，在喀麦隆用钱难不难？\\
Duì yóukè lái shuō, zài Kāmàilóng yòng qián nán bu nán?\\
« Pour un touriste, est-ce difficile d'utiliser l'argent au Cameroun ? »
\end{textebox}

\courssub{Production guidée : mini-exposé comparatif (Écrit)}

\begin{methode}[Rédiger un mini-exposé structuré (5 phrases)]
Suis ces 5 étapes pour couvrir toutes les fonctions langagières visées :
\begin{enumerate}
\item \textbf{Présenter les deux monnaies} (Vocabulaire précis : \zh{叫}, \zh{单位是}, \zh{是}).
\item \textbf{Comparer un usage} (Structure \zh{比} + adjectif factuel : \zh{普及}, \zh{多}, \zh{少}, \zh{方便}).
\item \textbf{Décrire une réalité locale} (Complément de lieu \zh{在}... avant verbe + fait socioculturel : marché, taxi, QR code).
\item \textbf{Exprimer ton avis nuancé} (\zh{我觉得} + \zh{不一样} / \zh{都有用} + \zh{但是}).
\item \textbf{Justifier / Conclure} (\zh{因为……所以……} + conseil \zh{应该} / \zh{建议}).
\end{enumerate}
\end{methode}

\begin{exoresolu}[Mini-exposé modèle — Production écrite (5 phrases)]
Rédige un texte de 5 phrases comparant la monnaie chinoise et la monnaie camerounaise. Contraintes : 1 phrase de présentation, 1 comparaison \zh{比}, 1 description d'usage \zh{在}..., 1 opinion \zh{我觉得}, 1 justification \zh{因为……所以……}.
\tcblower
\textbf{Production modèle :}\\[0.3em]
1. 中国的货币叫人民币，单位是元；喀麦隆用非洲法郎。\\
Zhōngguó de huòbì jiào rénmínbì, dānwèi shì yuán; Kāmàilóng yòng Fēizhōu fǎláng.\\
« La monnaie chinoise s'appelle le Renminbi, son unité est le yuan ; le Cameroun utilise le franc CFA. »\\[0.3em]
2. 在中国，手机支付比喀麦隆普及，用现金的人比较少。\\
Zài Zhōngguó, shǒujī zhīfù bǐ Kāmàilóng pǔjí, yòng xiànjīn de rén bǐjiào shǎo.\\
« En Chine, le paiement mobile est plus généralisé qu'au Cameroun, les gens utilisant des espèces sont relativement peu nombreux. »\\[0.3em]
3. 在喀麦隆，人们常用现金，特别是在市场和坐出租车时。\\
Zài Kāmàilóng, rénmen cháng yòng xiànjīn, tèbié shì zài shìchǎng hé zuò chūzūchē shí.\\
« Au Cameroun, on utilise souvent des espèces, surtout au marché et en prenant le taxi. »\\[0.3em]
4. 我觉得这两种支付习惯不一样，但是都适合自己的国家。\\
Wǒ juéde zhè liǎng zhǒng zhīfù xíguàn bù yíyàng, dànshì dōu shìhé zìjǐ de guójiā.\\
« Je trouve que ces deux habitudes de paiement sont différentes, mais chacune convient à son pays. »\\[0.3em]
5. 因为每个国家的经济环境不同，所以游客应该学会用当地的支付方式。\\
Yīnwèi měi ge guójiā de jīngjì huánjìng bù tóng, suǒyǐ yóukè yīnggāi xuéhuì yòng dāngdì de zhīfù fāngshì.\\
« Parce que l'environnement économique de chaque pays diffère, les touristes devraient apprendre à utiliser les modes de paiement locaux. »\\[0.4em]
\textbf{Commentaire linguistique détaillé :}\\
Phrase 1 : Vocabulaire technique exact (\zh{货币} huòbì, \zh{单位} dānwèi, \zh{叫} jiào). Point-virgule pour lier deux propositions indépendantes.\\
Phrase 2 : Double structure : \zh{比} (comparaison fait) + \zh{比较少} (degré). Respect de l'ordre \zh{在中国} (lieu) en tête.\\
Phrase 3 : \zh{在喀麦隆} (lieu) + \zh{特别是} (précision) + \zh{在...时} (temporel). Fait culturel ciblé (marché, taxi).\\
Phrase 4 : \zh{我觉得} (opinion) + \zh{不一样} (constat neutre) + \zh{但是} (concession) + \zh{适合} (évaluation fonctionnelle). Évite le jugement de valeur.\\
Phrase 5 : \zh{因为……所以……} (causalité complexe) + \zh{应该} (modalité déontique) + \zh{学会} (résultat). Conclusion actionnable.
\end{exoresolu}

\courssub{Bilan de la leçon et prolongement}

\begin{aretenir}[Bilan des acquis — Leçon 29]
\begin{itemize}
\item \textbf{Lexique :} Maîtrise du vocabulaire monétaire formel/oral (元/块, 角/毛, 分) et termes bancaires (汇率, 兑换, 人民币, 非洲法郎).
\item \textbf{Grammaire :} Comparatif \zh{比} (factuel), Opinion \zh{觉得}, Causalité \zh{因为……所以……}, Topicalisation lieu \zh{在}...V.
\item \textbf{Socioculturel :} Compréhension des écosystèmes de paiement contrastés (Cash-based vs Mobile-first) sans ethnocentrisme.
\item \textbf{Compétences :} Interaction orale (jeu de rôle 3 rôles), Production écrite courte structurée (5 phrases / 5 fonctions), Médiation (interview).
\end{itemize}
\end{aretenir}

\begin{savaistu}[Prolongement : L'avenir du paiement au Cameroun ?]
Le Cameroun connaît une croissance rapide du \zh{移动支付} (yídòng zhīfù, paiement mobile) via \zh{Mobile Money} (Orange Money, MTN MoMo, Express Union).
\begin{itemize}
\item Question débat : \zh{以后，喀麦隆人也会常用手机支付吗？} (Yǐhòu, Kāmàilóng rén yě huì cháng yòng shǒujī zhīfù ma ?)
\item Vocabulaire d'extension : \zh{手机钱包} (shǒujī qiánbāo, portefeuille mobile), \zh{扫码} (sǎomǎ, scanner un QR code), \zh{转账} (zhuǎnzhàng, virement), \zh{金融包容} (jīnróng bāoróng, inclusion financière).
\end{itemize}
Activité suggérée : Enquête de terrain (interview 3 commerçants à Yaoundé/Douala/Buea) → Présentation orale en chinois la semaine prochaine.
\end{savaistu}

\newpage

\coursec{Méthodes et exercices résolus}

\coursec{Leçon 29 — Éléments socioculturels : différence entre la monnaie chinoise et la monnaie camerounaise}

\courssub{Méthodes et exercices résolus}

\begin{methode}[Comparer deux systèmes monétaires : démarche en 4 étapes]
    \begin{enumerate}
        \item \textbf{Identifier les unités de base et les sous-unités.} Pour la Chine : \zh{元} (yuán), \zh{角} (jiǎo), \zh{分} (fēn) ; pour le Cameroun : \zh{法郎} (fǎláng) / FCFA. Noter l'absence de sous-unité courante pour le FCFA (le centime n'est plus utilisé).
        \item \textbf{Repérer le vocabulaire des transactions.} \zh{现金} (xiànjīn), \zh{支付} (zhīfù), \zh{扫码} (sǎomǎ), \zh{零钱} (língqián), \zh{汇率} (huìlǜ).
        \item \textbf{Structurer la comparaison.} Utiliser \zh{比较} (bǐjiào), \zh{不同} (bùtóng), \zh{相比之下} (xiāngbǐ zhīxià) ou la structure \zh{……，而……} (……, ér ……) pour opposer les faits.
        \item \textbf{Conclure sur l'usage quotidien.} Mentionner la prédominance du paiement mobile en Chine (\zh{微信支付}, \zh{支付宝}) vs l'usage mixte (espèces / Mobile Money / carte) au Cameroun.
    \end{enumerate}
\end{methode}

\begin{exoresolu}[Comparer les unités monétaires et les modes de paiement]
\textbf{Énoncé :} À partir du texte ci-dessous, complétez le tableau comparatif en chinois (caractères + pinyin) et répondez aux questions en phrases complètes.

\begin{textebox}[Texte support : Deux monnaies au quotidien]
在喀麦隆，法定货币是中非法郎（FCFA）。最小的单位是法郎，没有分。人们习惯用现金，也越来越多用手机转账（Mobile Money）。在中国，法定货币是人民币（RMB）。单位是元、角、分：1元=10角，1角=10分。现在中国人很少用现金，大多数人用微信支付或支付宝扫码付款，非常方便。
\par
\textbf{Pinyin :} Zài Kāmàilóng, fǎdìng huòbì shì Zhōngfēi Fǎláng (FCFA). Zuìxiǎo de dānwèi shì Fǎláng, méiyǒu fēn. Rénmen xíguàn yòng xiànjīn, yě yuè lái yuè duō yòng shǒujī zhuǎnzhàng (Mobile Money). Zài Zhōngguó, fǎdìng huòbì shì Rénmínbì (RMB). Dānwèi shì yuán, jiǎo, fēn: 1 yuán = 10 jiǎo, 1 jiǎo = 10 fēn. Xiànzài Zhōngguórén hěn shǎo yòng xiànjīn, dàduōshù rén yòng Wēixìn Zhīfù huò Zhīfùbǎo sǎomǎ fùkuǎn, fēicháng fāngbiàn.
\par
\textbf{Traduction :} Au Cameroun, la monnaie légale est le Franc CFA (FCFA). La plus petite unité est le franc, pas de centime. Les gens ont l'habitude d'utiliser du liquide, mais de plus en plus utilisent le transfert mobile (Mobile Money). En Chine, la monnaie légale est le Renminbi (RMB). Les unités sont yuan, jiao, fen : 1 yuan = 10 jiao, 1 jiao = 10 fen. Aujourd'hui, les Chinois utilisent peu le liquide, la plupart paient en scannant un code QR via WeChat Pay ou Alipay, c'est très pratique.
\end{textebox}

\textbf{Questions :}
\begin{enumerate}
    \item Complétez le tableau : Unités (Chine / Cameroun), Sous-unités, Paiement dominant.
    \item Traduisez en chinois : « Au Cameroun, on n'utilise pas de centimes. »
    \item Traduisez en chinois : « En Chine, le paiement mobile est plus pratique que l'argent liquide. »
    \item Exprimez une différence majeure avec la structure \zh{……，而……}.
\end{enumerate}

\tcblower
\textbf{Solution détaillée :}

\textbf{1. Tableau comparatif :}
\begin{center}
\begin{tabularx}{\linewidth}{|X|X|X|}
\hline
\rowcolor{popblueL} \textbf{Critère} & \textbf{Chine (中国)} & \textbf{Cameroun (喀麦隆)} \\ \hline
Unité principale & \zh{元} (yuán) & \zh{法郎} (fǎláng) / FCFA \\ \hline
Sous-unités & \zh{角} (jiǎo), \zh{分} (fēn) & \textit{Aucune en usage (pas de centime)} \\ \hline
Paiement dominant & \zh{扫码支付} (sǎomǎ zhīfù) / \zh{手机支付} (shǒujī zhīfù) & \zh{现金} (xiànjīn) / \zh{手机转账} (shǒujī zhuǎnzhàng) \\ \hline
\end{tabularx}
\end{center}

\textbf{2. Traduction : « Au Cameroun, on n'utilise pas de centimes. »}
\begin{itemize}
    \item \textbf{Analyse :} « On n'utilise pas » = \zh{不使用} (bù shǐyòng) ou \zh{没有} (méiyǒu) + objet. « Centimes » = \zh{分} (fēn) dans le contexte monétaire chinois, mais ici on parle de l'absence de sous-unité. On dit : « Il n'y a pas de \zh{分} ».
    \item \textbf{Réponse :} \zh{在喀麦隆，没有“分”。} (Zài Kāmàilóng, méiyǒu "fēn".)
    \item \textit{Alternative :} \zh{喀麦隆不使用分。} (Kāmàilóng bù shǐyòng fēn.)
\end{itemize}

\textbf{3. Traduction : « En Chine, le paiement mobile est plus pratique que l'argent liquide. »}
\begin{itemize}
    \item \textbf{Structure :} \zh{A 比 B + Adj.} (A est plus Adj que B).
    \item A = \zh{手机支付} (shǒujī zhīfù) / \zh{移动支付} (yídòng zhīfù). B = \zh{现金} (xiànjīn). Adj = \zh{方便} (fāngbiàn).
    \item \textbf{Réponse :} \zh{在中国，手机支付比现金方便。} (Zài Zhōngguó, shǒujī zhīfù bǐ xiànjīn fāngbiàn.)
    \item \textit{Note :} On peut ajouter \zh{得多} (de duō) pour « beaucoup plus » : \zh{手机支付比现金方便得多。}
\end{itemize}

\textbf{4. Différence majeure avec \zh{……，而……} (……, alors que …… / alors que ……) :}
\begin{itemize}
    \item \textbf{Structure :} Sujet 1 + Prédicat 1, \zh{而} (ér) + Sujet 2 + Prédicat 2. Marque une opposition forte.
    \item \textbf{Réponse :} \zh{中国人习惯扫码付款，而喀麦隆人还常用现金。} (Zhōngguórén xíguàn sǎomǎ fùkuǎn, ér Kāmàilóngrén hái chángyòng xiànjīn.) — Les Chinois ont l'habitude de payer en scannant, alors que les Camerounais utilisent encore souvent le liquide.
\end{itemize}

\textbf{Conclusion :} La maîtrise du vocabulaire monétaire (\zh{元/角/分} vs \zh{法郎}) et des structures de comparaison (\zh{比}, \zh{而}) est essentielle pour réussir l'expression écrite et orale sur ce thème.
\end{exoresolu}

\begin{methode}[Rédiger un mini-exposé comparatif (Expression écrite / Orale)]
    \begin{enumerate}
        \item \textbf{Introduction (1 phrase) :} Présenter le thème. \zh{今天我要比较中国和喀麦隆的货币差异。} (Jīntiān wǒ yào bǐjiào Zhōngguó hé Kāmàilóng de huòbì chāyì.)
        \item \textbf{Corps : 2 à 3 arguments structurés.} Utiliser les connecteurs logiques :
            \begin{itemize}
                \item \zh{首先……} (Shǒuxiān...) — Premièrement...
                \item \zh{其次……} (Qícì...) — Deuxièmement...
                \item \zh{另外……} (Lìngwài...) — En outre...
            \end{itemize}
            Chaque argument = Fait Chine + \zh{而/但是} (ér/dànshì) + Fait Cameroun.
        \item \textbf{Vocabulaire clé à insérer :} \zh{汇率} (huìlǜ), \zh{通货膨胀} (tōnghuò péngzhàng - optionnel niveau 4), \zh{无现金社会} (wú xiànjīn shèhuì), \zh{金融包容} (jīnróng bāoróng).
        \item \textbf{Conclusion (1 phrase) :} Avis personnel ou synthèse. \zh{总的来说，中国的支付方式更现代化。} (Zǒngde lái shuō, Zhōngguó de zhīfù fāngshì gèng xiàndàihuà.)
    \end{enumerate}
\end{methode}

\begin{exoresolu}[Rédiger un paragraphe comparatif structuré]
\textbf{Énoncé :} Rédigez un paragraphe d'environ 60-80 caractères chinois sur le thème : « Les différences de paiement entre la Chine et le Cameroun ». Utilisez au moins 3 connecteurs logiques (\zh{首先}, \zh{其次}, \zh{最后} ou \zh{另外}) et la structure \zh{比} une fois.

\tcblower
\textbf{Solution détaillée :}

\textbf{Étape 1 : Planification des idées.}
\begin{itemize}
    \item Idée 1 (Unités) : Chine a yuan/jiao/fen, Cameroun seulement franc.
    \item Idée 2 (Méthodes) : Chine = mobile dominant (WeChat/Alipay), Cameroun = espèces + Mobile Money.
    \item Idée 3 (Praticité) : Chine plus pratique (sans contact).
\end{itemize}

\textbf{Étape 2 : Rédaction en chinois (avec pinyin et traduction).}

\zh{首先，中国的货币单位有元、角、分，而喀麦隆只有法郎，没有分。}
(Shǒuxiān, Zhōngguó de huòbì dānwèi yǒu yuán, jiǎo, fēn, ér Kāmàilóng zhǐyǒu fǎláng, méiyǒu fēn.)
\textit{Premièrement, la monnaie chinoise a des unités yuan, jiao, fen, alors que le Cameroun n'a que le franc, pas de centime.}

\zh{其次，中国人大多用微信或支付宝扫码付款，非常快；但是喀麦隆人还常用现金和Mobile Money。}
(Qícì, Zhōngguórén dàduō yòng Wēixìn huò Zhīfùbǎo sǎomǎ fùkuǎn, fēicháng kuài; dànshì Kāmàilóngrén hái chángyòng xiànjīn hé Mobile Money.)
\textit{Deuxièmement, la plupart des Chinois paient en scannant avec WeChat ou Alipay, c'est très rapide ; mais les Camerounais utilisent encore souvent le liquide et le Mobile Money.}

\zh{最后，我觉得手机支付比用现金方便得多，中国已经接近无现金社会。}
(Zuìhòu, wǒ juéde shǒujī zhīfù bǐ yòng xiànjīn fāngbiàn de duō, Zhōngguó yǐjīng jiējìn wú xiànjīn shèhuì.)
\textit{Enfin, je trouve que le paiement mobile est bien plus pratique que l'espèce, la Chine est déjà proche d'une société sans numéraire.}

\textbf{Étape 3 : Vérification des contraintes.}
\begin{itemize}
    \item Connecteurs : \zh{首先}, \zh{其次}, \zh{最后} (\checkmark).
    \item Structure \zh{比} : \zh{手机支付比用现金方便得多} (\checkmark).
    \item Vocabulaire précis : \zh{扫码付款}, \zh{无现金社会} (\checkmark).
    \item Longueur : ~75 caractères (\checkmark).
\end{itemize}

\textbf{Version finale à présenter (caractères uniquement) :}
\zh{首先，中国的货币单位有元、角、分，而喀麦隆只有法郎，没有分。其次，中国人大多用微信或支付宝扫码付款，非常快；但是喀麦隆人还常用现金和Mobile Money。最后，我觉得手机支付比用现金方便得多，中国已经接近无现金社会。}
\end{exoresolu}

\begin{methode}[Comprendre un texte authentique : stratégie de lecture en 3 temps]
    \begin{enumerate}
        \item \textbf{Lecture globale (Skimming) :} Identifier le sujet (monnaie), le type de texte (comparatif informatif), la source (manuel / article). Repérer les noms propres : \zh{人民币}, \zh{中非法郎}, \zh{微信支付}, \zh{Mobile Money}.
        \item \textbf{Lecture détaillée (Scanning) :} Chercher les chiffres (taux de change approximatif 1 CNY $\approx$ 80-90 FCFA), les pourcentages (taux de pénétration du mobile payment), les marqueurs de contraste (\zh{但是}, \zh{然而}, \zh{相比之下}).
        \item \textbf{Inférence lexicale :} Deviner le sens des mots inconnus par le contexte ou la composition morphologique.
            \begin{itemize}
                \item \zh{法定货币} (fǎdìng huòbì) = \zh{法定} (légal) + \zh{货币} (monnaie) $\rightarrow$ Monnaie légale.
                \item \zh{零钱} (língqián) = \zh{零} (éparpillé/zéro) + \zh{钱} (argent) $\rightarrow$ Monnaie, petite monnaie.
                \item \zh{普及} (pǔjí) = \zh{普} (universel) + \zh{及} (atteindre) $\rightarrow$ Se généraliser, être populaire.
            \end{itemize}
    \end{enumerate}
\end{methode}

\begin{exoresolu}[Compréhension de texte et inférence lexicale]
\textbf{Énoncé :} Lisez le texte suivant et répondez aux questions.

\begin{textebox}[Article : L'évolution du paiement au Cameroun et en Chine]
随着中非贸易增长，了解对方货币很重要。中国官方货币是人民币（CNY），符号是¥。喀麦隆用中非法郎（XAF），固定汇率1欧元=655.957法郎。目前1元人民币约等于80至85法郎。
\par
在支付方式上，差异明显。中国已进入“无现金时代”：超市、出租车、甚至路边摊都支持扫码。数据显示，中国移动支付渗透率超85\%。相比之下，喀麦隆虽然推广Mobile Money（如Orange Money, MTN MoMo），现金仍占主导地位，尤其在农村市场买菜、坐摩托车（benskin）必须付现金。
\par
对于中国游客或商人在喀麦隆，建议携带美元或欧元现金兑换，因为直接兑换人民币网点较少。喀麦隆人去中国，建议办理银联卡或开通微信支付“钱包”功能绑定国际银行卡。
\par
\textbf{Pinyin :} Suízhe Zhōng-Fēi màoyì zēngzhǎng, liǎojiě duìfāng huòbì hěn zhòngyào. Zhōngguó guānfāng huòbì shì Rénmínbì (CNY), fúhào shì ¥. Kāmàilóng yòng Zhōngfēi Fǎláng (XAF), gùdìng huìlǜ 1 Ōuyuán = 655.957 Fǎláng. Mùqián 1 yuán Rénmínbì yuē děngyú 80 zhì 85 Fǎláng.
\par
Zài zhīfù fāngshì shàng, chāyì xiǎnxiàn. Zhōngguó yǐ jìnrù "wú xiànjīn shídài": chāoshì, chūzūchē, shènzhì lùbiāntān dōu zhīchí sǎomǎ. Shùjù xiǎnshì, Zhōngguó yídòng zhīfù shèntòulǜ chāo 85\%. Xiāngbǐ zhīxià, Kāmàilóng suīrán tuīguǎng Mobile Money (rú Orange Money, MTN MoMo), xiànjīn réng zhǔ zhǔdǎo dìwèi, yóuqí zài nóngcūn shìchǎng mǎicài, zuò mótuōchē (benskin) bìxū fù xiànjīn.
\par
Duìyú Zhōngguó yóukè huò shāngrén zài Kāmàilóng, jiànyí xiédài Měiyuán huò Ōuyuán xiànjīn duìhuàn, yīnwèi zhíjiē duìhuàn Rénmínbì wǎngdiǎn jiào shǎo. Kāmàilóngrén qù Zhōngguó, jiànyí bànlǐ Yínlián kǎ huò kāitōng Wēixìn Zhīfù "qiánbāo" gōngnéng bǎngdìng guójì yínhángkǎ.
\end{textebox}

\textbf{Questions :}
\begin{enumerate}
    \item Quel est le taux de change approximatif actuel 1 CNY = ? FCFA (d'après le texte) ?
    \item Citez deux marques de Mobile Money citées pour le Cameroun.
    \item Expliquez le sens de \zh{渗透率} (shèntòulǜ) dans ce contexte.
    \item Pourquoi le texte conseille-t-il aux Chinois d'apporter des USD/EUR au Cameroun ?
    \item Traduisez la phrase : \zh{喀麦隆人去中国，建议办理银联卡或开通微信支付“钱包”功能绑定国际银行卡。}
\end{enumerate}

\tcblower
\textbf{Solution détaillée :}

\textbf{1. Taux de change :} \zh{目前1元人民币约等于80至85法郎。} $\rightarrow$ \textbf{80 à 85 FCFA.}

\textbf{2. Marques Mobile Money :} \zh{Orange Money, MTN MoMo.} (Reconnaissance d'entités nommées dans le texte).

\textbf{3. Inférence \zh{渗透率} (shèntòulǜ) :}
\begin{itemize}
    \item Morphologie : \zh{渗透} (shèntòu = pénétrer / s'infiltrer) + \zh{率} (lǜ = taux).
    \item Contexte : « \zh{中国移动支付渗透率超85\%} » (Le taux de... du paiement mobile en Chine dépasse 85\%).
    \item Sens : \textbf{Taux de pénétration / Taux d'adoption / Part de marché.} Pourcentage de la population utilisant le service.
\end{itemize}

\textbf{4. Conseil USD/EUR :} \zh{因为直接兑换人民币网点较少。} (Parce que les points de change direct pour le RMB sont rares). Le FCFA est indexé sur l'Euro, le change EUR/FCFA est direct et fixe. Le change CNY/FCFA passe souvent par l'Euro ou le Dollar, d'où la recommandation.

\textbf{5. Traduction de la phrase :}
\begin{itemize}
    \item \zh{喀麦隆人去中国} (Kāmàilóngrén qù Zhōngguó) : Les Camerounais allant en Chine / Quand les Camerounais vont en Chine.
    \item \zh{建议} (jiànyí) : il est conseillé de / il faut.
    \item \zh{办理银联卡} (bànlǐ Yínlián kǎ) : faire une carte UnionPay (le réseau interbancaire chinois).
    \item \zh{或} (huò) : ou.
    \item \zh{开通微信支付“钱包”功能} (kāitōng Wēixìn Zhīfù "qiánbāo" gōngnéng) : activer la fonction "Portefeuille" de WeChat Pay.
    \item \zh{绑定国际银行卡} (bǎngdìng guójì yínhángkǎ) : lier/binder une carte bancaire internationale.
\end{itemize}
\textbf{Traduction finale :} « Pour les Camerounais qui se rendent en Chine, il est conseillé de se procurer une carte UnionPay ou d'activer la fonction "Portefeuille" de WeChat Pay en y liant une carte bancaire internationale. »
\end{exoresolu}

\begin{methode}[Exprimer un avis nuancé et débattre (Expression orale / Interaction)]
    \begin{enumerate}
        \item \textbf{Prendre position clairement.} \zh{我认为……} (Wǒ rènwéi...), \zh{我觉得……} (Wǒ juéde...), \zh{在我看来……} (Zài wǒ kànlái...).
        \item \textbf{Argumenter avec des connecteurs de cause/conséquence.}
            \begin{itemize}
                \item Cause : \zh{因为……所以……} (Yīnwèi... suǒyǐ...), \zh{由于……} (Yóuyú...).
                \item Conséquence : \zh{因此……} (Yīncǐ...), \zh{结果……} (Jiéguǒ...).
            \end{itemize}
        \item \textbf{Nuancer / Concéder.} \zh{虽然……，但是……} (Suīrán..., dànshì...), \zh{不过……} (Bùguò...), \zh{一方面……另一方面……} (Yīfāngmiàn... lìngyīfāngmiàn...).
        \item \textbf{Inviter l'interlocuteur.} \zh{你怎么看？} (Nǐ zěnme kàn?), \zh{你同意吗？} (Nǐ tóngyì ma?).
    \end{enumerate}
\end{methode}

\begin{exoresolu}[Débat : "Le Cameroun deviendra-t-il une société sans numéraire comme la Chine ?"]
\textbf{Énoncé :} En binôme, préparez un mini-débat (2 minutes chacun). L'élève A est \textbf{POUR} (Oui, c'est possible), l'élève B est \textbf{CONTRE} (Non, des obstacles persistent). Utilisez au moins 3 arguments chacun avec les structures ci-dessus. Écrivez vos notes en caractères + pinyin.

\tcblower
\textbf{Solution détaillée (Modèles de réponses types niveau HSK 3-4) :}

\textbf{Élève A (POUR) :}
\zh{我认为喀麦隆未来会变成无现金社会。}
(Wǒ rènwéi Kāmàilóng wèilái huì biànchéng wú xiànjīn shèhuì.)
\textit{Je pense que le Cameroun deviendra une société sans numéraire à l'avenir.}

\zh{首先，因为手机普及率很高，即使在农村也有网络。}
(Shǒuxiān, yīnwèi shǒujī pǔjí lǜ hěn gāo, jíshǐ zài nóngcūn yě yǒu wǎngluò.)
\textit{Premièrement, parce que le taux de pénétration des téléphones est élevé, même en zone rurale il y a internet.}

\zh{其次，Mobile Money像Orange Money已经很流行，年轻人习惯用手机转账。}
(Qícì, Mobile Money xiàng Orange Money yǐjīng hěn liúxíng, niánqīngrén xíguàn yòng shǒujī zhuǎnzhàng.)
\textit{Deuxièmement, le Mobile Money comme Orange Money est déjà très populaire, les jeunes ont l'habitude de transférer par téléphone.}

\zh{虽然现在现金还是主流，但是政府在推广金融数字化，所以我相信会改变。}
(Suīrán xiànzài xiànjīn háishì zhǔliú, dànshì zhèngfǔ zài tuīguǎng jīnróng shùzìhuà, suǒyǐ wǒ xiāngxìn huì gǎibiàn.)
\textit{Bien que le liquide reste dominant maintenant, le gouvernement promeut la finance numérique, donc je crois que cela changera.}

\textbf{Élève B (CONTRE) :}
\zh{我不太同意。我觉得喀麦隆很难完全无现金。}
(Wǒ bù tài tóngyì. Wǒ juéde Kāmàilóng hěn nán wánquán wú xiànjīn.)
\textit{Je ne suis pas tout à fait d'accord. Je trouve que le Cameroun aura du mal à être totalement sans numéraire.}

\zh{因为电力不稳定，网络在偏远地区很差，手机没电就不能付款。}
(Yīnwèi diànlì bù wěndìng, wǎngluò zài piānyuǎn dìqū hěn chā, shǒujī méidiàn jiù bù néng fùkuǎn.)
\textit{Parce que l'électricité est instable, le réseau est mauvais dans les zones reculées, téléphone déchargé = impossible de payer.}

\zh{另外，很多人没有银行账户，不信任数字货币，更喜欢摸到真钱。}
(Lìngwài, hěn duō rén méiyǒu yínháng zhànghù, bù xìnrèn shùzì huòbì, gèng xǐhuan mō dào zhēn qián.)
\textit{De plus, beaucoup n'ont pas de compte bancaire, ne font pas confiance à la monnaie numérique, préfèrent toucher de l'argent réel.}

\zh{不过，我在大城市如雅温得、杜阿拉会看到更多扫码支付。}
(Bùguò, wǒ zài dà chéngshì rú Yǎwēndé, Dù'ālā huì kàn dào gèng duō sǎomǎ zhīfù.)
\textit{Cependant, dans les grandes villes comme Yaoundé, Douala, on verra plus de paiement par QR code.}

\textbf{Analyse grammaticale des modèles :}
\begin{itemize}
    \item Structure \zh{因为……所以……} bien maîtrisée.
    \item Concession \zh{虽然……但是……} / \zh{不过} utilisée correctement.
    \item Vocabulaire spécifique : \zh{普及率} (taux de pénétration), \zh{偏远地区} (zones reculées), \zh{银行账户} (compte bancaire), \zh{金融数字化} (numérisation financière).
    \item Tons : Attention à \zh{普及} (pǔjí, 3-2), \zh{偏远} (piānyuǎn, 1-3), \zh{数字化} (shùzìhuà, 4-4-4).
\end{itemize}
\end{exoresolu}

\begin{methode}[Traduire des termes financiers et administratifs (Thème / Version)]
    \begin{enumerate}
        \item \textbf{Identifier le registre :} Administratif / Bancaire $\rightarrow$ Vocabulaire formel (\zh{书面语} shūmiànyǔ). Ex: \zh{兑换} (duìhuàn) au lieu de \zh{换} (huàn), \zh{办理} (bànlǐ) au lieu de \zh{做} (zuò).
        \item \textbf{Maîtriser les paires fixes (Collocations) :}
            \begin{center}
            \begin{tabularx}{\linewidth}{|X|X|X|}
            \hline
            \rowcolor{popblueL} \textbf{Français} & \textbf{Chinois (Verbe + Objet)} & \textbf{Pinyin} \\ \hline
            Changer de l'argent & \zh{兑换货币 / 兑换现金} & duìhuàn huòbì / duìhuàn xiànjīn \\ \hline
            Taux de change & \zh{汇率} & huìlǜ \\ \hline
            Ouvrir un compte & \zh{开户 / 开银行卡} & kāihù / kāi yínhángkǎ \\ \hline
            Retirer de l'argent & \zh{取款 / 提现} & qǔkuǎn / tíxiàn \\ \hline
            Transférer / Virer & \zh{转账} & zhuǎnzhàng \\ \hline
            Scanner (QR code) & \zh{扫码 / 扫一扫} & sǎomǎ / sǎo yī sǎo \\ \hline
            \end{tabularx}
            \end{center}
        \item \textbf{Attention aux "faux amis" numériques :}
            \zh{亿} (yì) = 100 millions (10 puissance 8) $\neq$ Billion (FR) / Milliard (US).
            \zh{兆} (zhào) = 10 puissance 12 (Trillion).
            En chinois, on groupe par 4 zéros (万 wàn, 亿 yì), en français par 3 (mille, million, milliard).
        \item \textbf{Structure de la phrase bancaire :} Sujet (Banque/Client) + \zh{为} (wèi) / \zh{向} (xiàng) + Bénéficiaire + \zh{办理/提供} + Service.
    \end{enumerate}
\end{methode}

\begin{exoresolu}[Traduction spécialisée : Vocabulaire bancaire et chiffres]
\textbf{Énoncé :}
\begin{enumerate}
    \item Traduisez en chinois :
        a) Le taux de change actuel est d'environ 1 Yuan pour 82 Francs CFA.
        b) Je dois aller à la banque pour ouvrir un compte et faire une carte UnionPay.
        c) Ce marchand n'accepte pas le liquide, il faut payer en scannant le code.
    \item Traduisez en français :
        a) \zh{请问，汇率是多少？我想兑换1000元人民币。}
        b) \zh{您的余额不足，无法完成转账，请先充值。}
        c) \zh{中国的GDP已经超过120万亿元人民币。} (Note : \zh{万亿} = 10 000 milliards = 10 puissance 12).
\end{enumerate}

\tcblower
\textbf{Solution détaillée :}

\textbf{1. Version (Français $\rightarrow$ Chinois) :}

\textbf{a) Le taux de change actuel est d'environ 1 Yuan pour 82 Francs CFA.}
\begin{itemize}
    \item Structure : Sujet (\zh{当前汇率}) + \zh{是} (shì) + Valeur.
    \item « Environ » = \zh{约} (yuē) / \zh{大概} (dàgài).
    \item « Pour » (équivalence) = \zh{兑换} (duìhuàn) ou structure \zh{1元 = 82法郎}.
    \item \textbf{Réponse :} \zh{当前汇率约为1元人民币兑换82法郎。} (Dāngqián huìlǜ yuē wéi 1 yuán Rénmínbì duìhuàn 82 Fǎláng.)
    \item \textit{Variante plus orale :} \zh{现在汇率大概是1块钱换82法郎。} (\zh{块} kuài = \zh{元} oral).
\end{itemize}

\textbf{b) Je dois aller à la banque pour ouvrir un compte et faire une carte UnionPay.}
\begin{itemize}
    \item « Devoir » = \zh{得} (dei - oral) / \zh{必须} (bìxū) / \zh{要} (yào).
    \item « Ouvrir un compte » = \zh{开户} (kāihù).
    \item « Faire une carte » = \zh{办卡} (bàn kǎ) / \zh{申请银行卡} (shēnqǐng yínhángkǎ).
    \item \textbf{Réponse :} \zh{我得去银行开户，再办一张银联卡。} (Wǒ dei qù yínháng kāihù, zài bàn yī zhāng Yínlián kǎ.)
    \item \textit{Note :} Classificateur \zh{张} (zhāng) pour les cartes.
\end{itemize}

\textbf{c) Ce marchand n'accepte pas le liquide, il faut payer en scannant le code.}
\begin{itemize}
    \item « Marchand » = \zh{商家} (shāngjiā) / \zh{老板} (lǎobǎn).
    \item « Accepter » = \zh{接受} (jiēshòu) / \zh{收} (shōu).
    \item « Il faut / Falloir » = \zh{得} (dei) / \zh{必须} (bìxū) / \zh{只能} (zhǐ néng - "ne peut que").
    \item \textbf{Réponse :} \zh{这家商家不收现金，只能扫码付款。} (Zhè jiā shāngjiā bù shōu xiànjīn, zhǐ néng sǎomǎ fùkuǎn.)
    \item \textit{Variante :} \zh{这个老板不接受现金，必须扫码支付。}
\end{itemize}

\textbf{2. Thème (Chinois $\rightarrow$ Français) :}

\textbf{a) \zh{请问，汇率是多少？我想兑换1000元人民币。}}
\begin{itemize}
    \item \zh{请问} (qǐngwèn) = Excusez-moi / Je voudrais savoir.
    \item \zh{汇率是多少} (huìlǜ shì duōshao) = Quel est le taux de change ?
    \item \zh{兑换} (duìhuàn) = Échanger / Convertir.
    \item \textbf{Traduction :} « Excusez-moi, quel est le taux de change ? Je voudrais échanger 1000 yuans RMB. »
\end{itemize}

\textbf{b) \zh{您的余额不足，无法完成转账，请先充值。}}
\begin{itemize}
    \item \zh{您的} (nín de) = Votre (formel).
    \item \zh{余额} (yúé) = Solde (compte).
    \item \zh{不足} (bùzú) = Insuffisant.
    \item \zh{无法} (wúfǎ) = Ne pas pouvoir (formel, écrit) = \zh{不能} (bù néng).
    \item \zh{完成} (wánchéng) = Effectuer / Terminer.
    \item \zh{充值} (chōngzhí) = Recharger / Créditer.
    \item \textbf{Traduction :} « Votre solde est insuffisant, impossible d'effectuer le virement, veuillez d'abord recharger (votre compte). »
\end{itemize}

\textbf{c) \zh{中国的GDP已经超过120万亿元人民币。}}
\begin{itemize}
    \item \zh{GDP} = PIB (Produit Intérieur Brut).
    \item \zh{已经} (yǐjīng) = Déjà.
    \item \zh{超过} (chāoguò) = Dépasser.
    \item \zh{120万亿} (120 wàn yì) = 120 $\times$ 10 000 milliards = 1 200 000 milliards = 120 000 milliards (échelle longue).
    \item \textbf{Traduction :} « Le PIB de la Chine a déjà dépassé 120 000 milliards de yuans (RMB). »
    \item \textbf{Piège chiffres :} Ne pas traduire \zh{万亿} par "120 milliards" (erreur majeure de facteur 1000). \zh{万} (wàn) = 10 000. \zh{亿} (yì) = 100 000 000. \zh{万亿} = 10 000 $\times$ 100 000 000 = 1 000 000 000 000 (10 puissance 12).
\end{itemize}
\end{exoresolu}

\begin{methode}[Écrire les caractères clés de la leçon : Ordre des traits et radicaux]
    \begin{enumerate}
        \item \zh{币} (bì) — Monnaie. \textbf{Radical :} \zh{巾} (jīn, tissu/étoffe) à gauche. \textbf{Traits :} 18. Ordre : partie haute \zh{罒} (wǎng, filet, 5 traits), puis \zh{巾} (jīn, 3 traits), puis \zh{弋} (yì, 3 traits) pour la partie phonétique simplifiée \zh{弊}. \textit{Astuce :} Monnaie en tissu (étoffe) autrefois.
        \item \zh{兑} (duì) — Échanger. \textbf{Radical :} \zh{儿} (ér, jambes/humain) en bas. \textbf{Traits :} 7. \zh{八} (huit) + \zh{儿} + trait horizontal final. \textit{Astuce :} Deux mains (\zh{八}) qui échangent quelque chose entre humains.
        \item \zh{汇} (huì) — Change / Rassembler. \textbf{Radical :} \zh{氵} (shuǐ, eau) à gauche. \textbf{Traits :} 13. Eau qui converge. Partie droite : \zh{彗} (huì, comète/balai) phonétique.
        \item \zh{率} (lǜ) — Taux / Ratio. \textbf{Radical :} \zh{玄} (xuán) en haut. \textbf{Traits :} 11. Complexe. Mémoriser : \zh{帅} (shuài, général) + \zh{玄} $\rightarrow$ Le général commande le taux.
        \item \zh{付} (fù) — Payer. \textbf{Radical :} \zh{亻} (rén, humain) à gauche. \textbf{Traits :} 5. Simple. Humain + \zh{寸} (cùn, pouce/mesure) $\rightarrow$ Mesurer pour payer.
        \item \zh{款} (kuǎn) — Fonds / Article / Paiement. \textbf{Radical :} \zh{欠} (qiàn, devoir) en bas à droite. \textbf{Traits :} 12. \zh{款} = partie haute (\zh{⺁} herbe + \zh{日} jour) + \zh{欠} (dette) $\rightarrow$ Régler sa dette jour après jour.
    \end{enumerate}
\end{methode}

\begin{exoresolu}[Dictée de caractères et reconnaissance visuelle]
\textbf{Énoncé :}
\begin{enumerate}
    \item Écrivez les caractères correspondant au pinyin suivant (attention aux tons) :
        a) huìlǜ \hspace{2cm} b) duìhuàn \hspace{2cm} c) fùkuǎn \hspace{2cm} d) xiànjīn \hspace{2cm} e) língqián
    \item Complétez le caractère manquant (radical / composant donné) :
        a) \zh{氵} + \zh{彗} = \_\_\_\_\_ (huì)
        b) \zh{亻} + \zh{寸} = \_\_\_\_\_ (fù)
        c) Radical \zh{巾} + composant phonétique \zh{弊} (simplifié : \zh{罒} + \zh{巾} + \zh{弋}) = \_\_\_\_\_ (bì)
    \item Choisissez le bon caractère / mot pour compléter la phrase :
        \zh{我没有\_\_\_\_\_，不能坐公交车。} (a) \zh{付} (b) \zh{币} (c) \zh{零钱} (d) \zh{率}
\end{enumerate}

\tcblower
\textbf{Solution détaillée :}

\textbf{1. Dictée (Caractères + Pinyin + Traduction) :}
\begin{center}
\begin{tabularx}{\linewidth}{|X|X|X|X|}
\hline
\rowcolor{popblueL} \textbf{Pinyin} & \textbf{Caractères} & \textbf{Nature} & \textbf{Traduction} \\ \hline
a) huìlǜ & \zh{汇率} & Nom & Taux de change \\ \hline
b) duìhuàn & \zh{兑换} & Verbe & Échanger (devises) \\ \hline
c) fùkuǎn & \zh{付款} & Verbe/Nom & Payer / Paiement \\ \hline
d) xiànjīn & \zh{现金} & Nom & Espèces / Liquide \\ \hline
e) língqián & \zh{零钱} & Nom & Monnaie / Petite monnaie \\ \hline
\end{tabularx}
\end{center}

\textbf{2. Composition par radicaux / composants :}
\begin{itemize}
    \item a) \zh{氵} + \zh{彗} = \zh{汇} (huì). \textit{Rappel :} Radicale eau (liquidité financière).
    \item b) \zh{亻} + \zh{寸} = \zh{付} (fù). \textit{Rappel :} Action humaine de mesurer/comptabiliser.
    \item c) \zh{币} (bì). Structure gauche-droite : radical \zh{巾} (jīn) + phonétique \zh{弊} (bì, simplifiée). \textit{Ordre traits :} \zh{巾} (3 traits) puis \zh{弊} (haut \zh{罒} 5 traits, milieu \zh{巾} 3 traits, bas \zh{弋} 3 traits). Total 14 traits (standard) ou 18 selon classification. Écrire lentement.
\end{itemize}

\textbf{3. Choix du bon mot (Contexte syntaxique) :}
\zh{我没有\_\_\_\_\_，不能坐公交车。} (Wǒ méiyǒu \_\_\_\_\_, bù néng zuò gōngjiāochē.)
\begin{itemize}
    \item (a) \zh{付} (fù) : Verbe "payer". \textit{Incorrect :} Il manque un objet ou structure \zh{付钱}.
    \item (b) \zh{币} (bì) : Nom "monnaie/unité". \textit{Incorrect :} Trop abstrait, on ne dit pas "je n'ai pas monnaie" avec \zh{币} seul.
    \item (c) \zh{零钱} (língqián) : Nom "petite monnaie". \textbf{Correct.} "Je n'ai pas de monnaie (pour l'appoint), je ne peux pas prendre le bus." Contexte réel : les bus n'ont souvent pas de rendu de monnaie sur gros billets.
    \item (d) \zh{率} (lǜ) : Nom "taux". \textit{Incorrect :} Hors sujet.
\end{itemize}
\textbf{Bonne réponse : (c) \zh{零钱}.}
\end{exoresolu}

\begin{attention}[Les 5 erreurs qui coûtent des points au Bac — Spécial Monnaie \& Paiement]
    \begin{enumerate}
        \item \textbf{Confusion \zh{元} / \zh{块} / \zh{角} / \zh{毛} / \zh{分} et ordre des mots.}
        \begin{itemize}
            \item \zh{元} (yuán) = formel (écrit, bancaire). \zh{块} (kuài) = oral (courant). \textbf{Ne pas mélanger dans une même phrase écrite formelle.}
            \item \zh{角} (jiǎo) = formel. \zh{毛} (máo) = oral.
            \item \textbf{Ordre :} \zh{12元5角} (12.50) ou \zh{12块5毛}. \textbf{Jamais} \zh{5角12元}.
            \item \textbf{Piège :} \zh{100分 = 1元}. Ne pas écrire \zh{100分 = 1块} (mélange registres).
        \end{itemize}
        \item \textbf{Mauvaise structure de comparaison : \zh{比} (bǐ) vs \zh{更} (gèng) vs \zh{比较} (bǐjiào).}
        \begin{itemize}
            \item \zh{中国比喀麦隆现代化。} (\checkmark Correct : A \zh{比} B Adj).
            \item \zh{中国更现代化比喀麦隆。} (\textbf{FAUX} : \zh{更} ne prend pas \zh{比} derrière, et l'ordre est inversé).
            \item \zh{中国比较现代化。} (\checkmark "Relativement moderne", sans comparaison explicite B).
            \item \textbf{Règle :} \zh{A 比 B Adj} (Adj \textbf{sans} \zh{很}, \zh{更}, \zh{最}). Si \zh{更/最} $\rightarrow$ Pas de structure \zh{比}.
        \end{itemize}
        \item \textbf{Place du complément de lieu / mode \textbf{AVANT} le verbe (Règle d'or ordre des mots).}
        \begin{itemize}
            \item \textbf{FAUX :} \zh{我付款用微信。} (Verbe + Objet + Mode).
            \item \textbf{JUSTE :} \zh{我用微信付款。} (Sujet + \zh{用} + Outil + Verbe).
            \item \textbf{JUSTE :} \zh{我在喀麦隆用法郎。} (Lieu avant verbe).
            \item \textbf{FAUX :} \zh{我用法郎在喀麦隆。} (Lieu après verbe = focalisation sur le lieu, rare ici).
        \end{itemize}
        \item \textbf{Classificateurs omis ou erronés pour l'argent / objets.}
        \begin{itemize}
            \item Billets : \zh{张} (zhāng) — \zh{一张百元钞票} (yī zhāng bǎi yuán chāopiào).
            \item Pièces / Monnaie : \zh{个} (gè) ou \zh{枚} (méi - formel) — \zh{一枚硬币} (yī méi yìngbì).
            \item Cartes (bancaires, téléphone) : \zh{张} (zhāng) — \zh{一张银行卡} (yī zhāng yínhángkǎ).
            \item \textbf{Erreur fatale :} \zh{一*个*钞票}, \zh{两*个*银行卡}. (Oubli du classificateur spécifique = 0 point en rédaction).
        \end{itemize}
        \item \textbf{Calque du français sur "Argent / Monnaie / Change".}
        \begin{itemize}
            \item « Changer de l'argent » $\neq$ \zh{换钱} (huàn qián - trop vague, oral).
            \item \textbf{Termes précis :} \zh{兑换} (duìhuàn - changer devise contre devise), \zh{兑换外汇} (duìhuàn wàihuì - changer devises), \zh{找零} (zhǎo líng - rendre la monnaie).
            \item « Monnaie » (pièces) = \zh{硬币} (yìngbì) ou \zh{零钱} (língqián).
            \item « Monnaie » (système) = \zh{货币} (huòbì) ou \zh{币制} (bìzhì).
            \item « Taux de change » = \zh{汇率} (huìlǜ), \textbf{pas} \zh{换率} (huàn lǜ - inexistant).
        \end{itemize}
    \end{enumerate}
\end{attention}

\newpage

\coursec{Exercices}

\courssub{Je vérifie mes connaissances}

\begin{exercice}[QCM : les deux monnaies]
Entoure la bonne réponse.
\begin{enumerate}
\item La monnaie officielle de la Chine s'appelle :
\begin{enumerate}
\item le franc CFA \item le renminbi (人民币) \item le dollar
\end{enumerate}
\item En Chine, la banque qui émet la monnaie est :
\begin{enumerate}
\item 中国人民银行 \item 中国银行 \item 非洲银行
\end{enumerate}
\item Au Cameroun, la monnaie utilisée est :
\begin{enumerate}
\item 人民币 \item 美元 \item 非洲法郎
\end{enumerate}
\item \zh{手机支付} (shǒujī zhīfù) signifie :
\begin{enumerate}
\item payer en espèces \item payer par téléphone portable \item payer par carte bancaire
\end{enumerate}
\end{enumerate}
\end{exercice}

\begin{exercice}[Vrai ou faux, à justifier]
Dis si chaque phrase est vraie ou fausse. Justifie ta réponse en français, puis corrige les phrases fausses en chinois.
\begin{enumerate}
\item Le franc CFA est utilisé dans plusieurs pays d'Afrique centrale.
\item En Chine, on paie toujours en espèces.
\item Le renminbi est la monnaie de la Chine.
\item Au Cameroun, tout le monde paie avec le téléphone portable.
\end{enumerate}
\end{exercice}

\begin{exercice}[Vocabulaire : traduis en français]
Traduis les mots suivants en français, puis écris-les trois fois en respectant l'ordre des traits.
\begin{enumerate}
\item \zh{钱} \item \zh{人民币} \item \zh{非洲法郎} \item \zh{银行} \item \zh{现金} \item \zh{手机支付}
\end{enumerate}
\end{exercice}

\begin{exercice}[Associe caractères et pinyin]
Relie chaque mot en caractères à son pinyin.
\begin{center}
\begin{tabularx}{\linewidth}{|X|X|}
\hline
\rowcolor{popblueL} Caractères & Pinyin \\
\hline
1. 人民币 & a. yínháng \\
\hline
2. 银行 & b. xiànjīn \\
\hline
3. 现金 & c. rénmínbì \\
\hline
4. 手机支付 & d. Fēizhōu fǎláng \\
\hline
5. 非洲法郎 & e. shǒujī zhīfù \\
\hline
\end{tabularx}
\end{center}
\end{exercice}

\courssub{Je m'entraîne}

\begin{exercice}[Texte à trous]
Complète le paragraphe comparatif avec les mots de la banque suivante : \zh{钱}、\zh{人民币}、\zh{非洲法郎}、\zh{银行}、\zh{现金}、\zh{手机支付}。

\begin{textebox}[Paragraphe à compléter]
在中国，人们用\_\_\_\_\_\_，这是中国的人民币。在喀麦隆，人们用\_\_\_\_\_\_。中国的\_\_\_\_\_\_发行人民币。在中国，很多人用\_\_\_\_\_\_，不太用\_\_\_\_\_\_。在喀麦隆，人们常常用现金买\_\_\_\_\_\_。
\end{textebox}
\end{exercice}

\begin{exercice}[Traduction : thème]
Traduis les phrases suivantes en chinois (caractères obligatoires).
\begin{enumerate}
\item Au Cameroun, on utilise le franc CFA.
\item En Chine, beaucoup de gens paient avec leur téléphone.
\item La Banque populaire de Chine émet le renminbi.
\item Je pense que le paiement par téléphone est très pratique.
\end{enumerate}
\end{exercice}

\begin{exercice}[Traduction : version]
Traduis le texte suivant en français.

\begin{textebox}[Texte à traduire]
中国人民银行发行人民币。在喀麦隆，人们常常用现金。\\
Zhōngguó Rénmín Yínháng fāxíng rénmínbì. Zài Kāmàilóng, rénmen chángcháng yòng xiànjīn.
\end{textebox}
\end{exercice}

\begin{exercice}[Remise en ordre]
Remets les mots dans l'ordre pour former des phrases correctes, puis traduis-les en français.
\begin{enumerate}
\item 用 / 人们 / 手机支付 / 在中国 / 常常
\item 是 / 非洲法郎 / 的 / 喀麦隆 / 钱
\item 觉得 / 我 / 现金 / 很方便 / 不
\end{enumerate}
\end{exercice}

\courssub{Je me prépare au Bac}

\begin{exercice}[Compréhension de texte type Bac]
Lis le texte suivant, puis réponds aux questions \textbf{en chinois}, par des phrases complètes.

\begin{textebox}[Deux monnaies, deux habitudes]
中国的钱叫人民币，中国人民银行发行人民币。现在在中国，很多人不用现金，他们用手机支付，买东西很方便。在喀麦隆，人们用非洲法郎。很多喀麦隆人常常用现金买东西，但是也有一些人开始用手机支付。我觉得两个国家的付钱方法不一样，但是都很有意思。\\
Zhōngguó de qián jiào rénmínbì, Zhōngguó Rénmín Yínháng fāxíng rénmínbì. Xiànzài zài Zhōngguó, hěn duō rén bú yòng xiànjīn, tāmen yòng shǒujī zhīfù, mǎi dōngxi hěn fāngbiàn. Zài Kāmàilóng, rénmen yòng Fēizhōu fǎláng. Hěn duō Kāmàilóng rén chángcháng yòng xiànjīn mǎi dōngxi, dànshì yě yǒu yìxiē rén kāishǐ yòng shǒujī zhīfù. Wǒ juéde liǎng ge guójiā de fù qián fāngfǎ bù yíyàng, dànshì dōu hěn yǒu yìsi.
\end{textebox}

Questions :
\begin{enumerate}
\item 中国的钱叫什么名字？
\item 谁发行人民币？
\item 在中国，现在很多人怎么付钱？
\item 在喀麦隆，人们常常用什么买东西？
\item 作者觉得两个国家的付钱方法怎么样？
\end{enumerate}

Question de langue : relève dans le texte deux mots du champ lexical de la monnaie et traduis-les en français.
\end{exercice}

\begin{exercice}[Thème et version]
\textbf{A. Thème.} Traduis en chinois (caractères obligatoires) :
\begin{enumerate}
\item Au Cameroun, on utilise le franc CFA. En Chine, beaucoup de gens paient avec leur téléphone.
\item Je pense que payer en espèces n'est pas très pratique, parce que c'est lent.
\end{enumerate}

\textbf{B. Version.} Traduis en français :

\begin{textebox}[Texte à traduire]
人民币是中国的钱。在中国，手机支付很方便。在喀麦隆，人们常常用现金，但是手机支付也越来越多。\\
Rénmínbì shì Zhōngguó de qián. Zài Zhōngguó, shǒujī zhīfù hěn fāngbiàn. Zài Kāmàilóng, rénmen chángcháng yòng xiànjīn, dànshì shǒujī zhīfù yě yuè lái yuè duō.
\end{textebox}
\end{exercice}

\begin{exercice}[Expression écrite et orale type Bac]
\textbf{A. Expression écrite.} Rédige en chinois un court texte de \textbf{5 à 8 phrases} (au moins 60 caractères) comparant la monnaie camerounaise et la monnaie chinoise. Ton texte doit contenir :
\begin{itemize}
\item le nom des deux monnaies ;
\item une différence d'habitude de paiement (现金 / 手机支付) ;
\item au moins une opinion personnelle avec la structure \zh{我觉得……因为……} (wǒ juéde... yīnwèi...).
\end{itemize}

\textbf{B. Expression orale.} Prépare un mini-exposé de 2 minutes intitulé « La monnaie de mon pays et la monnaie chinoise ». Appuie-toi sur un tableau comparatif (monnaie, banque émettrice, moyen de paiement habituel, avantages) que tu présenteras à la classe. Tu commenceras par \zh{大家好！今天我谈谈……} (Dàjiā hǎo! Jīntiān wǒ tántan...) et tu concluras par ton opinion personnelle.
\end{exercice}

\newpage

\coursec{Corrigés des exercices}

\coursec{Corrigés des exercices — Leçon 29 : Monnaies et paiements Cameroun / Chine}

\begin{corrige}[Exercice 1 — QCM : les deux monnaies]
\begin{enumerate}
\item \textbf{b.} le renminbi (\zh{人民币}, rénmínbì). C'est la monnaie officielle de la République populaire de Chine ; son unité de base est le \zh{元} (yuán).
\item \textbf{a.} \zh{中国人民银行} (Zhōngguó Rénmín Yínháng), la Banque populaire de Chine, banque centrale qui émet la monnaie. Attention : \zh{中国银行} (Zhōngguó Yínháng, Bank of China) est une banque commerciale, pas la banque centrale.
\item \textbf{c.} \zh{非洲法郎} (Fēizhōu fǎláng), le franc CFA, monnaie utilisée au Cameroun (zone CEMAC).
\item \textbf{b.} payer par téléphone portable : \zh{手机} (shǒujī, téléphone portable) + \zh{支付} (zhīfù, payer).
\end{enumerate}
\end{corrige}

\begin{corrige}[Exercice 2 — Vrai ou faux, à justifier]
\begin{enumerate}
\item \textbf{Vrai.} Le franc CFA (XAF) est la monnaie commune de six pays d'Afrique centrale (Cameroun, Tchad, Gabon, Congo, Guinée équatoriale, République centrafricaine) au sein de la CEMAC. Il existe aussi un franc CFA (XOF) en Afrique de l'Ouest (UEMOA), distinct mais à parité fixe.
\item \textbf{Faux.} En Chine, le paiement par téléphone portable (\zh{手机支付}) est très répandu ; une grande partie de la population n'utilise plus d'espèces au quotidien. Correction : \zh{在中国，很多人不用现金，他们用手机支付。} \emph{Zài Zhōngguó, hěn duō rén bú yòng xiànjīn, tāmen yòng shǒujī zhīfù.} « En Chine, beaucoup de gens n'utilisent pas d'espèces, ils paient par téléphone. »
\item \textbf{Vrai.} \zh{人民币} (rénmínbì) signifie littéralement « la monnaie du peuple » ; c'est bien la monnaie de la Chine.
\item \textbf{Faux.} Au Cameroun, le paiement en espèces (\zh{现金}) reste très courant, même si le paiement mobile (\zh{手机支付}, ex. Mobile Money) se développe rapidement. Correction : \zh{在喀麦隆，人们常常用现金，但是手机支付也越来越多。} \emph{Zài Kāmàilóng, rénmen chángcháng yòng xiànjīn, dànshì shǒujī zhīfù yě yuè lái yuè duō.} « Au Cameroun, les gens utilisent souvent des espèces, mais le paiement par téléphone est de plus en plus fréquent. »
\end{enumerate}
\end{corrige}

\begin{corrige}[Exercice 3 — Vocabulaire : traduis en français]
\begin{center}
\begin{tabularx}{\linewidth}{|X|X|X|X|}
\hline
\rowcolor{popblueL} Caractères & Pinyin & Nature & Traduction \\
\hline
钱 & qián & nom & l'argent, la monnaie \\
\hline
人民币 & rénmínbì & nom & le renminbi (monnaie chinoise) \\
\hline
非洲法郎 & Fēizhōu fǎláng & nom propre & le franc CFA (franc africain) \\
\hline
银行 & yínháng & nom & la banque \\
\hline
现金 & xiànjīn & nom & les espèces, le liquide \\
\hline
手机支付 & shǒujī zhīfù & nom (locution) & le paiement par téléphone portable \\
\hline
\end{tabularx}
\end{center}
\textbf{Ordre des traits (rappels) :}
\begin{itemize}
\item \zh{钱} : clé du métal \zh{钅} à gauche (5 traits, de haut en bas), puis \zh{戋} à droite ; on écrit toujours la partie gauche avant la partie droite.
\item \zh{银} : clé du métal \zh{钅} à gauche, puis \zh{艮} à droite.
\item \zh{行} : partie gauche \zh{彳} (deux traits obliques puis vertical), puis partie droite ; gauche avant droite.
\item \zh{现} : clé du jade \zh{王} à gauche (horizontal, horizontal, vertical, horizontal), puis \zh{见} à droite.
\item \zh{法} : clé de l'eau \zh{氵} (trois points, de haut en bas) à gauche, puis \zh{去} à droite.
\end{itemize}
L'élève recopie chaque mot trois fois en respectant ces ordres (haut $\rightarrow$ bas, gauche $\rightarrow$ droite).
\end{corrige}

\begin{corrige}[Exercice 4 — Associe caractères et pinyin]
\begin{center}
\begin{tabularx}{\linewidth}{|X|X|X|X|}
\hline
\rowcolor{popblueL} Caractères & Pinyin & Nature & Traduction \\
\hline
1. 人民币 & c. rénmínbì & nom & le renminbi \\
\hline
2. 银行 & a. yínháng & nom & la banque \\
\hline
3. 现金 & b. xiànjīn & nom & les espèces \\
\hline
4. 手机支付 & e. shǒujī zhīfù & nom (locution) & le paiement par téléphone \\
\hline
5. 非洲法郎 & d. Fēizhōu fǎláng & nom propre & le franc CFA \\
\hline
\end{tabularx}
\end{center}
\textbf{Points de vigilance :} ne pas confondre \zh{银行} (yínháng, banque) et \zh{现金} (xiànjīn, espèces) ; attention au ton de \zh{行}, prononcé \emph{háng} (2e ton) dans \zh{银行}, mais \emph{xíng} (2e ton aussi, sens « marcher / aller ») dans \zh{不行} (bù xíng, « ne pas aller / ne pas faire l'affaire »).
\end{corrige}

\begin{corrige}[Exercice 5 — Texte à trous]
Paragraphe complété :

\zh{在中国，人们用\textbf{人民币}，这是中国的钱。在喀麦隆，人们用\textbf{非洲法郎}。中国的\textbf{银行}发行人民币。在中国，很多人用\textbf{手机支付}，不太用\textbf{现金}。在喀麦隆，人们常常用现金买\textbf{东西}。}\\
\emph{Zài Zhōngguó, rénmen yòng rénmínbì, zhè shì Zhōngguó de qián. Zài Kāmàilóng, rénmen yòng Fēizhōu fǎláng. Zhōngguó de yínháng fāxíng rénmínbì. Zài Zhōngguó, hěn duō rén yòng shǒujī zhīfù, bú tài yòng xiànjīn. Zài Kāmàilóng, rénmen chángcháng yòng xiànjīn mǎi dōngxi.}

Traduction : « En Chine, les gens utilisent le renminbi, c'est la monnaie de la Chine. Au Cameroun, les gens utilisent le franc CFA. La banque de Chine émet le renminbi. En Chine, beaucoup de gens utilisent le paiement par téléphone et utilisent peu les espèces. Au Cameroun, les gens achètent souvent des choses en espèces. »

\textbf{Remarque :} le dernier trou attend un nom d'objet ; \zh{钱} de la banque de mots peut aussi convenir (\zh{用钱} « dépenser de l'argent »), mais \zh{东西} (dōngxi, choses) est la solution la plus naturelle ; accepter \zh{钱} si la phrase reste correcte.
\end{corrige}

\begin{corrige}[Exercice 6 — Traduction : thème]
\begin{enumerate}
\item \zh{在喀麦隆，人们用非洲法郎。} \emph{Zài Kāmàilóng, rénmen yòng Fēizhōu fǎláng.} Structure : complément de lieu \zh{在} + lieu en tête de phrase, puis sujet + verbe.
\item \zh{在中国，很多人用手机付钱。} \emph{Zài Zhōngguó, hěn duō rén yòng shǒujī fù qián.} On peut aussi dire \zh{用手机支付} ; \zh{付钱} (fù qián) = « payer ».
\item \zh{中国人民银行发行人民币。} \emph{Zhōngguó Rénmín Yínháng fāxíng rénmínbì.} Nom propre insécable ; \zh{发行} (fāxíng) = émettre.
\item \zh{我觉得手机支付很方便。} \emph{Wǒ juéde shǒujī zhīfù hěn fāngbiàn.} Structure d'opinion : \zh{我觉得} + proposition ; l'adjectif \zh{方便} est précédé de \zh{很} (pas de verbe « être » devant un adjectif).
\end{enumerate}
\textbf{Points de vigilance :} ne jamais traduire « est très pratique » par \zh{是很方便} ; en chinois, l'adjectif seul joue le rôle de prédicat : \zh{很方便}.
\end{corrige}

\begin{corrige}[Exercice 7 — Traduction : version]
Traduction : « La Banque populaire de Chine émet le renminbi. Au Cameroun, les gens utilisent souvent des espèces. »

\textbf{Découpage :}
\begin{itemize}
\item \zh{中国人民银行} : la Banque populaire de Chine (nom propre, banque centrale) ;
\item \zh{发行人民币} : émet le renminbi (\zh{发行} = émettre, mettre en circulation) ;
\item \zh{在喀麦隆} : au Cameroun (complément de lieu en tête) ;
\item \zh{常常用现金} : utilisent souvent des espèces (\zh{常常} = souvent ; \zh{用} = utiliser).
\end{itemize}
\end{corrige}

\begin{corrige}[Exercice 8 — Remise en ordre]
\begin{enumerate}
\item \zh{在中国，人们常常用手机支付。} \emph{Zài Zhōngguó, rénmen chángcháng yòng shǒujī zhīfù.} « En Chine, les gens paient souvent par téléphone. » Le complément de lieu \zh{在中国} se place en tête ; l'adverbe \zh{常常} se place avant le verbe \zh{用}.
\item \zh{非洲法郎是喀麦隆的钱。} \emph{Fēizhōu fǎláng shì Kāmàilóng de qián.} « Le franc CFA est la monnaie du Cameroun. » Structure : A \zh{是} B \zh{的} C ; le mot-outil \zh{的} marque la possession.
\item \zh{我觉得现金不方便。} \emph{Wǒ juéde xiànjīn bù fāngbiàn.} « Je trouve que les espèces ne sont pas pratiques. » La négation \zh{不} se place directement avant l'adjectif \zh{方便} ; ordre \zh{很方便} $\rightarrow$ \zh{不} + \zh{方便} (on ne dit pas \zh{很不方便} ici car la phrase donnée contient \zh{不} et \zh{很} : la forme \zh{不很方便} « pas très pratique » est aussi acceptable si l'on garde \zh{很}).
\end{enumerate}
\end{corrige}

\begin{corrige}[Exercice 9 — Compréhension de texte type Bac]
\textbf{Barème indicatif (10 points) :} 5 questions de compréhension $\times$ 1,5 pt = 7,5 pts ; question de langue = 2,5 pts (1 pt par mot relevé + traduction correcte, 0,5 pt pour la pertinence).

\textbf{Réponses attendues (en chinois, phrases complètes) :}
\begin{enumerate}
\item \zh{中国的钱叫人民币。} \emph{Zhōngguó de qián jiào rénmínbì.} « La monnaie de la Chine s'appelle le renminbi. »
\item \zh{中国人民银行发行人民币。} \emph{Zhōngguó Rénmín Yínháng fāxíng rénmínbì.} « C'est la Banque populaire de Chine qui émet le renminbi. »
\item \zh{现在很多人用手机支付付钱。} \emph{Xiànzài hěn duō rén yòng shǒujī zhīfù fù qián.} « Maintenant, beaucoup de gens paient par téléphone portable. »
\item \zh{在喀麦隆，人们常常用现金买东西。} \emph{Zài Kāmàilóng, rénmen chángcháng yòng xiànjīn mǎi dōngxi.} « Au Cameroun, les gens achètent souvent des choses en espèces. »
\item \zh{作者觉得两个国家的付钱方法不一样，但是都很有意思。} \emph{Zuòzhě juéde liǎng ge guójiā de fù qián fāngfǎ bù yíyàng, dànshì dōu hěn yǒu yìsi.} « L'auteur trouve que les méthodes de paiement des deux pays sont différentes, mais toutes les deux très intéressantes. »
\end{enumerate}

\textbf{Question de langue :} par exemple \zh{人民币} (rénmínbì) = « le renminbi » et \zh{现金} (xiànjīn) = « les espèces » ; autres réponses possibles : \zh{钱} (qián, argent), \zh{手机支付} (shǒujī zhīfù, paiement par téléphone), \zh{非洲法郎} (Fēizhōu fǎláng, franc CFA).

\textbf{Points de vigilance :} répondre par des phrases complètes avec sujet ; recopier les noms propres sans les déformer (\zh{中国人民银行}) ; \zh{不一样} = « pas pareil, différent ».
\end{corrige}

\begin{corrige}[Exercice 10 — Thème et version]
\textbf{Barème indicatif (10 points) :} thème 5 pts (2,5 + 2,5) ; version 5 pts. Pour chaque phrase : 1 pt pour l'exactitude des caractères, 1 pt pour la grammaire, 0,5 pt pour le respect du sens.

\textbf{A. Thème — corrigé :}
\begin{enumerate}
\item \zh{在喀麦隆，人们用非洲法郎。在中国，很多人用手机付钱。} \emph{Zài Kāmàilóng, rénmen yòng Fēizhōu fǎláng. Zài Zhōngguó, hěn duō rén yòng shǒujī fù qián.}
\item \zh{我觉得用现金不太方便，因为很慢。} \emph{Wǒ juéde yòng xiànjīn bú tài fāngbiàn, yīnwèi hěn màn.} Structure d'opinion : \zh{我觉得} + proposition, cause introduite par \zh{因为}.
\end{enumerate}

\textbf{B. Version — corrigé :} « Le renminbi est la monnaie de la Chine. En Chine, le paiement par téléphone portable est très pratique. Au Cameroun, les gens utilisent souvent des espèces, mais le paiement par téléphone est aussi de plus en plus répandu. »

\textbf{Points de vigilance :} \zh{越来越} + adjectif = « de plus en plus » ; \zh{但是} = « mais » ; ne pas oublier le complément de lieu \zh{在} + pays en tête de phrase ; pas de verbe \zh{是} devant l'adjectif \zh{方便}.
\end{corrige}

\begin{corrige}[Exercice 11 — Expression écrite et orale type Bac]
\textbf{Barème indicatif (10 points) :} respect de la consigne et du nombre de phrases (2 pts) ; correction des caractères (2 pts) ; grammaire et structures (3 pts) ; vocabulaire du thème (2 pts) ; opinion personnelle avec \zh{我觉得……因为……} (1 pt).

\textbf{A. Proposition de rédaction modèle :}

\zh{中国的钱叫人民币，喀麦隆的钱是非洲法郎。中国人民银行发行人民币。在中国，很多人用手机支付，不太用现金，买东西很方便。在喀麦隆，人们常常用现金买东西，但是现在手机支付也越来越多。我觉得手机支付很方便，因为很快。我也觉得现金很有用，因为每个人都能用。两个国家的付钱方法不一样，但是都很有意思。}\\
\emph{Zhōngguó de qián jiào rénmínbì, Kāmàilóng de qián shì Fēizhōu fǎláng. Zhōngguó Rénmín Yínháng fāxíng rénmínbì. Zài Zhōngguó, hěn duō rén yòng shǒujī zhīfù, bú tài yòng xiànjīn, mǎi dōngxi hěn fāngbiàn. Zài Kāmàilóng, rénmen chángcháng yòng xiànjīn mǎi dōngxi, dànshì xiànzài shǒujī zhīfù yě yuè lái yuè duō. Wǒ juéde shǒujī zhīfù hěn fāngbiàn, yīnwèi hěn kuài. Wǒ yě juéde xiànjīn hěn yǒuyòng, yīnwèi měi gè rén dōu néng yòng. Liǎng gè guójiā de fù qián fāngfǎ bù yíyàng, dànshì dōu hěn yǒu yìsi.}

Traduction : « La monnaie de la Chine s'appelle le renminbi, celle du Cameroun est le franc CFA. La Banque populaire de Chine émet le renminbi. En Chine, beaucoup de gens paient par téléphone et utilisent peu les espèces ; faire ses achats est très pratique. Au Cameroun, les gens achètent souvent en espèces, mais aujourd'hui le paiement par téléphone est de plus en plus fréquent. Je trouve que le paiement par téléphone est très pratique, parce que c'est rapide. Je trouve aussi que les espèces sont très utiles, parce que tout le monde peut les utiliser. Les méthodes de paiement des deux pays sont différentes, mais toutes les deux très intéressantes. » (8 phrases, environ 120 caractères.)

\textbf{B. Expression orale — grille de préparation :}
\begin{itemize}
\item Ouverture obligatoire : \zh{大家好！今天我谈谈我们国家的钱和中国的钱。} \emph{Dàjiā hǎo! Jīntiān wǒ tán tan wǒmen guójiā de qián hé Zhōngguó de qián.} « Bonjour à tous ! Aujourd'hui, je parle de la monnaie de notre pays et de celle de la Chine. »
\item Développement appuyé sur le tableau comparatif : nom des monnaies (\zh{人民币} / \zh{非洲法郎}), banque émettrice (\zh{中国人民银行}), moyen de paiement habituel (\zh{现金} / \zh{手机支付}), avantages (\zh{方便}、\zh{快}、\zh{有用}).
\item Conclusion avec opinion : \zh{我觉得……因为……。谢谢大家！} \emph{Wǒ juéde... yīnwèi... Xièxie dàjiā!}
\end{itemize}

\textbf{Points de vigilance :} bien prononcer les tons de \zh{人民币} (rén-mín-bì : 2-2-4) ; ne pas confondre \zh{钱} (argent) et \zh{前} (devant/avant), homophones \emph{qián} ; employer \zh{用} + instrument (\zh{用手机}、\zh{用现金}) ; conclure par une opinion personnelle claire.
\end{corrige}

\newpage

\coursec{Fiche bilan}

\begin{popfigure}
\begin{tikzpicture}[mindmap, grow cyclic, every node/.style={concept, font=\small}, concept color=popblue, root concept/.append style={concept color=popblue, font=\bfseries}, level 1/.append style={level distance=3.4cm, sibling angle=60}, level 2/.append style={level distance=2.2cm}]
\node[root concept, align=center]{人民币 与 中非法郎\\ \emph{Monnaies Chine / Cameroun}}
  child[concept color=poporange]{ node[align=center]{人民币 Rénmínbì\\ 元 yuán · 角 jiǎo · 分 fēn} }
  child[concept color=popgreen]{ node[align=center]{中非法郎\\ franc CFA (XAF)\\ 银行 yínháng} }
  child[concept color=poppurple]{ node[align=center]{换钱 huàn qián\\ 汇率 huìlǜ\\ 多少钱 duōshao qián} }
  child[concept color=poppink]{ node[align=center]{Grammaire\\ 比 bǐ + adj\\ 和……一样} }
  child[concept color=popgold]{ node[align=center]{Faits comparés\\ 现金 / 移动支付\\ 支付宝 · Orange Money} }
  child[concept color=popblueL]{ node[align=center]{Point de vue\\ 我觉得……\\ 因为……所以……} };
\end{tikzpicture}
\legende{Carte mentale de la leçon : deux monnaies, deux réalités}
\end{popfigure}

\begin{bilan}
\textbf{1. Lexique clé à connaître par cœur}
\begin{itemize}
  \item \zh{钱} qián : argent ; \zh{货币} huòbì : monnaie ; \zh{现金} xiànjīn : espèces, liquide
  \item \zh{人民币} Rénmínbì : monnaie chinoise ; \zh{元} yuán : yuan (unité) ; \zh{块} kuài : yuan (langue parlée)
  \item \zh{中非法郎} Zhōngfēi fǎláng : franc CFA ; \zh{硬币} yìngbì : pièce ; \zh{纸币} zhǐbì : billet
  \item \zh{银行} yínháng : banque ; \zh{换钱} huàn qián : changer de l'argent ; \zh{汇率} huìlǜ : taux de change
  \item \zh{移动支付} yídòng zhīfù : paiement mobile ; \zh{贵} guì : cher ; \zh{便宜} piányi : bon marché
\end{itemize}

\textbf{2. Structures de grammaire à maîtriser}
\begin{center}
\begin{tabularx}{\linewidth}{|X|X|X|}
\hline
\rowcolor{popblueL} Structure & Exemple (caractères + pinyin) & Traduction \\
\hline
A + 比 + B + adjectif & \zh{人民币比中非法郎稳定。} Rénmínbì bǐ Zhōngfēi fǎláng wěndìng. & Le yuan est plus stable que le franc CFA. \\
\hline
A + 和 + B + 一样 (+ adj) & \zh{中国的银行和喀麦隆的银行一样多。} Zhōngguó de yínháng hé Kāmàilóng de yínháng yíyàng duō. & Il y a autant de banques en Chine qu'au Cameroun. \\
\hline
我觉得 + phrase & \zh{我觉得移动支付很方便。} Wǒ juéde yídòng zhīfù hěn fāngbiàn. & Je trouve que le paiement mobile est pratique. \\
\hline
因为……所以…… & \zh{因为汇率会变，所以换钱要小心。} Yīnwèi huìlǜ huì biàn, suǒyǐ huàn qián yào xiǎoxīn. & Parce que le taux change, il faut être prudent en changeant. \\
\hline
多少钱？/ 几 + classifieur & \zh{这个多少钱？} Zhège duōshao qián? & Combien ça coûte ? \\
\hline
\end{tabularx}
\end{center}

\textbf{3. Faits comparés à retenir}
\begin{itemize}
  \item Chine : le \zh{人民币} (yuan), émis par la Banque populaire de Chine ; usage massif du paiement mobile (支付宝, 微信支付).
  \item Cameroun : le \zh{中非法郎} (franc CFA, zone CEMAC), monnaie commune à plusieurs pays d'Afrique centrale ; le liquide reste très utilisé, le mobile money (Orange Money, MTN MoMo) progresse.
  \item Le taux de change entre les deux monnaies varie : on change l'argent à la banque ou dans un bureau de change.
\end{itemize}

\textbf{4. Méthode express — mini-exposé comparatif}
\begin{enumerate}
  \item Introduire : \zh{今天我说一说中国和喀麦隆的钱。}
  \item Comparer avec 比 et 和……一样 (au moins 3 phrases).
  \item Donner son avis : \zh{我觉得……，因为……。}
  \item Conclure : \zh{谢谢大家！} Xièxie dàjiā !
\end{enumerate}
\end{bilan}

\begin{aretenir}[Checklist avant le Bac]
\begin{itemize}
  \item $\square$ Je connais le nom des deux monnaies : \zh{人民币} et \zh{中非法郎}.
  \item $\square$ Je sais écrire \zh{钱}、\zh{元}、\zh{银行} avec le bon ordre des traits.
  \item $\square$ Je sais demander un prix : \zh{这个多少钱？}
  \item $\square$ Je maîtrise la comparaison avec \zh{比} (A + 比 + B + adjectif).
  \item $\square$ Je sais dire « autant que » avec \zh{和……一样}.
  \item $\square$ Je distingue \zh{元} (écrit) et \zh{块} (parlé).
  \item $\square$ Je peux citer deux moyens de paiement : \zh{现金} et \zh{移动支付}.
  \item $\square$ Je sais exprimer mon avis avec \zh{我觉得……因为……}.
  \item $\square$ Je peux expliquer ce qu'est un taux de change (\zh{汇率}).
  \item $\square$ Je suis capable de présenter un mini-exposé comparatif de 5 phrases en chinois.
\end{itemize}
\end{aretenir}

\findecours

\end{document}
